\documentclass[11pt]{article}
\usepackage[utf8]{inputenc}
\usepackage[T1]{fontenc}
\usepackage{amsmath,amssymb,amsthm}
\usepackage[margin=1.1in]{geometry}
\allowdisplaybreaks
\emergencystretch=1.5em

\newtheorem{theorem}{Theorem}[section]
\newtheorem{proposition}[theorem]{Proposition}
\newtheorem{lemma}[theorem]{Lemma}
\newtheorem{corollary}[theorem]{Corollary}
\theoremstyle{definition}
\newtheorem{definition}[theorem]{Definition}
\theoremstyle{remark}
\newtheorem{remark}[theorem]{Remark}
\theoremstyle{plain}
\newtheorem{auxlemma}{Auxiliary Lemma}

\newcommand{\R}{\mathbb{R}}
\newcommand{\C}{\mathbb{C}}
\newcommand{\Z}{\mathbb{Z}}
\newcommand{\cO}{\mathcal{O}}
\newcommand{\cP}{\mathcal{P}}
\newcommand{\cQ}{\mathcal{Q}}
\newcommand{\cU}{\mathcal{U}}
\newcommand{\cM}{\mathcal{M}}
\newcommand{\cG}{\mathcal{G}}
\newcommand{\cN}{\mathcal{N}}
\newcommand{\cH}{\mathcal{H}}
\newcommand{\cT}{\mathcal{T}}
\newcommand{\cF}{\mathcal{F}}
\newcommand{\cR}{\mathcal{R}}
\newcommand{\cV}{\mathcal{V}}
\newcommand{\cW}{\mathcal{W}}
\newcommand{\Sym}{\operatorname{Sym}}
\newcommand{\rank}{\operatorname{rank}}
\newcommand{\adj}{\operatorname{adj}}
\newcommand{\acc}{\operatorname{acc}}
\newcommand{\logit}{\operatorname{logit}}
\newcommand{\sig}{\sigma}
\newcommand{\TF}{\mathsf{TF}}
\newcommand{\SM}{\operatorname{SoftMaxCol}^{\mathrm{c}}}
\newcommand{\spn}{\operatorname{span}}
\newcommand{\dist}{\operatorname{dist}}
\newcommand{\clo}{\operatorname{cl}}

\title{Identifiability of deep causal attention-only transformers\\ with skip connections}
\author{}
\date{}

\begin{document}
\maketitle

\begin{abstract}
We prove that an $L$-layer attention-only transformer with causal masking, a single
attention head per layer, and additive skip connections is identifiable from its
input--output map: for all parameters $\theta'$ outside an explicit finite union of proper
algebraic subsets of parameter space, any parameter $\theta$ computing the same function as
$\theta'$ equals $\theta'$ exactly, layer by layer.  At the level of the weight
factorizations $V_\ell=W_O^{(\ell)}W_V^{(\ell)}$ and
$A_\ell=(W_K^{(\ell)})^{\top}W_Q^{(\ell)}$ the fiber of the realization map is the product
over layers of the two general-linear orbits, with no cross-layer symmetries.  The proof is
by induction on depth.  Complex-analytic continuation in an inverse-temperature variable
identifies the innermost attention matrix through an $L$-tier hierarchy of accumulating
singularities, in the spirit of Fefferman's reconstruction of $\tanh$ networks
\cite{fefferman1994reconstructing}; everything thereafter is real-variable: exact saturated
limits of the attention gates along dialed input paths, a structural trichotomy that
eliminates all exotic saturation profiles by linear algebra, and a cancellation identity
(``$K=I$'') created by the skip connections, which removes any dependence on the unknown
parameter's saturation pattern.  All genericity conditions are explicit polynomial
nonvanishing conditions, with verified witnesses.
\end{abstract}

\section{Introduction}\label{sec:intro}

Fix integers $L\ge 1$, $d\ge 2$ and $r\ge 2$.  We study the map computed by an $L$-layer,
single-head, causal, attention-only transformer with skip connections on sequences of
$T:=r+1$ tokens in $\R^d$.  The parameter of the network is a tuple
\[
\theta=(V_1,A_1,\dots,V_L,A_L)\in\cP:=(\R^{d\times d})^{2L},
\]
where $V_\ell$ is the (combined output--value) mixing matrix and $A_\ell$ the (combined
key--query) attention matrix of layer $\ell$.  The question we answer is: \emph{does the
input--output map determine $\theta$?}  Our main result
(Theorem~\ref{thm:main}) is that it does, exactly, for all $\theta$ outside a finite union
of proper algebraic subsets of $\cP$, provided $d$ is at least an explicit polynomial in
$L$.  When the matrices $V_\ell,A_\ell$ are further factored into the usual four weight
matrices, the factors are determined up to the obvious per-layer general-linear changes of
basis (Corollary~\ref{cor:factors}); because the layers are ordered and there is one head
per layer, no permutation symmetries occur.

Two features of the architecture are essential to the proof and are worth highlighting.
\emph{Causal masking} makes a two-token-type probe family close under the layer map: on an
input whose first $r$ tokens are equal, every layer acts on the probe by an explicit
two-dimensional recursion driven by scalar sigmoid gates
(Proposition~\ref{prop:probe}).  \emph{Skip connections} force the algebraic identity
$B_\ell-V_\ell=I$ for the matrices $B_\ell:=I+V_\ell$ transporting the token streams; this
single identity is responsible both for the transparency of fully attending layers (used to
control the generic parameter's saturation cascade) and for the cancellation that
identifies the first layer of the \emph{unknown} parameter without any control over its
saturation pattern (Proposition~\ref{prop:matching}).

The proof is an induction on depth over a strengthened statement
(Theorem~\ref{thm:IDL}): agreement of the two networks along all
\emph{$\tau$-dependent probe paths} whose limit points sweep an arbitrary open set meeting
the anchor set forces equality of parameters.  The strengthening is what makes the
induction close: after the first layer is identified, the deeper layers are compared along
the image paths of a first-layer ``sweep'', and that image is an open set one does not get
to choose.  The inductive step has three parts.

\smallskip
\emph{Step 1 (\S\ref{sec:A1}).}  Along constant probes, the observable extends to a
meromorphic-type function of the complex inverse temperature $\tau$, with a hierarchy of
singular sets: the first gate has poles on a vertical arithmetic progression whose real
part is $-b/\lambda_1$; the poles of each deeper gate accumulate at poles of the previous
one.  Running this hierarchy $L$ tiers deep on the generic side and intersecting with the
($\le L$)-fold stratification of the unknown side pins down the innermost slope
$\lambda_1$, hence $A_1$.  This is the only complex analysis in the paper; it follows
Fefferman's singular-set strategy for $\tanh$ networks \cite{fefferman1994reconstructing},
with arithmetic progressions of poles of the logistic sigmoid playing the role of the poles
of $\tanh$.

\smallskip
\emph{Step 2 (\S\ref{sec:cascade}).}  Real saturated limits along \emph{dialed} paths ---
paths based at the quadric $\{w^\top A_1 v=0\}$ whose first gate converges to a freely
chosen value $t\in(0,1)$ --- reduce both networks to affine limit maps.  The generic
parameter's gates all saturate to $1$ on a suitable sign region; the unknown parameter's
gates saturate to \emph{some} constants in $\{0,1,\alpha\}$ after a trichotomy
(Proposition~\ref{prop:trichotomy}) kills every other possibility by linear algebra against
$\det A_1\neq0$, $\Sym A_1\neq 0$ and a depth certificate $V_1\neq 0$ produced in Step~1.
Matching the two affine limit maps and using $B-V=I$ on both sides cancels the unknown
saturation data and yields $V_1=V_1'$, $B_1=B_1'$.

\smallskip
\emph{Step 3 (\S\ref{sec:descent}).}  With the first layer common, an exact peeling
identity transports the hypothesis to the depth-$(L-1)$ networks along effective paths.  A
realization lemma (a scalar Banach fixed point, Lemma~\ref{lem:realization}) shows that
\emph{every} path near the swept image arises as an effective path, and an inverse-function
argument shows the sweep covers an open set; the induction hypothesis then finishes the
proof.

\smallskip
Section~\ref{sec:generic} constructs the genericity set: each condition is the
nonvanishing of finitely many explicit polynomials of $\theta'$, and each is verified at an
explicit witness.  The anchor points needed by Steps 2--3 at every level of the recursion
are produced by a single linear solve against a rank certificate; this costs the dimension
hypothesis $d\ge\binom{L}{2}+2(L-1)$, which we have not tried to optimize
(Remark~\ref{rem:dimension}).

\subsection*{Relation to prior work}
Fefferman \cite{fefferman1994reconstructing} proved identifiability of $\tanh$
multilayer perceptrons from their output map using analytic continuation and the
arithmetic of pole progressions; one hidden layer was treated earlier by Sussmann
\cite{sussmann1992} and Albertini--Sontag \cite{albertinisontag}.  For attention, the
one-layer multi-head non-causal case (without skip connections) can be handled by a
linear-independence argument for sigmoids of different slopes; the present paper treats
arbitrary depth for the causal single-head skip architecture.  The rank-factorization
lemma used for the factor-level statement appears in \cite{henry2025geometry}.

\section{The model, the probe family, and the recursion}\label{sec:model}

\subsection{Notation}\label{subsec:notation}
$\sig(x):=(1+e^{-x})^{-1}$ denotes the logistic sigmoid and
$\logit:(0,1)\to\R$, $\logit t=\log\frac{t}{1-t}$, its inverse.  We fix once and for all an
integer $r\ge2$ and set
\[
T:=r+1,\qquad b:=\log r>0,\qquad \alpha:=\sig(b)=\tfrac{r}{r+1}\in(0,1).
\]
For a square matrix $M$, $\Sym(M):=\tfrac12(M+M^\top)$; note that
$x^\top Mx=x^\top\Sym(M)x$ for every real vector $x$, and that a symmetric matrix whose
quadratic form vanishes identically is zero, by polarization:
$2x^\top My=(x+y)^\top M(x+y)-x^\top Mx-y^\top My$.  For matrices $M_i$ indexed by layers
we write $M_{j:i}:=M_jM_{j-1}\cdots M_i$, with the convention $M_{j:i}=I$ when $j<i$; empty
sums are $0$ and empty products are $I$ (for scalars, $1$).  Throughout,
$\Pi:=(2\Z+1)\pi i\subset\C$ denotes the set of odd multiples of $\pi i$; $\Pi$ is closed,
$\Pi\cap\R=\varnothing$ (its points are nonzero and purely imaginary), and distinct points
of $\Pi$ are at distance at least $2\pi$.

We record the elementary sigmoid calculus used throughout.  For real $x$,
$\sig'(x)=\sig(x)\bigl(1-\sig(x)\bigr)\in(0,\tfrac14]$, so $\sig$ is strictly increasing on
$\R$ and $|\sig'|\le\tfrac14$ there; by definition $\logit\circ\sig=\mathrm{id}_\R$ and
$\sig\circ\logit=\mathrm{id}_{(0,1)}$; and both $\sig|_\R$ and $\logit$ are real-analytic
with real Taylor coefficients at every point of $\R$, resp.\ of $(0,1)$ ($\sig$ because it
is holomorphic on a neighbourhood of $\R$ --- see Lemma~\ref{lem:sigmoid} --- and real on
$\R$; $\logit$ because $\logit t=\log t-\log(1-t)$).

\subsection{The network}
Causal column softmax: for $M\in\R^{T\times T}$,
\[
\big[\SM(M)\big]_{ij}:=\frac{e^{M_{ij}}\,\mathbf 1[i\le j]}{\sum_{i'\le j}e^{M_{i'j}}}.
\]
Thus, with $M=X^\top AX$, the convention is $M_{ij}=X_{:,i}^\top A\,X_{:,j}$, and the
$j$-th output column of the attention attends only to input columns $i\le j$.  A single
causal attention layer with parameters $(V,A)\in(\R^{d\times d})^2$ and skip connection
acts on $X\in\R^{d\times T}$ by
\begin{equation}\label{eq:layer}
\mathrm{Layer}_{V,A}(X):=X+V\,X\,\SM\!\big(X^\top A X\big).
\end{equation}
The network with parameter $\theta=(V_1,A_1,\dots,V_L,A_L)$ is
\[
\TF_\theta:=\mathrm{Layer}_{V_L,A_L}\circ\cdots\circ\mathrm{Layer}_{V_1,A_1}
:\R^{d\times T}\to\R^{d\times T}.
\]
We set $B_\ell:=I+V_\ell$.  Two trivial identities are used constantly:
\begin{equation}\label{eq:transparency}
B_\ell-1\cdot V_\ell=I,\qquad B_\ell-0\cdot V_\ell=B_\ell .
\end{equation}

\subsection{Probes and the two-type recursion}
For $(w,v)\in\R^d\times\R^d$ and $\tau>0$ put $u:=w+v$ and define the probe input
\[
X_{w,v}(\tau):=\sqrt{\tau}\,\big[\,\underbrace{u\ \cdots\ u}_{r}\ \ v\,\big]\in\R^{d\times T}.
\]

\begin{proposition}[Probe reduction]\label{prop:probe}
Fix $\theta\in\cP$, $(w,v)\in\R^{2d}$ and $\tau>0$.  Define $w_0:=w$, $v_0:=v$,
$u_0:=u=w+v$ and, for $1\le\ell\le L$,
\begin{equation}\label{eq:recursion}
\lambda_\ell:=w_{\ell-1}^\top A_\ell v_{\ell-1},\quad
s_\ell:=\sig(\tau\lambda_\ell+b),\quad
u_\ell:=B_\ell u_{\ell-1},\quad
v_\ell:=B_\ell v_{\ell-1}+s_\ell V_\ell w_{\ell-1},\quad
w_\ell:=u_\ell-v_\ell .
\end{equation}
Then $w_\ell=(B_\ell-s_\ell V_\ell)w_{\ell-1}$, and for every $0\le\ell\le L$ the hidden
state after $\ell$ layers is
$X^{(\ell)}=\sqrt\tau\,[\,u_\ell\ \cdots\ u_\ell\ \ v_\ell\,]$.  In particular the last
column of the output determines
\begin{equation}\label{eq:F}
F^{(L)}_\theta(w,v,\tau):=\tfrac{1}{\sqrt\tau}\big[\TF_\theta(X_{w,v}(\tau))\big]_{:,T}=v_L .
\end{equation}
\textup{(}If $V_1=0$, then $B_1=I$, so $w_1=w$, $v_1=v$: the first layer leaves the probe
data unchanged, the gate $s_1$ drops out of the recursion, and consequently
$F^{(m)}_\vartheta(w,v,\cdot)=F^{(m-1)}_{\vartheta_{\ge2}}(w,v,\cdot)$ for every $m$-tuple
$\vartheta$ with $V_1=0$, in the notation introduced after the proof.\textup{)}
\end{proposition}

\begin{proof}
Write $X^{(0)}:=X_{w,v}(\tau)$ and let $X^{(\ell)}$ denote the hidden state after $\ell$
layers, so $X^{(\ell)}=\mathrm{Layer}_{V_\ell,A_\ell}(X^{(\ell-1)})$.  We prove by
induction on $\ell\in\{0,1,\dots,L\}$ the statement
\[
(P_\ell):\qquad
X^{(\ell)}=\sqrt\tau\,[\,\underbrace{u_\ell\ \cdots\ u_\ell}_{r}\ \ v_\ell\,]
\quad\text{and}\quad u_\ell=w_\ell+v_\ell .
\]

\emph{Base case $\ell=0$.}  By definition of the probe input,
$X^{(0)}=\sqrt\tau\,[\,u_0\cdots u_0\ v_0\,]$ with $u_0=w+v=w_0+v_0$.

\emph{Inductive step.}  Assume $(P_{\ell-1})$ and write $X:=X^{(\ell-1)}=\sqrt\tau\,C$,
where $C=[\,c_1\ \cdots\ c_T\,]$ has $c_1=\dots=c_r=u_{\ell-1}$ and $c_T=v_{\ell-1}$
($T=r+1$).

\emph{Scores.}  $S:=X^\top A_\ell X=(\sqrt\tau)^2\,C^\top A_\ell C=\tau\,C^\top A_\ell C$,
i.e.\ $S_{ij}=\tau\,c_i^\top A_\ell c_j$ in the convention
$S_{ij}=X_{:,i}^\top A_\ell X_{:,j}$.

\emph{The layer, column by column.}  Put $\Gamma:=\SM(S)$.  By the causal convention
$\Gamma_{ij}=0$ for $i>j$, while for $i\le j$
\[
\Gamma_{ij}=\frac{e^{S_{ij}}}{\sum_{i'\le j}e^{S_{i'j}}}>0,
\qquad\sum_{i\le j}\Gamma_{ij}=1 .
\]
Since $(X\Gamma)_{:,j}=\sum_{i=1}^{T}X_{:,i}\Gamma_{ij}=\sqrt\tau\sum_{i\le j}\Gamma_{ij}c_i$,
\[
\big[\mathrm{Layer}_{V_\ell,A_\ell}(X)\big]_{:,j}
=X_{:,j}+V_\ell\,(X\Gamma)_{:,j}
=\sqrt\tau\,\big(c_j+V_\ell\,\bar c_j\big),
\qquad
\bar c_j:=\sum_{i\le j}\Gamma_{ij}\,c_i,
\]
and $\bar c_j$ is a convex combination of $c_1,\dots,c_j$ (nonnegative weights summing
to~$1$).

\emph{Columns $j\le r$.}  Here $c_1=\dots=c_j=u_{\ell-1}$, so for \emph{any} softmax
weights
\[
\bar c_j=\Big(\sum_{i\le j}\Gamma_{ij}\Big)\,u_{\ell-1}=u_{\ell-1},
\]
independently of the scores.  Hence column $j$ becomes
$\sqrt\tau\,(u_{\ell-1}+V_\ell u_{\ell-1})=\sqrt\tau\,B_\ell u_{\ell-1}
=\sqrt\tau\,u_\ell$, the same vector for every $j\le r$: the $r$-fold repetition is
preserved.

\emph{Column $j=T=r+1$.}  This column attends to the whole context ($i\le T$).  With
\[
\beta:=u_{\ell-1}^\top A_\ell v_{\ell-1},
\qquad
\delta:=v_{\ell-1}^\top A_\ell v_{\ell-1},
\]
the scores in column $T$ are $S_{i,T}=\tau\beta$ for $i\le r$ and $S_{T,T}=\tau\delta$.
The softmax denominator is $D:=r\,e^{\tau\beta}+e^{\tau\delta}$, and
$\Gamma_{iT}=e^{\tau\beta}/D$ for $i\le r$, $\Gamma_{TT}=e^{\tau\delta}/D$.
The total weight on the $u$-block is therefore
\begin{align*}
p:=\sum_{i\le r}\Gamma_{iT}
=\frac{r\,e^{\tau\beta}}{r\,e^{\tau\beta}+e^{\tau\delta}}
&=\frac{1}{1+\dfrac{e^{\tau\delta}}{r\,e^{\tau\beta}}}
=\frac{1}{1+e^{\tau(\delta-\beta)-\log r}}
=\sig\big(\tau(\beta-\delta)+\log r\big)
=\sig\big(\tau(\beta-\delta)+b\big),
\end{align*}
using $\sig(x)=(1+e^{-x})^{-1}$.  This is the only point at which the constant $b$ enters;
it equals $\log r$ \emph{exactly}, forced by the multiplicity $r$ of the identical key
columns inside the softmax normalization over the $T=r+1$ admissible rows of column $T$.
Moreover, by $(P_{\ell-1})$ and $w_{\ell-1}=u_{\ell-1}-v_{\ell-1}$,
\[
\beta-\delta=(u_{\ell-1}-v_{\ell-1})^\top A_\ell v_{\ell-1}
=w_{\ell-1}^\top A_\ell v_{\ell-1}=\lambda_\ell,
\]
so $p=\sig(\tau\lambda_\ell+b)=s_\ell$, and
\[
\bar c_T=p\,u_{\ell-1}+(1-p)\,v_{\ell-1}
=v_{\ell-1}+p\,(u_{\ell-1}-v_{\ell-1})
=v_{\ell-1}+s_\ell\,w_{\ell-1}.
\]
Column $T$ therefore becomes
\[
\sqrt\tau\,\big(v_{\ell-1}+V_\ell(v_{\ell-1}+s_\ell w_{\ell-1})\big)
=\sqrt\tau\,\big((I+V_\ell)v_{\ell-1}+s_\ell V_\ell w_{\ell-1}\big)
=\sqrt\tau\,\big(B_\ell v_{\ell-1}+s_\ell V_\ell w_{\ell-1}\big)
=\sqrt\tau\,v_\ell .
\]

\emph{Closing the induction.}  The two displays give
$X^{(\ell)}=\sqrt\tau\,[\,u_\ell\cdots u_\ell\ v_\ell\,]$.  Next, by definition
$w_\ell:=u_\ell-v_\ell$, so using $u_{\ell-1}=w_{\ell-1}+v_{\ell-1}$,
\[
w_\ell
=B_\ell u_{\ell-1}-B_\ell v_{\ell-1}-s_\ell V_\ell w_{\ell-1}
=B_\ell(u_{\ell-1}-v_{\ell-1})-s_\ell V_\ell w_{\ell-1}
=(B_\ell-s_\ell V_\ell)\,w_{\ell-1},
\]
which is the claimed closed form, and $u_\ell=w_\ell+v_\ell$ holds by the definition of
$w_\ell$.  This is $(P_\ell)$.

\emph{Output.}  By $(P_L)$,
$\big[\TF_\theta(X_{w,v}(\tau))\big]_{:,T}=X^{(L)}_{:,T}=\sqrt\tau\,v_L$, so the
$\sqrt\tau$ scaling cancels and $F^{(L)}_\theta(w,v,\tau)=v_L$.  The argument is verbatim
the same for any $m$-tuple $\vartheta$, giving $F^{(m)}_\vartheta(w,v,\tau)=v_m$; in
particular $F^{(0)}_{()}(w,v,\tau)=v_0=v$ for the empty tuple.  Finally, the parenthetical
claim is immediate: $V_1=0$ gives $B_1=I$, $w_1=(I-s_1\cdot 0)w=w$ and
$v_1=Iv+s_1\cdot 0\cdot w=v$, after which the recursion for layers $2,\dots,m$ coincides
with the $(m-1)$-layer recursion of $\vartheta_{\ge2}$ started at $(w,v)$ at the same
$\tau$.
\end{proof}

For any $m$-tuple $\vartheta=(V_1,A_1,\dots,V_m,A_m)$ we write
$F^{(m)}_\vartheta(w,v,\tau):=v_m$ for the quantity defined by the recursion
\eqref{eq:recursion} run for $m$ layers.  The recursion immediately gives the exact peeling
identity, valid pointwise:

\begin{lemma}[Peeling]\label{lem:peeling}
For every $\theta$, $(w,v)\in\R^{2d}$, $\tau>0$, with
$s_1=\sig(\tau\,w^\top A_1v+b)$,
\[
F^{(L)}_\theta(w,v,\tau)
=F^{(L-1)}_{\theta_{\ge2}}\big((B_1-s_1V_1)w,\;B_1v+s_1V_1w,\;\tau\big),
\]
where $\theta_{\ge2}:=(V_2,A_2,\dots,V_L,A_L)$.
\end{lemma}

\begin{proof}
By Proposition~\ref{prop:probe} (the case $\ell=1$), the hidden state after the first
layer is $X^{(1)}=\sqrt\tau\,[\,u_1\cdots u_1\ v_1\,]$ ($r$ copies of $u_1$, then $v_1$),
with
\[
\lambda_1=w_0^\top A_1v_0=w^\top A_1v,\qquad
s_1=\sig(\tau\lambda_1+b),\qquad
w_1=(B_1-s_1V_1)w,\qquad
v_1=B_1v+s_1V_1w,
\]
and $u_1=B_1(w+v)$.  The probe form is preserved exactly:
\[
w_1+v_1=(B_1-s_1V_1)w+B_1v+s_1V_1w=B_1w+B_1v=B_1(w+v)=u_1,
\]
so $X^{(1)}=\sqrt\tau\,[\,(w_1+v_1)\cdots(w_1+v_1)\ \ v_1\,]=X_{w_1,v_1}(\tau)$ --- by
definition of the probe input, this is precisely the probe at the pair $(w_1,v_1)$ and the
\emph{same} inverse temperature $\tau$, with the same $r$-then-$1$ column layout and the
same $\sqrt\tau$ scaling.

Since $\TF_\theta=\TF_{\theta_{\ge2}}\circ\mathrm{Layer}_{V_1,A_1}$ (composition of the
layer maps in order), we have
$\TF_\theta(X_{w,v}(\tau))=\TF_{\theta_{\ge2}}\big(X^{(1)}\big)
=\TF_{\theta_{\ge2}}\big(X_{w_1,v_1}(\tau)\big)$.
Applying Proposition~\ref{prop:probe} to the $(L-1)$-tuple $\theta_{\ge2}$ and the probe
$X_{w_1,v_1}(\tau)$, and reading off the last column,
\[
F^{(L)}_\theta(w,v,\tau)
=\tfrac1{\sqrt\tau}\big[\TF_{\theta_{\ge2}}(X_{w_1,v_1}(\tau))\big]_{:,T}
=F^{(L-1)}_{\theta_{\ge2}}(w_1,v_1,\tau),
\]
which is the claim.  (For $L=1$ the right-hand side is $F^{(0)}_{()}(w_1,v_1,\tau)=v_1$,
consistent with the convention $F^{(m)}_\vartheta=v_m$.)
\end{proof}

\subsection{Slope polynomials}\label{subsec:slope}
Let $\vartheta$ be an $m$-layer parameter and let $z=(z_1,z_2,\dots)$ be indeterminates.
Define formal streams $w_0(z):=w$, $v_0(z):=v$ and
\[
w_j(z):=(B_j-z_jV_j)\,w_{j-1}(z),\qquad
v_j(z):=B_jv_{j-1}(z)+z_jV_jw_{j-1}(z),
\]
and the \emph{slope polynomials}
\[
\phi^\vartheta_\ell(z_1,\dots,z_{\ell-1};w,v):=w_{\ell-1}(z)^\top A_\ell\,v_{\ell-1}(z),
\qquad \phi^\vartheta_1(w,v):=w^\top A_1v=:q_\vartheta(w,v).
\]
We suppress $\vartheta$ when clear.  An immediate induction shows that $w_j(z)$ and
$v_j(z)$ are vector-valued polynomials in $z_1,\dots,z_j$ alone (each recursion step is
affine in the newly introduced variable, with matrix coefficients), so $\phi_\ell$ is a
genuine polynomial in $z_1,\dots,z_{\ell-1}$ whose coefficients depend on
$(\vartheta,w,v)$.

\begin{lemma}[Structure of the slope polynomials]\label{lem:phi}
Write $W_{j}(z):=(B_{j}-z_{j}V_{j})\cdots(B_1-z_1V_1)$ (so $W_0=I$) and
$T_{j}(z):=\sum_{i=1}^{j}z_i\,B_{j:i+1}V_iW_{i-1}(z)$.  Then:
\begin{enumerate}
\item[(i)] $w_j(z)=W_j(z)w$, $v_j(z)=B_{j:1}v+T_j(z)w$, and
\[
\phi_\ell(z;w,v)=w^\top X_\ell(z)\,v+w^\top Y_\ell(z)\,w,\qquad
X_\ell(z):=W_{\ell-1}(z)^\top A_\ell B_{\ell-1:1},\quad
Y_\ell(z):=W_{\ell-1}(z)^\top A_\ell T_{\ell-1}(z);
\]
in particular $\phi_\ell$ is an affine function of $v$ with $v$-gradient
$\nabla_v\phi_\ell=(B_{\ell-1:1})^\top A_\ell^\top W_{\ell-1}(z)w$.
\item[(ii)] For $\ell\ge2$: $\deg_{z_i}\phi_\ell\le 2$ for each $i\le\ell-1$, and the
coefficient of $z_1^2z_2^2\cdots z_{\ell-1}^2$ equals
$-\,w^\top \Sym\!\big(V_{\ell-1:1}^\top A_\ell V_{\ell-1:1}\big)w$.
\item[(iii)] For $\ell\ge2$: the coefficient of $z_{\ell-1}^2$ in $\phi_\ell$ is the
polynomial
$f_\ell(z_1,\dots,z_{\ell-2}):=-\,w_{\ell-2}(z)^\top V_{\ell-1}^\top A_\ell
V_{\ell-1}w_{\ell-2}(z)$, whose coefficient of $z_1^2\cdots z_{\ell-2}^2$ equals
$-\,w^\top\Sym\!\big(V_{\ell-1:1}^\top A_\ell V_{\ell-1:1}\big)w$ as well.
\item[(iv)] (\emph{Transparency.})  For $\ell\ge2$: with $z_2=\dots=z_{\ell-1}=1$ one has
$W_{\ell-1}(z_1,1,\dots,1)=B_1-z_1V_1$ by \eqref{eq:transparency}.
\item[(v)] (\emph{Dressing.})  For $m\ge1$: the vector polynomial
$g_m(z):=V_mW_{m-1}(z)w$ has $\deg_{z_i}\le1$ and the coefficient of
$z_1\cdots z_{m-1}$ equals $(-1)^{m-1}V_{m:1}w$.
\end{enumerate}
\end{lemma}

\begin{proof}
\emph{Step 0: multilinear expansion of $W_j$, and a unique-split rule.}
For $E\subseteq[j]:=\{1,\dots,j\}$ let $M^{(j)}_E:=C_jC_{j-1}\cdots C_1$ denote the ordered
product with $C_i:=V_i$ for $i\in E$ and $C_i:=B_i$ for $i\notin E$.  Expanding each factor
of $W_j(z)=(B_j-z_jV_j)\cdots(B_1-z_1V_1)$ and collecting terms,
\begin{equation}\label{eq:Wexp}
W_j(z)=\sum_{E\subseteq[j]}(-1)^{|E|}\Big(\prod_{i\in E}z_i\Big)\,M^{(j)}_E .
\end{equation}
Distinct subsets $E$ produce distinct monomials $\prod_{i\in E}z_i$, so \eqref{eq:Wexp} is
the monomial expansion of $W_j$.  Consequences: \textbf{(a)} $W_j$ has degree $\le1$ in
each of $z_1,\dots,z_j$ (\emph{multilinear}); \textbf{(b)} its constant term
($E=\varnothing$) is $B_{j:1}$; \textbf{(c)} its coefficient of the top monomial
$z_1\cdots z_j$ ($E=[j]$) is $(-1)^{j}V_{j:1}$.

\emph{Unique-split rule.}  Let $a(z)=\sum_m a_m\,m$ and $t(z)=\sum_m t_m\,m$ be
$\R^d$-valued polynomials in $z_1,\dots,z_k$, each of degree $\le1$ in every variable, the
sums running over square-free monomials $m$, with vector coefficients $a_m,t_m\in\R^d$.
For any $A\in\R^{d\times d}$, the coefficient of $z_1^2\cdots z_k^2$ in $a(z)^\top A\,t(z)$
equals $a_{z_1\cdots z_k}^\top A\,t_{z_1\cdots z_k}$: indeed the coefficient is
$\sum_{m_1m_2=z_1^2\cdots z_k^2}a_{m_1}^\top A\,t_{m_2}$ over pairs of square-free
monomials, and a square-free factorization of $z_1^2\cdots z_k^2$ forces each $m_1,m_2$ to
contain every $z_i$ exactly once, i.e.\ $m_1=m_2=z_1\cdots z_k$.

\emph{(i)}  The product form $w_j(z)=W_j(z)w$ is immediate by induction from
$w_j(z)=(B_j-z_jV_j)w_{j-1}(z)$ and $w_0=w$.  For $v_j$, induct on $j$: the base is
$v_0=v=B_{0:1}v+T_0w$ ($B_{0:1}=I$ as an empty product, $T_0=0$ as an empty sum), and
\[
v_j=B_jv_{j-1}+z_jV_jw_{j-1}
=B_j\big(B_{j-1:1}v+T_{j-1}w\big)+z_jV_jW_{j-1}w
=B_{j:1}v+\big(B_jT_{j-1}+z_jV_jW_{j-1}\big)w,
\]
where
\[
B_jT_{j-1}+z_jV_jW_{j-1}
=\sum_{i=1}^{j-1}z_i\,B_jB_{j-1:i+1}V_iW_{i-1}+z_j\,B_{j:j+1}V_jW_{j-1}
=\sum_{i=1}^{j}z_i\,B_{j:i+1}V_iW_{i-1}
=T_j,
\]
using the telescoping $B_jB_{j-1:i+1}=B_{j:i+1}$ and the empty product
$B_{j:j+1}=I$.  Hence
\[
\phi_\ell=w_{\ell-1}(z)^\top A_\ell\,v_{\ell-1}(z)
=(W_{\ell-1}w)^\top A_\ell\big(B_{\ell-1:1}v+T_{\ell-1}w\big)
=w^\top X_\ell\,v+w^\top Y_\ell\,w
\]
with $X_\ell,Y_\ell$ as stated.  As $\phi_\ell$ is affine in $v$ with linear part
$w^\top X_\ell v$,
$\nabla_v\phi_\ell=X_\ell^\top w
=(W_{\ell-1}^\top A_\ell B_{\ell-1:1})^\top w
=(B_{\ell-1:1})^\top A_\ell^\top W_{\ell-1}(z)\,w$.

\emph{(ii)}  Fix $\ell\ge2$ and $1\le i\le\ell-1$.

\emph{Degrees.}  By Step~0(a), $W_{\ell-1}(z)$ is multilinear, hence so is
$w^\top X_\ell v=(W_{\ell-1}w)^\top A_\ell B_{\ell-1:1}v$ ($B_{\ell-1:1}$ is constant in
$z$): degree $\le1$ in each $z_i$.  In
$T_{\ell-1}=\sum_{i'=1}^{\ell-1}z_{i'}B_{\ell-1:i'+1}V_{i'}W_{i'-1}(z)$, the $i'$-th
summand involves only $z_1,\dots,z_{i'}$, with degree exactly $1$ in $z_{i'}$ (the
explicit factor; $W_{i'-1}$ is free of $z_{i'}$) and degree $\le1$ in each $z_{i''}$,
$i''<i'$ (multilinearity of $W_{i'-1}$).  Hence $T_{\ell-1}$ has degree $\le1$ in each
$z_i$, and $Y_\ell=W_{\ell-1}^\top A_\ell T_{\ell-1}$, a product of two factors each of
degree $\le1$ in $z_i$, has degree $\le2$ in each $z_i$.  Therefore
$\deg_{z_i}\phi_\ell\le2$.

\emph{Top coefficient.}  Since $w^\top X_\ell v$ is multilinear, it contains no monomial
with a squared variable; the monomial $z_1^2\cdots z_{\ell-1}^2$ can occur only in
$w^\top Y_\ell w$.  Set $a(z):=W_{\ell-1}(z)w$ and $t(z):=T_{\ell-1}(z)w$, so that
$w^\top Y_\ell w=a(z)^\top A_\ell\,t(z)$ with $a,t$ multilinear in $z_1,\dots,z_{\ell-1}$
(shown above).  By the unique-split rule (Step~0 with $k=\ell-1$), the coefficient of
$z_1^2\cdots z_{\ell-1}^2$ in $a^\top A_\ell t$ is
$a_{\mathrm{top}}^\top A_\ell\,t_{\mathrm{top}}$, where $a_{\mathrm{top}}$,
$t_{\mathrm{top}}$ are the respective coefficients of $z_1\cdots z_{\ell-1}$.  These are:
$a_{\mathrm{top}}=(-1)^{\ell-1}V_{\ell-1:1}\,w$, by Step~0(c) with $j=\ell-1$; and
$t_{\mathrm{top}}$, computed as follows: the $i'$-th summand of $T_{\ell-1}w$ involves
only $z_1,\dots,z_{i'}$, so only $i'=\ell-1$ can contribute to the coefficient of
$z_1\cdots z_{\ell-1}$; that summand is
$z_{\ell-1}B_{\ell-1:\ell}V_{\ell-1}W_{\ell-2}(z)w=z_{\ell-1}V_{\ell-1}W_{\ell-2}(z)w$
(empty product $B_{\ell-1:\ell}=I$), and by Step~0(c) with $j=\ell-2$ its coefficient of
$z_1\cdots z_{\ell-1}$ is
$V_{\ell-1}\cdot(-1)^{\ell-2}V_{\ell-2:1}\,w=(-1)^{\ell-2}V_{\ell-1:1}\,w$.
Hence the coefficient of $z_1^2\cdots z_{\ell-1}^2$ in $\phi_\ell$ equals
\[
\big[(-1)^{\ell-1}V_{\ell-1:1}w\big]^\top A_\ell\,\big[(-1)^{\ell-2}V_{\ell-1:1}w\big]
=(-1)^{2\ell-3}\,w^\top V_{\ell-1:1}^\top A_\ell V_{\ell-1:1}w
=-\,w^\top V_{\ell-1:1}^\top A_\ell V_{\ell-1:1}w,
\]
since $2\ell-3$ is odd; and $w^\top Mw=w^\top\Sym(M)w$ gives the stated symmetrized form.

\emph{(iii)}  Fix $\ell\ge2$ and write $z'=(z_1,\dots,z_{\ell-2})$.  We isolate the
dependence on $z_{\ell-1}$.  On the left factor,
\[
a(z)=W_{\ell-1}(z)w=(B_{\ell-1}-z_{\ell-1}V_{\ell-1})W_{\ell-2}(z')w
=a_0(z')+z_{\ell-1}\,a_1(z'),\quad
a_0=B_{\ell-1}W_{\ell-2}(z')w,\quad
a_1=-\,V_{\ell-1}W_{\ell-2}(z')w,
\]
since $W_{\ell-2}$ involves only $z'$.  On the right factor,
$h(z):=B_{\ell-1:1}v+T_{\ell-1}(z)w$: the term $B_{\ell-1:1}v$ is constant in $z$; the
summands $i'<\ell-1$ of $T_{\ell-1}$ involve only $z_1,\dots,z_{i'}\subseteq z'$; and the
summand $i'=\ell-1$ is $z_{\ell-1}V_{\ell-1}W_{\ell-2}(z')w$, exactly linear in
$z_{\ell-1}$ with $z_{\ell-1}$-free coefficient.  Hence
\[
h(z)=h_0(z')+z_{\ell-1}\,h_1(z'),\qquad
h_1=V_{\ell-1}W_{\ell-2}(z')w,\qquad
h_0=B_{\ell-1:1}v+\sum_{i'=1}^{\ell-2}z_{i'}B_{\ell-1:i'+1}V_{i'}W_{i'-1}(z')w .
\]
Since $a_0,a_1,h_0,h_1$ are free of $z_{\ell-1}$,
\[
\phi_\ell=a^\top A_\ell h
=a_0^\top A_\ell h_0
+z_{\ell-1}\big(a_0^\top A_\ell h_1+a_1^\top A_\ell h_0\big)
+z_{\ell-1}^2\,a_1^\top A_\ell h_1,
\]
so the coefficient of $z_{\ell-1}^2$ is
\[
a_1^\top A_\ell h_1
=-\,\big(W_{\ell-2}(z')w\big)^\top V_{\ell-1}^\top A_\ell V_{\ell-1}
\big(W_{\ell-2}(z')w\big)
=-\,w_{\ell-2}(z)^\top V_{\ell-1}^\top A_\ell V_{\ell-1}\,w_{\ell-2}(z)
=f_\ell(z'),
\]
using $w_{\ell-2}(z)=W_{\ell-2}(z')w$ from (i); in particular $f_\ell$ depends only on
$z_1,\dots,z_{\ell-2}$, as claimed.  For its top coefficient, set
$q(z'):=W_{\ell-2}(z')w$, multilinear in $z'$ by Step~0(a), so
$f_\ell=-\,q^\top\big(V_{\ell-1}^\top A_\ell V_{\ell-1}\big)q$.  By the unique-split rule
(Step~0 with $k=\ell-2$) and Step~0(c) ($q_{\mathrm{top}}=(-1)^{\ell-2}V_{\ell-2:1}w$),
the coefficient of $z_1^2\cdots z_{\ell-2}^2$ in $f_\ell$ is
\[
-\,\big[(-1)^{\ell-2}V_{\ell-2:1}w\big]^\top V_{\ell-1}^\top A_\ell V_{\ell-1}
\big[(-1)^{\ell-2}V_{\ell-2:1}w\big]
=-\,w^\top V_{\ell-1:1}^\top A_\ell V_{\ell-1:1}\,w,
\]
since $[(-1)^{\ell-2}]^2=1$ and $V_{\ell-1}V_{\ell-2:1}=V_{\ell-1:1}$; symmetrizing as
before yields the stated form, matching (ii).  (For $\ell=2$: $z'$ is the empty tuple,
$W_0=I$, $q=w$, and $f_2\equiv-\,w^\top V_1^\top A_2V_1w$ is the constant whose ``empty
top monomial'' coefficient is itself --- consistent.)

\emph{(iv)}  For $\ell\ge2$, substituting $z_2=\dots=z_{\ell-1}=1$,
\[
W_{\ell-1}(z_1,1,\dots,1)
=(B_{\ell-1}-V_{\ell-1})(B_{\ell-2}-V_{\ell-2})\cdots(B_2-V_2)\,(B_1-z_1V_1),
\]
and each factor $B_j-V_j$ ($2\le j\le\ell-1$) equals $I$ by \eqref{eq:transparency}, so
the product collapses to $B_1-z_1V_1$.

\emph{(v)}  For $m\ge1$, $g_m(z)=V_mW_{m-1}(z)w$ is multilinear in $z_1,\dots,z_{m-1}$ by
Step~0(a) (left multiplication by the constant $V_m$ does not change degrees), so
$\deg_{z_i}g_m\le1$; and by Step~0(c) with $j=m-1$ its coefficient of
$z_1\cdots z_{m-1}$ is $V_m\cdot(-1)^{m-1}V_{m-1:1}\,w=(-1)^{m-1}V_{m:1}\,w$.
(For $m=1$ this reads: the constant $g_1=V_1w=V_{1:1}w$, consistent under the
empty-monomial convention.)
\end{proof}

\begin{lemma}[Slope identity]\label{lem:lambdaphi}
Along the recursion \eqref{eq:recursion} of Proposition~\ref{prop:probe} (for any $m$-tuple
$\vartheta$, any $(w,v)\in\R^{2d}$ and any $\tau>0$), one has
$w_j=w_j(z)\big|_{z=(s_1,\dots,s_j)}$ and $v_j=v_j(z)\big|_{z=(s_1,\dots,s_j)}$ for
$0\le j\le m$; in particular
\begin{equation}\label{eq:lambdaphi}
\lambda_\ell=\phi_\ell\big(s_1,\dots,s_{\ell-1};w,v\big)\qquad(1\le\ell\le m),
\end{equation}
and $F^{(m)}_\vartheta(w,v,\tau)=v_m(z)\big|_{z=(s_1,\dots,s_m)}$ is a fixed polynomial
expression in the gates $(s_1,\dots,s_m)$ with coefficients depending only on
$(\vartheta,w,v)$.
\end{lemma}

\begin{proof}
The recursion of Proposition~\ref{prop:probe},
$w_j=(B_j-s_jV_j)w_{j-1}$, $v_j=B_jv_{j-1}+s_jV_jw_{j-1}$,
and the formal recursion defining $w_j(z),v_j(z)$ are given by the same affine maps, the
former with the numerical value $s_j$ in place of the indeterminate $z_j$.  Since the
initializations agree ($w_0=w_0(z)=w$, $v_0=v_0(z)=v$), induction on $j$ gives
$w_j=w_j(z)|_{z_i=s_i,\,1\le i\le j}$ and $v_j=v_j(z)|_{z_i=s_i,\,1\le i\le j}$ for
$0\le j\le m$.  Taking $j=\ell-1$ and substituting into the definition of $\lambda_\ell$,
\[
\lambda_\ell=w_{\ell-1}^\top A_\ell v_{\ell-1}
=\Big(w_{\ell-1}(z)^\top A_\ell\,v_{\ell-1}(z)\Big)\Big|_{z=(s_1,\dots,s_{\ell-1})}
=\phi_\ell(s_1,\dots,s_{\ell-1};w,v),
\]
and taking $j=m$ gives the last assertion via $F^{(m)}_\vartheta(w,v,\tau)=v_m$
(Proposition~\ref{prop:probe}).
\end{proof}

\begin{remark}[Real coefficients]\label{rem:real}
Since $B_j=I+V_j$, $V_j$, $A_\ell$ are real matrices and $w,v$ are real vectors, an
immediate induction on $j$ along the defining recursion shows that every component of
$w_j(z)$ and of $v_j(z)$ lies in $\R[z_1,\dots,z_j]$ (real matrices acting on
real-coefficient polynomial vectors, and multiplication by the indeterminate $z_j$,
preserve real coefficients).  Consequently $\phi_\ell\in\R[z_1,\dots,z_{\ell-1}]$ for
every $\ell$, and every component of $v_m(z)$ lies in $\R[z_1,\dots,z_m]$.  This is used
in the proof of Proposition~\ref{prop:strat}.
\end{remark}

\subsection{Probe paths and observables}\label{subsec:paths}

\begin{definition}[Path classes]\label{def:paths}
For $T_0>0$ let $\cH_{T_0}$ be the set of functions
$P=(w,v):\{\tau\in\C:|\tau|>T_0\}\to\C^d\times\C^d$ that are holomorphic and bounded, and
real-valued on $(T_0,\infty)$.  Every $P\in\cH_{T_0}$ extends holomorphically to
$\tau=\infty$ (Riemann's removable-singularity theorem, applied componentwise in the
variable $1/\tau$); write $P(\infty)\in\R^{2d}$ for the limit (real, being a limit of real
values along $(T_0,\infty)$).  For an open set $\cO\subseteq\R^{2d}$ put
\[
\cP(\cO):=\bigcup_{T_0>0}\big\{P\in\cH_{T_0}:\ P(\infty)\in\cO\big\}.
\]
Constant paths $P\equiv(w,v)$ with $(w,v)\in\cO$ belong to $\cP(\cO)$.
\end{definition}

For $P\in\cH_{T_0}$, $\theta\in\cP$ and \emph{real} $\tau>T_0$, all quantities in
\eqref{eq:recursion} evaluated at $(w(\tau),v(\tau))$ are well defined (the arguments of
$\sig$ are real), and we set
\[
G^P_\theta(\tau):=F^{(L)}_\theta\big(w(\tau),v(\tau),\tau\big),\qquad \tau\in(T_0,\infty).
\]
If $\TF_\theta=\TF_{\theta'}$ on all probe inputs $X_{w,v}(\tau)$, then
$G^P_\theta=G^P_{\theta'}$ on $(T_0,\infty)$ for every path $P$, by \eqref{eq:F} (for real
$\tau>T_0$ the point $(w(\tau),v(\tau))$ is real, so $X_{w(\tau),v(\tau)}(\tau)$ is a
genuine probe input).

\section{Setup and main results}\label{sec:statements}

Throughout, $\cP^{(m)}:=(\R^{d\times d}\times\R^{d\times d})^m$ denotes the space of
depth-$m$ parameters, and $\cP:=\cP^{(L)}$.

\subsection{Anchor sets}

\begin{definition}[Anchor sets]\label{def:anchor}
Let $\vartheta=(V_1,A_1,\dots,V_m,A_m)$ and write $q(w,v)=w^\top A_1v$,
$\pi(w,v):=A_1^\top w-A_1v$, $\Phi^\vartheta_t(w,v):=\big((B_1-tV_1)w,\ B_1v+tV_1w\big)$.
Set $\cM_1(\vartheta):=\R^{2d}$, and for $m\ge2$ let $\cM_m(\vartheta)$ be the set of
$(w,v)\in\R^{2d}$ such that:
\begin{enumerate}
\item[(M1)] $q(w,v)=0$, $A_1^\top w\neq0$ and $\pi(w,v)\neq0$;
\item[(M2)] there is $s\in(0,1)$ with $\phi^\vartheta_\ell(s,1,\dots,1;w,v)>0$ for all
$\ell=2,\dots,m$;
\item[(M3)] there is $t\in(0,1)$ with $\det(B_1-tV_1)\neq0$,
$\ \pi(w,v)^\top(B_1-tV_1)^{-1}V_1w\neq0$, and
$\Phi^\vartheta_t(w,v)\in\cM_{m-1}(\vartheta_{\ge2})$.
\end{enumerate}
\end{definition}

\subsection{Genericity sets}

\begin{definition}[Genericity]\label{def:generic}
$\cG^{(1)}:=\{(V_1,A_1):V_1\neq0,\ A_1\neq0\}$.  For $L\ge2$, $\theta'\in\cG^{(L)}$ iff
$\theta'_{\ge2}\in\cG^{(L-1)}$ and:
\begin{enumerate}
\item[(G1)] $\det A_1'\neq0$ and $\Sym(A_1')\neq0$;
\item[(G2)] $\Sym\!\big(V_{j-1:1}'^{\,\top}A_j'V_{j-1:1}'\big)\neq0$ for every
$j=2,\dots,L$, and $V_{L:1}'\neq0$;
\item[(G3)] $\det B_1'\neq0$;
\item[(G4)] $\theta'$ satisfies the anchor certificate of
Definition~\ref{def:certificate}.
\end{enumerate}
\end{definition}

Each of (G1)--(G4) is the condition that finitely many polynomials of $\theta'$ do not all
vanish (for (G4) see Auxiliary Lemma~\ref{lem:grad}); see \S\ref{sec:generic}, where we
prove (Proposition~\ref{prop:genericnonempty}) that $\cG^{(L)}$ contains the complement of
a finite union of proper algebraic subsets of $\cP$, and
(Proposition~\ref{prop:anchorsexist}) that $\cM_L(\theta')\neq\varnothing$ for every
$\theta'\in\cG^{(L)}$ provided $d\ge d_*(L):=\binom L2+2(L-1)$.  Throughout we write
$N:=d_*(L)=\binom L2+2(L-1)$ for this integer (cf.\ Definition~\ref{def:certificate};
Remark~\ref{rem:dlessN} shows the bound $d\ge N$ is intrinsic to (G4)).

\subsection{Main results}

\begin{theorem}[$\mathrm{ID}_L$]\label{thm:IDL}
Let $L\ge1$, $d\ge2$, $r\ge2$.  Let $\theta'\in\cG^{(L)}$, let $\theta\in\cP$ be arbitrary,
and let $\cO\subseteq\R^{2d}$ be open and nonempty with
$\cO\cap\cM_L(\theta')\neq\varnothing$ (no condition when $L=1$).  Suppose that for every
$P\in\cP(\cO)$ there is $T_P$ such that
\[
G^P_\theta(\tau)=G^P_{\theta'}(\tau)\qquad\text{for all real }\tau>T_P .
\]
Then $\theta=\theta'$.
\end{theorem}

Theorem~\ref{thm:IDL} is proved in \S\ref{sec:descent}, by induction on $L$, using
\S\ref{sec:A1} and \S\ref{sec:cascade} for the inductive step.

\begin{theorem}[Main theorem]\label{thm:main}
Let $L\ge1$, $r\ge2$ and $d\ge\max(2,d_*(L))$ with $d_*(L)=\binom L2+2(L-1)$.  There is a
set $\cN\subset\cP$, a finite union of proper algebraic subsets (hence closed and Lebesgue
null), such that for every $\theta'\in\cP\setminus\cN$ and every $\theta\in\cP$:
\[
\TF_\theta(X)=\TF_{\theta'}(X)\ \ \text{for all }X\in\R^{d\times(r+1)}
\quad\Longrightarrow\quad \theta=\theta'.
\]
Moreover equality is only needed on the probe inputs $X_{w,v}(\tau)$,
$(w,v,\tau)\in\R^{2d}\times(0,\infty)$, and only in the last output column.
\end{theorem}

The proof of Theorem~\ref{thm:main}, from Theorem~\ref{thm:IDL} and the results of
\S\ref{sec:generic}, is given in \S\ref{subsec:wrapup}.

\begin{lemma}[Rank factorization, cf.\ \cite{henry2025geometry}]\label{lem:lift}
Let $A,A'\in\R^{d\times \rho}$, $B,B'\in\R^{\rho\times d}$ with $AB=A'B'$ and
$\rank(AB)=\rho$.  Then there is a unique $G\in GL_\rho(\R)$ with $A'=AG$, $B'=G^{-1}B$.
\end{lemma}

\begin{proof}
\emph{Full column ranks.}  Always $\mathrm{col}(AB)\subseteq\mathrm{col}(A)$, and
$\dim\mathrm{col}(AB)=\rank(AB)=\rho$, while $\dim\mathrm{col}(A)\le\rho$ ($A$ has $\rho$
columns); hence $\dim\mathrm{col}(A)=\rho$, the $\rho$ columns of $A$ are linearly
independent ($A$ has full column rank, so $X\mapsto AX$ is injective on matrices with
$\rho$ rows), and $\mathrm{col}(A)=\mathrm{col}(AB)$.  The same argument applied to
$A'B'=AB$ gives that $A'$ has full column rank and
$\mathrm{col}(A')=\mathrm{col}(A'B')=\mathrm{col}(AB)=\mathrm{col}(A)$.

\emph{Existence and uniqueness of $G$.}  Each column of $A'$ lies in $\mathrm{col}(A)$, so
$A'=AG$ for some $G\in\R^{\rho\times\rho}$, and $G$ is unique by injectivity of
$X\mapsto AX$; explicitly $G=(A^\top A)^{-1}A^\top A'$, the matrix $A^\top A$ being
invertible by full column rank.

\emph{Invertibility of $G$.}  By symmetry of the roles there is $G'$ with $A=A'G'$.  Then
$A=AGG'$, so $A(GG'-I)=0$ and $GG'=I$; likewise $G'G=I$.  Hence $G\in GL_\rho(\R)$.

\emph{The second identity.}  $AB=A'B'=AGB'$ gives $A(B-GB')=0$, hence $B=GB'$ and
$B'=G^{-1}B$.
\end{proof}

\begin{corollary}[Factor-level identifiability]\label{cor:factors}
Assume the hypotheses of Theorem~\ref{thm:main}.  Fix $1\le d_v\le d$ and let each layer
be parametrized by
$W_O^{(\ell)}\in\R^{d\times d_v}$, $W_V^{(\ell)}\in\R^{d_v\times d}$,
$W_K^{(\ell)},W_Q^{(\ell)}\in\R^{d\times d}$, with
$V_\ell=W_O^{(\ell)}W_V^{(\ell)}$ and $A_\ell=(W_K^{(\ell)})^\top W_Q^{(\ell)}$.  There is
a finite union $\cN_{\mathrm{fac}}$ of proper algebraic subsets of the factor parameter
space such that: if the factor tuple of $\theta'$ lies outside $\cN_{\mathrm{fac}}$ and
$\TF_\theta=\TF_{\theta'}$ on $\R^{d\times(r+1)}$, then for every $\ell$ there exist unique
$G_\ell^{(v)}\in GL_{d_v}(\R)$ and $G^{(a)}_\ell\in GL_d(\R)$ with
\[
W_O^{(\ell)}=W_O'^{(\ell)}G_\ell^{(v)},\quad
W_V^{(\ell)}=(G_\ell^{(v)})^{-1}W_V'^{(\ell)},\quad
(W_K^{(\ell)})^\top=(W_K'^{(\ell)})^\top G^{(a)}_\ell,\quad
W_Q^{(\ell)}=(G^{(a)}_\ell)^{-1}W_Q'^{(\ell)} .
\]
Conversely any such change of factors leaves $\TF$ unchanged.  (There are no
permutation symmetries: layers are ordered and have one head each.)
\end{corollary}

The proof of Corollary~\ref{cor:factors}, from Theorem~\ref{thm:main},
Lemma~\ref{lem:lift} and the witnesses of \S\ref{sec:generic}, is also given in
\S\ref{subsec:wrapup}.

\section{Toolkit}\label{sec:toolkit}

\subsection{Point-set preliminaries}

We write $D(a,r):=\{z\in\C:|z-a|<r\}$.  A \emph{punctured neighbourhood} of $x$ is a set
$W\setminus\{x\}$ with $W$ a neighbourhood of $x$.  For $A\subseteq\C$ let
\[
\acc(A):=\{\tau\in\C:\ \text{every punctured neighbourhood of $\tau$ meets } A\},
\qquad \acc^{(0)}(A):=A,\quad \acc^{(k+1)}:=\acc\circ\acc^{(k)} .
\]
For $U\subseteq\C$, an \emph{accumulation point of $A$ in $U$} is a point of
$\acc(A)\cap U$.  We say $A\subseteq U$ is \emph{closed and discrete in $U$} if $A$ is
closed in the subspace $U$ and every point of $A$ is isolated in $A$.  Closures
$\clo(\cdot)$ are taken in $\C$ unless subscripted; recall the subspace identity
$\clo_U(A)=\clo(A)\cap U$, valid for any $A\subseteq U\subseteq\C$.  An empty union of
sets is $\varnothing$.

\begin{auxlemma}[Discrete sets]\label{lem:discrete}
Let $A\subseteq U\subseteq\C$.
\begin{enumerate}
\item[(i)] $A$ has no accumulation point in $U$ if and only if $A$ is closed and discrete
in $U$.
\item[(ii)] If $U$ is open and $A$ has no accumulation point in $U$, then $A\cap K$ is
finite for every compact $K\subseteq U$; in particular $A$ is countable.
\item[(iii)] If $A\subseteq B\subseteq\C$ then $\acc(A)\subseteq\acc(B)$; hence
$\acc^{(k)}(A)\subseteq\acc^{(k)}(B)$ for every $k\ge0$.
\item[(iv)] For any $A\subseteq\C$, $\clo(A)=A\cup\acc(A)$; in particular, if $A$ is
closed in $\C$ then $\acc(A)\subseteq A$.
\end{enumerate}
\end{auxlemma}

\begin{proof}
(i) ($\Leftarrow$) Suppose $A$ is closed and discrete in $U$ and, for contradiction, that
$x\in\acc(A)\cap U$.  Every neighbourhood of $x$ meets $A$, so
$x\in\clo(A)\cap U=\clo_U(A)=A$.  But $x\in A$ is isolated in $A$: there is a
neighbourhood $V$ of $x$ with $V\cap A=\{x\}$, so the punctured neighbourhood
$V\setminus\{x\}$ misses $A$, contradicting $x\in\acc(A)$.
($\Rightarrow$) Suppose $A$ has no accumulation point in $U$.  Each $a\in A\subseteq U$ is
not in $\acc(A)$, i.e.\ some punctured neighbourhood of $a$ misses $A$; hence $a$ is
isolated in $A$.  If $x\in\clo_U(A)\setminus A$, then every neighbourhood of $x$ meets
$A$, necessarily in points $\ne x$, so $x\in\acc(A)\cap U$, a contradiction; hence
$\clo_U(A)=A$, i.e.\ $A$ is closed in $U$.

(ii) Let $K\subseteq U$ be compact and suppose $A\cap K$ is infinite.  Choose pairwise
distinct $a_\nu\in A\cap K$; by compactness a subsequence converges to some
$x\in K\subseteq U$.  Since the $a_\nu$ are pairwise distinct, every punctured
neighbourhood of $x$ contains some $a_\nu\ne x$, so $x\in\acc(A)\cap U$, a contradiction.
For countability, exhaust $U$: set
$K_n:=\{z\in\C:\ |z|\le n\ \text{and}\ \dist(z,\C\setminus U)\ge 1/n\}$ (with
$\dist(z,\varnothing):=\infty$ when $U=\C$).  Each $K_n$ is closed and bounded, hence
compact, $K_n\subseteq U$, and $U=\bigcup_n K_n$ because
$\dist(z,\C\setminus U)>0$ for every $z$ in the open set $U$.  Then
$A=\bigcup_n(A\cap K_n)$ is a countable union of finite sets, hence countable.

(iii) If every punctured neighbourhood of $x$ meets $A$, it meets $B\supseteq A$; this
gives $\acc(A)\subseteq\acc(B)$, and the iterated statement follows by induction on $k$
($\acc^{(0)}$ is the identity; if $\acc^{(k)}$ is monotone then
$\acc^{(k+1)}(A)=\acc(\acc^{(k)}(A))\subseteq\acc(\acc^{(k)}(B))=\acc^{(k+1)}(B)$).

(iv) If $x\in\clo(A)\setminus A$, every neighbourhood of $x$ meets $A$ in a point
$\ne x$, so $x\in\acc(A)$; conversely $A\subseteq\clo(A)$ and
$\acc(A)\subseteq\clo(A)$ directly.  If $A$ is closed, $\acc(A)\subseteq\clo(A)=A$.
\end{proof}

\subsection{Complex-analytic lemmas}

\begin{lemma}[Sigmoid basics]\label{lem:sigmoid}
$\sig$ is meromorphic on $\C$, holomorphic exactly off $\Pi=(2\Z+1)\pi i$, with simple
poles at each point of $\Pi$.  For real $x$, $\sig(x)\in(0,1)$ and
\begin{equation}\label{eq:tail}
0<1-\sig(x)\le e^{-x},\qquad 0<\sig(x)\le e^{x}.
\end{equation}
If $\lambda\in\R\setminus\{0\}$, then $\tau\mapsto\sig(\lambda\tau+b)$ is meromorphic on
$\C$ with pole set $\big\{\frac{-b+(2n+1)\pi i}{\lambda}:n\in\Z\big\}$, which is contained
in the vertical line $\{\operatorname{Re}\tau=-b/\lambda\}$ and is closed and discrete in
$\C$.
\end{lemma}

\begin{proof}
The function $z\mapsto 1+e^{-z}$ is entire and not identically zero, and
$\sig=1/(1+e^{-z})$, so $\sig$ is meromorphic on $\C$.  Since the numerator $1$ never
vanishes, there is no cancellation: $\sig$ is holomorphic exactly off the zero set of the
denominator, and at each zero of the denominator $\sig$ has a pole whose order equals the
order of the zero.  Now $1+e^{-z}=0$ iff $e^{-z}=-1=e^{\pi i}$ iff $-z\in\pi i+2\pi i\Z$
iff $z\in(2\Z+1)\pi i=\Pi$; and $\frac{d}{dz}(1+e^{-z})=-e^{-z}\neq0$ everywhere, so each
such zero is simple.  Hence $\sig$ is holomorphic exactly on $\C\setminus\Pi$ with a
simple pole at each point of $\Pi$.

For real $x$, $e^{-x}\in(0,\infty)$, so $\sig(x)=(1+e^{-x})^{-1}\in(0,1)$.  Moreover
$1-\sig(x)=\frac{e^{-x}}{1+e^{-x}}$ and $\sig(x)=\frac{e^{x}}{1+e^{x}}$; since
$1+e^{-x}\ge1$ and $1+e^{x}\ge1$, both upper bounds in \eqref{eq:tail} follow, and the
lower bounds are strict because $e^{-x},e^{x}>0$.

Let $\lambda\in\R\setminus\{0\}$.  The map $\alpha(\tau):=\lambda\tau+b$ is an affine
automorphism of $\C$, so $\sig\circ\alpha$ is meromorphic on $\C$, holomorphic exactly at
those $\tau$ with $\alpha(\tau)\notin\Pi$, with (simple) poles exactly on
\[
P:=\alpha^{-1}(\Pi)=\Bigl\{\tfrac{-b+(2n+1)\pi i}{\lambda}:\ n\in\Z\Bigr\}.
\]
Since $b,\lambda\in\R$, every point of $P$ has real part $-b/\lambda$, so
$P\subseteq\{\operatorname{Re}\tau=-b/\lambda\}$.  Distinct points of $P$ are at distance
$2\pi|n-n'|/|\lambda|\ge 2\pi/|\lambda|$, so $P$ has no accumulation point in $\C$; by
Auxiliary Lemma~\ref{lem:discrete}(i) (with $U=\C$), $P$ is closed and discrete in $\C$.
\end{proof}

\begin{lemma}[Countable complements]\label{lem:connect}
Let $E\subset\C$ be countable.  Then $\C\setminus E$ is path-connected.  Consequently, if
$E$ is also closed, $\C\setminus E$ is a domain, and two functions holomorphic on open
sets containing $\C\setminus E$ that agree on some interval
$(T,\infty)\subset\C\setminus E$ agree on all of $\C\setminus E$.
\end{lemma}

\begin{proof}
\emph{Path-connectedness.}  Let $a,b\in\C\setminus E$.  If $a=b$ the constant path works,
so assume $a\neq b$, and let $\ell$ be the perpendicular bisector of the segment $[a,b]$,
an (uncountable) line.  Note that $a\notin\ell$ and $b\notin\ell$: points of $\ell$ are
equidistant from $a$ and $b$, while $\dist(a,a)=0<|a-b|=\dist(a,b)$, and symmetrically for
$b$.  For $z\in\ell$ let $P_z$ be the broken path $[a,z]\cup[z,b]$, and for $e\in E$ set
\[
Z_e:=\{z\in\ell:\ e\in[a,z]\cup[z,b]\}.
\]
We claim $|Z_e|\le2$.  First, $e\neq a$ and $e\neq b$ since $a,b\notin E$.  If
$e\in[a,z]$, then $e=a+s(z-a)$ for some $s\in(0,1]$ (the case $s=1$ being the corner
$e=z$), hence $z=a+s^{-1}(e-a)$ lies on the line $L_{a,e}$ through $a$ and $e$.  Since
$a\in L_{a,e}\setminus\ell$, we have $L_{a,e}\neq\ell$, and two distinct lines in the
plane meet in at most one point (in none if they are parallel); so at most one $z\in\ell$
satisfies $e\in[a,z]$.  Symmetrically (using $e\neq b$ and $b\notin\ell$), at most one
$z\in\ell$ satisfies $e\in[z,b]$.  This proves the claim.  Hence
$\bigcup_{e\in E}Z_e$ is countable, while $\ell$ is uncountable, so there exists
$z\in\ell\setminus\bigcup_{e\in E}Z_e$.  For this $z$, no point of $E$ lies on $P_z$
(interior points and the corner $z$ by the choice of $z$; the endpoints because
$a,b\notin E$), so $P_z$ joins $a$ to $b$ inside $\C\setminus E$.  Thus $\C\setminus E$ is
path-connected, hence connected.

\emph{Identity-theorem consequence.}  Suppose in addition that $E$ is closed.  Then
$\C\setminus E$ is open, nonempty (indeed co-countable in the uncountable plane) and
connected, i.e.\ a domain.  Let $F$, $G$ be holomorphic on open sets
$U_F,U_G\supseteq\C\setminus E$, with $F=G$ on a nonempty interval
$(T,\infty)\subseteq\C\setminus E$.  The restrictions of $F$ and $G$ to $\C\setminus E$
are holomorphic on this domain and agree on $(T,\infty)$, a set with an accumulation
point in $\C\setminus E$ (e.g.\ the point $T+1$, which is an accumulation point of
$(T,\infty)$ and lies in $\C\setminus E$).  By the identity theorem on the domain
$\C\setminus E$ (applied componentwise in the vector-valued case), $F=G$ on all of
$\C\setminus E$.
\end{proof}

\begin{lemma}[Pole transfer]\label{lem:transfer}
Let $E_F,E_G\subset\C$ be countable and closed, let $F,G$ be holomorphic on
$\C\setminus E_F$, $\C\setminus E_G$ respectively, and suppose $F=G$ on some
$(T,\infty)$ disjoint from $E_F\cup E_G$.  If $\tau_*$ is an isolated point of $E_G$ with
$|G(\tau)|\to\infty$ as $\tau\to\tau_*$, $\tau\notin E_G$, then $\tau_*\in E_F$.
\end{lemma}

\begin{proof}
The set $E:=E_F\cup E_G$ is countable (union of two countable sets) and closed (union of
two closed sets), so $\Omega:=\C\setminus E$ is a domain by Lemma~\ref{lem:connect}; by
hypothesis $(T,\infty)\subseteq\Omega$.  The functions $F$ and $G$ are holomorphic on the
open sets $\C\setminus E_F\supseteq\Omega$ and $\C\setminus E_G\supseteq\Omega$
respectively and agree on $(T,\infty)$; by the identity-theorem clause of
Lemma~\ref{lem:connect} (applied to this $E$), $F=G$ on $\Omega$.

Suppose, for contradiction, that $\tau_*\notin E_F$.  Since $E_F$ is closed, there is
$\delta>0$ with $\overline{D(\tau_*,\delta)}\subseteq\C\setminus E_F$, and
$M:=\sup_{\overline{D(\tau_*,\delta)}}|F|<\infty$, $F$ being continuous on a compact set.
For each $\nu\ge1$ the punctured disc $D(\tau_*,\min(\delta,1/\nu))\setminus\{\tau_*\}$ is
uncountable, hence meets $\Omega$ (its complement $E$ is countable); choose $\tau_\nu$ in
this intersection.  Then $\tau_\nu\to\tau_*$, $\tau_\nu\neq\tau_*$, and
$\tau_\nu\notin E_G$, so $|G(\tau_\nu)|\to\infty$ by hypothesis.  But
$\tau_\nu\in\Omega\cap D(\tau_*,\delta)$, so $|G(\tau_\nu)|=|F(\tau_\nu)|\le M$ for all
$\nu$, a contradiction.  Hence $\tau_*\in E_F$.  (The hypothesis that $\tau_*$ is an
isolated point of $E_G$ is not used by this proof; it is retained in the statement
because the downstream applications supply it.)
\end{proof}

\begin{lemma}[Accumulation of $\Pi$-preimages at a pole]\label{lem:T4}
Let $D$ be a disc around $\xi$ and $H$ holomorphic on $D\setminus\{\xi\}$ with
$|H(\tau)|\to\infty$ as $\tau\to\xi$.  Then for every $\varepsilon>0$ the set
$\{\tau:0<|\tau-\xi|<\varepsilon,\ H(\tau)\in\Pi\}$ is infinite.
\end{lemma}

\begin{proof}
\emph{Step 1: $\xi$ is a pole of $H$ of finite order.}  Since $|H(\tau)|\to\infty$ as
$\tau\to\xi$, there is $r_0>0$ with $D(\xi,r_0)\subseteq D$ and $|H|\ge1$ on
$D(\xi,r_0)\setminus\{\xi\}$.  There $H$ has no zeros, so $\eta:=1/H$ is holomorphic on
$D(\xi,r_0)\setminus\{\xi\}$ and bounded by $1$; by Riemann's removable-singularity
theorem, $\eta$ extends holomorphically to $D(\xi,r_0)$, with
$\eta(\xi)=\lim_{\tau\to\xi}1/H(\tau)=0$.  The extension is not identically zero (it is
nonzero off $\xi$), so its zero at $\xi$ has finite order $m\ge1$:
$\eta(\tau)=(\tau-\xi)^m g(\tau)$ with $g$ holomorphic on $D(\xi,r_0)$ and
$g(\xi)\neq0$.  Choose $r_1\in(0,r_0]$ with $g\neq0$ on $D(\xi,r_1)$ and set $h:=1/g$,
holomorphic and zero-free on $D(\xi,r_1)$; then
\[
H(\tau)=h(\tau)\,(\tau-\xi)^{-m}\qquad(0<|\tau-\xi|<r_1).
\]

\emph{Step 2: solving $H=\lambda$ near $\xi$ via Rouch\'e.}  Fix $\varepsilon>0$ and set
$\rho_0:=\min(\varepsilon,r_1/2)$, so that
$\overline{D(\xi,\rho_0)}\subset D(\xi,r_1)$ and $\rho_0\le\varepsilon$.  Let
$M:=\max_{|\tau-\xi|\le\rho_0}|h(\tau)|$; then $0<M<\infty$ (continuity on a compact set;
positivity since $h$ is zero-free).  Let $\lambda\in\C$ with $|\lambda|>2M\rho_0^{-m}$ and
put $\rho_\lambda:=(2M/|\lambda|)^{1/m}\in(0,\rho_0)$.  On the circle
$|\tau-\xi|=\rho_\lambda$,
\[
|\lambda(\tau-\xi)^m|=|\lambda|\,\rho_\lambda^{\,m}=2M>M\ge|h(\tau)|,
\]
so by Rouch\'e's theorem (both functions are holomorphic on a neighbourhood of
$\overline{D(\xi,\rho_\lambda)}$, and the dominant one is zero-free on the circle),
$\lambda(\tau-\xi)^m-h(\tau)$ has the same number of zeros, with multiplicity, in
$D(\xi,\rho_\lambda)$ as $\lambda(\tau-\xi)^m$, namely $m\ge1$.  Any such zero $\hat\tau$
satisfies $\hat\tau\neq\xi$, because the value at $\xi$ is $-h(\xi)\neq0$; hence
$0<|\hat\tau-\xi|<\rho_\lambda<\rho_0\le\varepsilon$ and
$H(\hat\tau)=h(\hat\tau)(\hat\tau-\xi)^{-m}=\lambda$.

\emph{Step 3: infinitely many $\Pi$-preimages.}  Apply Step 2 to
$\lambda_n:=(2n+1)\pi i\in\Pi$ for every $n\in\Z$ with
$|\lambda_n|=(2|n|+1)\pi>2M\rho_0^{-m}$ --- there are infinitely many such $n$ ---
obtaining for each one a point $\tau_n$ with $0<|\tau_n-\xi|<\varepsilon$ and
$H(\tau_n)=\lambda_n$.  For $n\neq n'$ we have
$H(\tau_n)=\lambda_n\neq\lambda_{n'}=H(\tau_{n'})$, so $\tau_n\neq\tau_{n'}$: the points
are pairwise distinct.  Hence
$\{\tau:0<|\tau-\xi|<\varepsilon,\ H(\tau)\in\Pi\}$ contains the infinite set
$\{\tau_n\}$, as claimed.
\end{proof}

\begin{lemma}[Discrete level sets]\label{lem:levelsets}
Let $\Omega\subseteq\C$ be a domain containing an interval $(T,\infty)$, and $H$
holomorphic on $\Omega$ and real on $(T,\infty)$.  Then $S:=H^{-1}(\Pi)$ is closed in
$\Omega$ and has no accumulation points in $\Omega$.
\end{lemma}

\begin{proof}
The interval $(T,\infty)$ is nonempty; fix $\tau_0\in(T,\infty)$.

If $H$ is constant, then $H\equiv H(\tau_0)\in\R$, and since $\Pi\cap\R=\varnothing$ we
get $S=\varnothing$, which is closed in $\Omega$ and has no accumulation points.

Assume $H$ is nonconstant.  The set $S=H^{-1}(\Pi)$ is closed in $\Omega$ as the preimage
of the closed set $\Pi$ under the continuous map $H:\Omega\to\C$.  Suppose, for
contradiction, that $S$ has an accumulation point $\tau_*\in\Omega$.  Choose pairwise
distinct $\tau_\nu\in S\setminus\{\tau_*\}$ with $\tau_\nu\to\tau_*$: this is possible
since every punctured disc $D(\tau_*,\rho)\setminus\{\tau_*\}$ meets $S$, so one may pick
the $\tau_\nu$ recursively with radii $\rho_\nu<\min(1/\nu,\,|\tau_{\nu-1}-\tau_*|)$,
forcing distinctness.  By continuity $H(\tau_\nu)\to H(\tau_*)$, and $H(\tau_*)\in\Pi$
since $\Pi$ is closed.  Distinct points of $\Pi$ are at distance $\ge2\pi$, so
$D(H(\tau_*),\pi)\cap\Pi=\{H(\tau_*)\}$; for all large $\nu$, $H(\tau_\nu)$ lies in this
intersection, hence $H(\tau_\nu)=H(\tau_*)=:(2n+1)\pi i$.  Thus the zero set of the
holomorphic function $H-(2n+1)\pi i$ contains the pairwise distinct points $\tau_\nu$
($\nu$ large), which accumulate at $\tau_*\in\Omega$.  Since $\Omega$ is a domain, the
identity theorem yields $H\equiv(2n+1)\pi i$ on $\Omega$.  But then
$H(\tau_0)=(2n+1)\pi i\notin\R$, contradicting realness of $H$ on $(T,\infty)$.  Hence
$S$ has no accumulation point in $\Omega$.

By Auxiliary Lemma~\ref{lem:discrete}(i), the conclusion is equivalent to: $S$ is closed
and discrete in $\Omega$ --- the form in which it is consumed by
Lemma~\ref{lem:strataacc} and by Proposition~\ref{prop:strat}.
\end{proof}

\begin{lemma}[Stratified accumulation]\label{lem:strataacc}
Let $S=S^1\cup\dots\cup S^m\subset\C$, where $S^1$ is closed and discrete in $\C$ and, for
$j\ge2$, $S^j\subseteq\C\setminus(S^1\cup\dots\cup S^{j-1})$ is closed and discrete in
$\C\setminus(S^1\cup\dots\cup S^{j-1})$.  Then each partial union
$S^1\cup\dots\cup S^j$ is closed in $\C$, and
\[
\acc^{(k)}(S)\subseteq S^1\cup\dots\cup S^{m-k}\quad(0\le k\le m),\qquad
\acc^{(m)}(S)=\varnothing .
\]
\end{lemma}

\begin{proof}
Write $T_j:=S^1\cup\dots\cup S^j$ ($T_0:=\varnothing$) and $U_j:=\C\setminus T_{j-1}$, so
that the hypotheses read: for $1\le j\le m$, $S^j\subseteq U_j$ is closed and discrete in
$U_j$ (with $U_1=\C$); equivalently, by Auxiliary Lemma~\ref{lem:discrete}(i), $S^j$ has
no accumulation point in $U_j$.

\emph{Step 1: each $T_j$ is closed in $\C$.}  Induction on $j$.  For $j\le1$:
$T_0=\varnothing$ and $T_1=S^1$ is closed in $\C$ by hypothesis.  Let $j\ge2$ and assume
$T_{j-1}$ is closed, so $U_j$ is open.  Since
$\clo(T_j)\subseteq\clo(T_{j-1})\cup\clo(S^j)=T_{j-1}\cup\clo(S^j)$, let
$x\in\clo(T_j)$.  If $x\in T_{j-1}$, then $x\in T_j$.  Otherwise
$x\in\clo(S^j)\cap U_j=\clo_{U_j}(S^j)=S^j\subseteq T_j$, where we used the subspace
identity $\clo_{U_j}(S^j)=\clo(S^j)\cap U_j$ (valid because $S^j\subseteq U_j$) and the
closedness of $S^j$ in $U_j$.  Hence $\clo(T_j)=T_j$.  In particular, every
$U_j=\C\setminus T_{j-1}$ is open in $\C$.

\emph{Step 2: $\acc^{(k)}(S)\subseteq T_{m-k}$ for $0\le k\le m$.}  We prove this by
induction on the number of strata $m\ge1$, for all systems satisfying the hypotheses.

\emph{Base $m=1$.}  The case $k=0$ is trivial ($\acc^{(0)}(S)=S=T_1$).  For $k=1$:
$\acc(S^1)\cap\C=\varnothing$ by Auxiliary Lemma~\ref{lem:discrete}(i) applied with
$U=\C$ ($S^1$ closed and discrete in $\C$ means no accumulation points in $\C$), i.e.\
$\acc(S)=\varnothing=T_0$.

\emph{Step $m\ge2$.}  We first prove the one-step inclusion
$\acc(S)\subseteq T_{m-1}$.  Let $x\in\acc(S)$ and suppose $x\notin T_{m-1}$.  Then
$x\in U_m=\C\setminus T_{m-1}$, which is open by Step~1.  Let $W\setminus\{x\}$ be any
punctured neighbourhood of $x$.  Then $(W\cap U_m)\setminus\{x\}$ is again a punctured
neighbourhood of $x$ (as $U_m$ is an open neighbourhood of $x$), so it meets $S$; and
$S\cap U_m=S^m$, since $T_{m-1}\cap U_m=\varnothing$ and $S^m\subseteq U_m$.  Hence every
punctured neighbourhood of $x$ meets $S^m$, i.e.\ $x\in\acc(S^m)\cap U_m$ ---
contradicting the hypothesis that $S^m$ has no accumulation point in $U_m$.  (Note that
the discreteness used here for $S^m$ is the \emph{relative} one, in $U_m$; the
\emph{absolute} discreteness in $\C$ is used only for $S^1$, in the base case.)

The truncated system $S^1,\dots,S^{m-1}$ satisfies the hypotheses of the lemma with $m-1$
strata --- its hypotheses are literally a sublist of the given ones --- so by the
inductive hypothesis,
\[
\acc^{(k)}(T_{m-1})\subseteq T_{(m-1)-k}\qquad(0\le k\le m-1).
\]
For $1\le k\le m$, monotonicity of $\acc^{(k-1)}$ (Auxiliary
Lemma~\ref{lem:discrete}(iii)) applied to the one-step inclusion
$\acc(S)\subseteq T_{m-1}$ gives
\[
\acc^{(k)}(S)=\acc^{(k-1)}\bigl(\acc(S)\bigr)\subseteq\acc^{(k-1)}(T_{m-1})
\subseteq T_{(m-1)-(k-1)}=T_{m-k}.
\]
Together with the trivial case $k=0$ this proves the inclusion for all $0\le k\le m$; for
$k=m$ it reads $\acc^{(m)}(S)\subseteq T_0=\varnothing$, i.e.\
$\acc^{(m)}(S)=\varnothing$.
\end{proof}

\subsection{The singular stratification}

\begin{proposition}[Singular stratification of a constant-probe observable]
\label{prop:strat}
Let $\vartheta$ be any $m$-layer parameter and fix $(w,v)\in\R^{2d}$.  There are sets
$S^1,\dots,S^m\subseteq\C$ with the structure of Lemma~\ref{lem:strataacc} such that,
writing $\Omega_j:=\C\setminus(S^1\cup\dots\cup S^{j})$:
\begin{enumerate}
\item[(a)] $(S^1\cup\dots\cup S^m)\cap\R=\varnothing$;
\item[(b)] $s_1$ extends to a meromorphic function on $\C$ with pole set $S^1$; for
$j\ge2$, $H_j(\tau):=\tau\,\phi_j(s_1(\tau),\dots,s_{j-1}(\tau);w,v)+b$ is holomorphic on
$\Omega_{j-1}$, $S^j=H_j^{-1}(\Pi)\subseteq\Omega_{j-1}$, and
$s_j:=\sig(H_j)$ is holomorphic on $\Omega_j$;
\item[(c)] $\tau\mapsto F^{(m)}_\vartheta(w,v,\tau)$ extends holomorphically from
$(0,\infty)$ to $\Omega_m=\C\setminus S_\vartheta$, where
$S_\vartheta:=S^1\cup\dots\cup S^m$ is countable and closed;
\item[(d)] if $\lambda_1:=w^\top A_1v\neq0$ then
$S^1=\{(-b+(2n+1)\pi i)/\lambda_1:n\in\Z\}\subseteq\{\operatorname{Re}\tau=-b/\lambda_1\}$,
and if $\lambda_1=0$ then $S^1=\varnothing$;
\item[(e)] if $V_1=0$, then $F^{(m)}_\vartheta(w,v,\cdot)=F^{(m-1)}_{\vartheta_{\ge2}}
(w,v,\cdot)$, which admits the same structure with $m-1$ strata.
\end{enumerate}
\end{proposition}

\begin{proof}
Throughout, $\vartheta=(V_1,A_1,\dots,V_m,A_m)$ consists of real matrices and
$(w,v)\in\R^{2d}$ is real.  Recall from Remark~\ref{rem:real} that
$\phi_i\in\R[z_1,\dots,z_{i-1}]$ for every $i$ and that the components of $v_m(z)$ lie in
$\R[z_1,\dots,z_m]$, and recall $\Pi\cap\R=\varnothing$.  Set
$T_j:=S^1\cup\dots\cup S^j$ with $T_0:=\varnothing$, $\Omega_j:=\C\setminus T_j$ with
$\Omega_0:=\C$, and $\lambda_1:=w^\top A_1v=\phi_1(w,v)\in\R$.

We construct the sets $S^j$ and holomorphic functions $s_j$ (denoted by the same symbol
as the recursion gates of Proposition~\ref{prop:probe}, which they will be shown to
extend) by induction on $j\in\{1,\dots,m\}$, maintaining the following hypotheses:
\begin{itemize}
\item[(I1)$_j$] (\emph{Stratification})  $S^1$ is closed and discrete in $\C$; for
$2\le i\le j$, $S^i\subseteq\Omega_{i-1}$ and $S^i$ has no accumulation point in
$\Omega_{i-1}$ --- equivalently (Auxiliary Lemma~\ref{lem:discrete}(i)), $S^i$ is closed
and discrete in $\Omega_{i-1}$.
\item[(I2)$_j$] (\emph{Topology})  For $1\le i\le j$: $S^i$ is countable and
$S^i\cap\R=\varnothing$; the partial union $T_i$ is countable, closed in $\C$, and
$T_i\cap\R=\varnothing$; and $\Omega_i$ is a domain with $\R\subseteq\Omega_i$.
\item[(I3)$_j$] (\emph{Holomorphy})  $s_1$ is meromorphic on $\C$ with pole set exactly
$S^1$ and is holomorphic on $\Omega_1$; for $2\le i\le j$,
$H_i(\tau):=\tau\,\phi_i(s_1(\tau),\dots,s_{i-1}(\tau);w,v)+b$ is holomorphic on
$\Omega_{i-1}$, $S^i=H_i^{-1}(\Pi)\subseteq\Omega_{i-1}$, and $s_i:=\sig\circ H_i$ is
holomorphic on $\Omega_i$.
\item[(I4)$_j$] (\emph{Realness})  For every $\tau\in\R$: $H_i(\tau)\in\R$ ($2\le i\le j$)
and $s_i(\tau)\in(0,1)$ ($1\le i\le j$).
\item[(I5)$_j$] (\emph{Compatibility})  For every $\tau\in(0,\infty)$ and $1\le i\le j$,
$s_i(\tau)$ equals the gate value $s_i$ produced by the recursion of
Proposition~\ref{prop:probe} at inverse temperature $\tau$.
\end{itemize}
The point of carrying (I2)$_j$ along the induction is precisely that the domain property
of $\Omega_{j-1}$, needed \emph{before} Lemma~\ref{lem:levelsets} can be applied at stage
$j$, requires $T_{j-1}$ to be countable and closed.

\emph{Base $j=1$.}  Define $s_1(\tau):=\sig(\lambda_1\tau+b)$ on $\C$.  If
$\lambda_1\neq0$, Lemma~\ref{lem:sigmoid} shows that $s_1$ is meromorphic on $\C$ with
pole set
\[
S^1:=\Bigl\{\tfrac{-b+(2n+1)\pi i}{\lambda_1}:n\in\Z\Bigr\}
\subseteq\{\operatorname{Re}\tau=-b/\lambda_1\},
\]
closed and discrete in $\C$; it is countable (indexed by $\Z$) and disjoint from $\R$,
since each of its points has imaginary part $(2n+1)\pi/\lambda_1\neq0$.  If
$\lambda_1=0$, then $s_1\equiv\sig(b)=r/(r+1)\in(0,1)$ is entire and we set
$S^1:=\varnothing$.  In both cases (I1)$_1$ holds; $T_1=S^1$ is countable, closed, and
avoids $\R$; and $\Omega_1=\C\setminus S^1$ is a domain by Lemma~\ref{lem:connect}
($S^1$ countable and closed; if $S^1=\varnothing$ then $\Omega_1=\C$) containing $\R$ ---
this is (I2)$_1$.  By Lemma~\ref{lem:sigmoid}, $s_1$ is holomorphic exactly off its pole
set, in particular on $\Omega_1$: (I3)$_1$.  For $\tau\in\R$, $\lambda_1\tau+b\in\R$, so
$s_1(\tau)=\sig(\lambda_1\tau+b)\in(0,1)$: (I4)$_1$.  For $\tau>0$,
Proposition~\ref{prop:probe} defines the first gate as
$\sig(\tau\lambda_1+b)=s_1(\tau)$: (I5)$_1$.  This base case also proves item (d).

\emph{Inductive step $j\ge2$.}  Assume (I1)$_{j-1}$--(I5)$_{j-1}$.  Since
$T_i\subseteq T_{j-1}$, we have $\Omega_{j-1}\subseteq\Omega_i$ for every $i\le j-1$, so
each $s_i$ ($1\le i\le j-1$) is holomorphic on $\Omega_{j-1}$.  Hence
\[
H_j(\tau):=\tau\,\phi_j\bigl(s_1(\tau),\dots,s_{j-1}(\tau);w,v\bigr)+b
\]
is holomorphic on $\Omega_{j-1}$, being a polynomial expression with constant
coefficients in $\tau$ and finitely many holomorphic functions.  For
$\tau\in\R\subseteq\Omega_{j-1}$ (the inclusion by (I2)$_{j-1}$): the arguments satisfy
$s_i(\tau)\in(0,1)\subseteq\R$ by (I4)$_{j-1}$, the polynomial $\phi_j$ has real
coefficients (Remark~\ref{rem:real}), and $\tau,b\in\R$; hence $H_j(\tau)\in\R$.  In
particular $H_j$ is real on $(0,\infty)$.

By (I2)$_{j-1}$, $\Omega_{j-1}$ is a domain containing the interval $(0,\infty)$, and
$H_j$ is holomorphic on it and real on $(0,\infty)$.  Lemma~\ref{lem:levelsets} (with
$T=0$) therefore yields: $S^j:=H_j^{-1}(\Pi)\subseteq\Omega_{j-1}$ is closed in
$\Omega_{j-1}$ and has no accumulation point in $\Omega_{j-1}$ --- equivalently, $S^j$ is
closed and discrete in $\Omega_{j-1}$ (Auxiliary Lemma~\ref{lem:discrete}(i)).  Together
with (I1)$_{j-1}$ this establishes (I1)$_j$, which is exactly the hypothesis system of
Lemma~\ref{lem:strataacc} for the strata $S^1,\dots,S^j$.  Moreover
$S^j\cap\R=\varnothing$, since $H_j(\R)\subseteq\R$ and $\Pi\cap\R=\varnothing$; and
$S^j$ is countable by Auxiliary Lemma~\ref{lem:discrete}(ii), being closed and discrete
in the open set $\Omega_{j-1}$.

Consequently $T_j=T_{j-1}\cup S^j$ is countable and avoids $\R$, and it is closed in $\C$
by the first assertion (Step~1) of Lemma~\ref{lem:strataacc}, applied to the system
$S^1,\dots,S^j$, whose hypotheses are (I1)$_j$, just verified.  Hence
$\Omega_j=\C\setminus T_j$ is open and, by Lemma~\ref{lem:connect} ($T_j$ countable and
closed), a domain; it contains $\R$ because $T_j\cap\R=\varnothing$.  This is (I2)$_j$.
Note that no fact about $T_j$ was used before being established: the domain property of
$\Omega_{j-1}$ came from (I2)$_{j-1}$, and that of $\Omega_j$ is derived only now.

Define $s_j:=\sig\circ H_j$.  By Lemma~\ref{lem:sigmoid}, $\sig$ is holomorphic on
$\C\setminus\Pi$, and $H_j$ maps $\Omega_j=\Omega_{j-1}\setminus S^j$ into
$\C\setminus\Pi$ by the definition of $S^j$; hence $s_j$ is holomorphic on $\Omega_j$:
(I3)$_j$.  For $\tau\in\R$ we have $H_j(\tau)\in\R$, so
$s_j(\tau)=\sig(H_j(\tau))\in(0,1)$: (I4)$_j$.  Finally, fix $\tau>0$ and let
$s_1,\dots,s_m$ momentarily denote the gates of Proposition~\ref{prop:probe} at this
$\tau$; by (I5)$_{j-1}$ the constructed functions agree with the gates
$s_1,\dots,s_{j-1}$ at $\tau$, so by \eqref{eq:lambdaphi},
\[
H_j(\tau)=\tau\,\phi_j(s_1,\dots,s_{j-1};w,v)+b=\tau\lambda_j+b,
\]
whence $s_j(\tau)=\sig(\tau\lambda_j+b)$ is exactly the $j$-th gate of the recursion:
(I5)$_j$.  This completes the induction.

\emph{Conclusion.}  Take $j=m$.  By (I1)$_m$ (and Auxiliary
Lemma~\ref{lem:discrete}(i)), the sets $S^1,\dots,S^m$ have exactly the structure
required by Lemma~\ref{lem:strataacc}.  Item (a) is the statement
$T_m\cap\R=\varnothing$ contained in (I2)$_m$ --- it holds on all of $\R$, not merely on
$(0,\infty)$, because (I4) gives realness of every $H_i$ and $s_i$ on the whole real line
and the first stratum avoids $\R$ by its explicit form.  Item (b) is (I3)$_m$, combined
with (I5)$_m$: the meromorphic function $s_1$ coincides on $(0,\infty)$ with the first
recursion gate, hence extends it, and likewise each $s_j$ extends the $j$-th gate from
$(0,\infty)$ to $\Omega_j$.  Item (d) was proved in the base step.

\emph{(c).}  By Proposition~\ref{prop:probe} and Lemma~\ref{lem:lambdaphi}, for every
$\tau>0$,
\[
F^{(m)}_\vartheta(w,v,\tau)=v_m=P\bigl(s_1,\dots,s_m\bigr),\qquad P(z):=v_m(z),
\]
where $P:\C^m\to\C^d$ is a fixed polynomial map with constant
($\vartheta,w,v$-dependent) coefficients and the arguments are the gates at $\tau$; by
(I5)$_m$ these equal $s_1(\tau),\dots,s_m(\tau)$.  Define
\[
\widehat F(\tau):=P\bigl(s_1(\tau),\dots,s_m(\tau)\bigr),\qquad \tau\in\Omega_m .
\]
Each $s_i$ is holomorphic on $\Omega_i\supseteq\Omega_m$, so every component of
$\widehat F$ is holomorphic on $\Omega_m$, and
$\widehat F=F^{(m)}_\vartheta(w,v,\cdot)$ on $(0,\infty)\subseteq\Omega_m$.  By
(I2)$_m$, $S_\vartheta=T_m$ is countable and closed in $\C$, so
$\Omega_m=\C\setminus S_\vartheta$ is a domain (Lemma~\ref{lem:connect}).  Thus
$\widehat F$ is a holomorphic extension of $F^{(m)}_\vartheta(w,v,\cdot)$ from
$(0,\infty)$ to $\Omega_m$; moreover it is the unique one, since any two holomorphic
extensions agree on $(0,\infty)$ and hence on all of $\C\setminus S_\vartheta$ by the
identity-theorem clause of Lemma~\ref{lem:connect}.  This proves (c).

\emph{(e).}  If $V_1=0$ then $B_1=I$, so by the parenthetical statement of
Proposition~\ref{prop:probe} we have $w_1=w$, $v_1=v$, the gate $s_1$ drops out of the
recursion, and $F^{(m)}_\vartheta(w,v,\cdot)=F^{(m-1)}_{\vartheta_{\ge2}}(w,v,\cdot)$ on
$(0,\infty)$.  The construction above is uniform in the layer tuple and in the number of
layers, and the proofs of items (a)--(d) make no use of item (e); applying the
construction to the $(m-1)$-tuple $\vartheta_{\ge2}=(V_2,A_2,\dots,V_m,A_m)$ and the same
$(w,v)$ therefore produces strata $\widetilde S^1,\dots,\widetilde S^{m-1}$ with the
structure of Lemma~\ref{lem:strataacc} and properties (a)--(d) for
$F^{(m-1)}_{\vartheta_{\ge2}}(w,v,\cdot)$, the first stratum now governed by
$\widetilde\lambda_1=w^\top A_2v$; no quantity involving $s_1$ enters, since $w_1=w$ and
$v_1=v$ hold identically.  (Degenerate case $m=1$: $F^{(0)}(w,v,\cdot)\equiv v$ is entire
and the structure is the empty stratification, all assertions holding vacuously.)  This
proves (e) and completes the proof.
\end{proof}

\subsection{Algebraic lemmas}

\begin{lemma}[Zero sets of polynomials]\label{lem:zariski}
Let $p\not\equiv0$ be a polynomial on $\R^n$.  Then $\{p=0\}$ is closed, has empty
interior and Lebesgue measure zero; in particular $\{p\neq0\}$ is open and dense, and a
polynomial vanishing on a nonempty open subset of $\R^n$ vanishes identically.
Consequently finitely many nonzero polynomials have a common nonvanishing point in any
nonempty open subset of $\R^n$.  Finally, a finite union of proper linear subspaces of
$\R^d$ has empty interior.
\end{lemma}

\begin{proof}
\emph{Closedness.}  $\{p=0\}$ is the preimage of the closed set $\{0\}$ under the
continuous map $p$.

\emph{Empty interior.}  Suppose $p$ vanishes on a ball $B(x_0,\varepsilon)$.  A
polynomial equals its (finite) Taylor expansion at $x_0$; vanishing on a neighbourhood of
$x_0$ forces all partial derivatives of $p$ at $x_0$ to vanish, hence all Taylor
coefficients, hence $p\equiv0$ --- contradiction.  The contrapositive is the assertion
that a polynomial vanishing on a nonempty open set is identically zero, and density of
$\{p\neq0\}$ is the same statement.

\emph{Measure zero.}  Induction on $n$.  For $n=1$, a nonzero polynomial has at most
$\deg p$ roots, a finite (null) set.  For $n\ge2$, write
$p(x',x_n)=\sum_{k\le D}c_k(x')\,x_n^k$ with $x'\in\R^{n-1}$ and not all $c_k\equiv0$;
fix $k_0$ with $c_{k_0}\not\equiv0$.  Let $A:=\{p=0\}$ (a closed, hence Borel, set) and
$Z':=\{x'\in\R^{n-1}:\ p(x',\cdot)\equiv0\}=\{x':\ c_k(x')=0\ \forall k\}\subseteq
\{c_{k_0}=0\}$, which is Lebesgue null in $\R^{n-1}$ by the inductive hypothesis.  For
$x'\notin Z'$ the slice $A_{x'}=\{x_n:\ p(x',x_n)=0\}$ is finite, hence null in $\R$.  By
Fubini--Tonelli applied to the indicator of $A$,
\[
\lambda_n(A)=\int_{\R^{n-1}}\lambda_1(A_{x'})\,dx'
=\int_{Z'}\lambda_1(A_{x'})\,dx'+\int_{\R^{n-1}\setminus Z'}0\,dx'=0,
\]
since $\lambda_{n-1}(Z')=0$.

\emph{Finitely many polynomials.}  If $p_1,\dots,p_k\not\equiv0$, then
$p_1\cdots p_k\not\equiv0$ ($\R[x_1,\dots,x_n]$ is an integral domain); its nonvanishing
set is open and dense, hence meets any nonempty open set, and at any such point all
$p_i\neq0$.

\emph{Union of proper subspaces.}  Each proper linear subspace $V_i\subsetneq\R^d$ is
contained in a hyperplane $\{\ell_i=0\}$ for some nonzero linear form $\ell_i$ (extend a
basis of $V_i$ and take a coordinate functional vanishing on $V_i$).  The union is
contained in $\{\ell_1\cdots\ell_k=0\}$, the zero set of a nonzero polynomial, which has
empty interior by the above.
\end{proof}

\begin{definition}\label{def:dtc}
A polynomial $p\in\C[z_1,\dots,z_K]$ has a \emph{dominant top coefficient} if, with
$D_i:=\deg_{z_i}p$, the coefficient of $\prod_iz_i^{D_i}$ in $p$ is nonzero.  (If $p$ does
not involve $z_i$ then $D_i=0$; in particular every nonzero constant has a dominant top
coefficient, and a polynomial with a dominant top coefficient is nonzero.)
\end{definition}

\begin{lemma}[Nested largeness with continuous thresholds]\label{lem:nested}
Let $p_1,\dots,p_M\in\C[z_1,\dots,z_K]$ each have a dominant top coefficient.  Then there
exist $R_1>0$ and continuous functions $R_k:\cN_{k-1}\to(0,\infty)$ for
$2\le k\le K$, where $\cN_1:=\{z_1\in\C:|z_1|>R_1\}$ and
$\cN_k:=\{(z',z_k):z'\in\cN_{k-1},\ |z_k|>R_k(z')\}$, such that for every
$1\le k\le K$ and every index $a$ with $p_a\in\C[z_1,\dots,z_k]$ one has
$p_a(z)\neq0$ for all $z\in\cN_k$.
\end{lemma}

\begin{proof}
\emph{Step 0: products of polynomials with dominant top coefficients again have dominant
top coefficients.}  Let $p,q$ have dominant top coefficients, with per-variable degrees
$D_i,E_i$ and top coefficients $c_p,c_q\neq0$.  Viewing $p,q$ as one-variable polynomials
in $z_i$ over the integral domain $\C[z_1,\dots,\widehat{z_i},\dots,z_K]$, their leading
coefficients are nonzero and multiply, so $\deg_{z_i}(pq)=D_i+E_i$ for every $i$.  The
coefficient of $\prod_iz_i^{D_i+E_i}$ in $pq$ is
$\sum a_{m_1}b_{m_2}$ over factorizations $m_1m_2=\prod_iz_i^{D_i+E_i}$ with
$m_1\in\operatorname{supp}p$, $m_2\in\operatorname{supp}q$; since
$\deg_{z_i}m_1\le D_i$ and $\deg_{z_i}m_2\le E_i$ with sum $D_i+E_i$, equality is forced
in each variable, i.e.\ $m_1=\prod z_i^{D_i}$, $m_2=\prod z_i^{E_i}$ --- a unique split
--- so the coefficient equals $c_pc_q\neq0$.  Iterating, $P:=p_1\cdots p_M$ has a
dominant top coefficient.

\emph{Step 1: iterated leading coefficients.}  For a polynomial $p$ with dominant top
coefficient and $K\ge1$, let $\operatorname{lc}_{z_K}(p)\in\C[z_1,\dots,z_{K-1}]$ denote
the coefficient of $z_K^{\deg_{z_K}p}$.  Then $\operatorname{lc}_{z_K}(p)$ again has a
dominant top coefficient, with the \emph{same} per-variable degrees $D_i$ ($i<K$) and the
\emph{same} top coefficient: indeed the coefficient of $\prod_{i<K}z_i^{D_i}$ in
$\operatorname{lc}_{z_K}(p)$ equals the top coefficient of $p$, which is nonzero; this
forces $\deg_{z_i}\operatorname{lc}_{z_K}(p)\ge D_i$ for $i<K$, while trivially $\le D_i$.
Define the iterated leading coefficients $P_K:=P$ and
$P_k:=\operatorname{lc}_{z_{k+1}}(P_{k+1})\in\C[z_1,\dots,z_k]$ for $k=K-1,\dots,1$; each
$P_k$ has a dominant top coefficient.  Moreover, leading coefficients are multiplicative
over an integral domain:
$\operatorname{lc}_{z_K}(pq)=\operatorname{lc}_{z_K}(p)\operatorname{lc}_{z_K}(q)$, and if
$p$ does not involve $z_K$ then $\operatorname{lc}_{z_K}(p)=p$.  Iterating, for every
$1\le k\le K$,
\[
P_k=\prod_{a=1}^M p_a^{(k)},\qquad
p_a^{(k)}:=\text{the iterated leading coefficient of }p_a\text{ down to level }k,
\]
and $p_a^{(k)}=p_a$ whenever $p_a\in\C[z_1,\dots,z_k]$.  In particular, at every point
where $P_k\neq0$, each factor is nonzero; so $p_a\neq0$ wherever $P_k\neq0$, for each
qualifying $a$.

\emph{Step 2: construction of the regions.}  We prove, by induction on $K\ge1$: for any
$Q\in\C[z_1,\dots,z_K]$ with a dominant top coefficient there are regions
$\cN_1,\dots,\cN_K$ with thresholds as in the statement such that the iterated leading
coefficients $Q_k$ of $Q$ satisfy $Q_k\neq0$ everywhere on $\cN_k$, for every
$1\le k\le K$.

For $K=1$: $Q\in\C[z_1]$ is nonzero, hence has finitely many roots; take $R_1$ larger
than the modulus of every root (and $\ge1$), so $Q=Q_1\neq0$ on
$\cN_1=\{|z_1|>R_1\}$.

For $K\ge2$: write $Q=\sum_{i\le D}c_i(z')z_K^i$ with $D:=\deg_{z_K}Q$ and
$c_D=\operatorname{lc}_{z_K}(Q)=Q_{K-1}$.  By Step~1, $c_D$ has a dominant top
coefficient, and the iterated leading coefficients of $c_D$ are exactly
$Q_1,\dots,Q_{K-1}$.  By the inductive hypothesis applied to $c_D$, obtain
$\cN_1,\dots,\cN_{K-1}$ with $Q_k\neq0$ on $\cN_k$ for $k\le K-1$; in particular
$c_D\neq0$ on $\cN_{K-1}$.  If $D=0$, then $Q=c_D$ does not involve $z_K$; set
$R_K:\equiv1$, and $Q=Q_{K-1}\neq0$ on $\cN_{K-1}$, hence on $\cN_K$.  If $D\ge1$, define
\[
R_K(z'):=1+\sum_{i<D}\frac{|c_i(z')|}{|c_D(z')|},
\]
continuous on $\cN_{K-1}$ (the $c_i$ are polynomials, hence continuous, and $c_D$ is
continuous and zero-free there).  For $z'\in\cN_{K-1}$ and $|z_K|>R_K(z')\ge1$,
\[
|c_D(z')z_K^D|=|c_D|\,|z_K|\cdot|z_K|^{D-1}
>|c_D|\,R_K(z')\,|z_K|^{D-1}
=\Big(|c_D|+\sum_{i<D}|c_i|\Big)|z_K|^{D-1}
\ge|c_D||z_K|^{D-1}+\sum_{i<D}|c_i|\,|z_K|^{i},
\]
using $|z_K|^i\le|z_K|^{D-1}$ for $i\le D-1$ ($|z_K|\ge1$); hence
$|c_Dz_K^D|>\big|\sum_{i<D}c_iz_K^i\big|$ and $Q(z',z_K)\neq0$.  Thus $Q_K=Q\neq0$ on
$\cN_K$, completing the induction.

\emph{Conclusion.}  Apply Step~2 to $Q:=P=p_1\cdots p_M$ (a polynomial with a dominant
top coefficient by Step~0).  For $1\le k\le K$ and $p_a\in\C[z_1,\dots,z_k]$, Step~1
gives that $p_a$ is a factor of $P_k$, which is nonvanishing on $\cN_k$; hence
$p_a\neq0$ on $\cN_k$.
\end{proof}

Definition~\ref{def:dtc} and Lemma~\ref{lem:nested} are consumed in Step~1 and Step~3 of
the proof of Proposition~\ref{prop:A1} (and nowhere else).

\subsection{Quadric rigidity}

Throughout this subsection, $A_1\in\R^{d\times d}$ is invertible and $d\ge2$.  The
\emph{probe quadric} is $\cQ:=\{(w,v)\in\R^{2d}:w^\top A_1v=0\}$, the zero set of
$q(w,v):=w^\top A_1v$.  For $w\in\R^d$ put
$H_w:=\{v\in\R^d:(A_1^\top w)^\top v=0\}=\ker\big((A_1^\top w)^\top\big)$, a linear
hyperplane when $A_1^\top w\neq0$; note $(w,v)\in\cQ\iff v\in H_w$.  A set
$U\subseteq\cQ$ (or $\subseteq\cQ\times(0,1)$) is \emph{relatively open} if it is open in
the subspace topology.

\begin{auxlemma}[Affine vanishing on a hyperplane slice]\label{lem:auxaffine}
Let $d\ge2$, $a\in\R^d\setminus\{0\}$, $H:=\{v\in\R^d:a^\top v=0\}$, and let
$f(v):=b^\top v+\beta$ with $b\in\R^d$, $\beta\in\R$.  If $f$ vanishes on a nonempty
relatively open subset $S\subseteq H$, then $f$ vanishes on all of $H$; consequently
\[
\beta=0\qquad\text{and}\qquad b\in H^\perp=\spn\{a\}.
\]
\end{auxlemma}

\begin{proof}
$H$ is a linear subspace with $\dim H=d-1\ge1$.  Choose a basis $u_1,\dots,u_{d-1}$ of
$H$ and define $\phi:\R^{d-1}\to H$, $\phi(t):=\sum_{i=1}^{d-1}t_iu_i$.  Then $\phi$ is a
linear bijection onto $H$, and it is a homeomorphism for the subspace topology on $H$
(every linear map between finite-dimensional normed spaces is continuous; this applies to
$\phi$ and to its inverse, the restriction to $H$ of any linear extension of
$\phi^{-1}$).  Hence $\phi^{-1}(S)\subseteq\R^{d-1}$ is open, and it is nonempty since
$S\neq\varnothing$.  The map $t\mapsto f(\phi(t))$ is a polynomial (affine) on
$\R^{d-1}$ vanishing on the nonempty open set $\phi^{-1}(S)$, hence identically zero by
Lemma~\ref{lem:zariski}; that is, $f\equiv0$ on $H$.

Since $0\in H$, $\beta=f(0)=0$.  Then $b^\top v=0$ for all $v\in H$, i.e.\
$b\in H^\perp$.  Finally, $\spn\{a\}\subseteq H^\perp$ by definition of $H$, and
$\dim H^\perp=d-\dim H=1=\dim\spn\{a\}$, so $H^\perp=\spn\{a\}$.
\end{proof}

\begin{lemma}[Slices and parallel-rigidity]\label{lem:slices}
Let $A_1\in\R^{d\times d}$ be invertible, $d\ge2$, and $U\subseteq\cQ$ a nonempty
relatively open set containing a point $(w_*,v_*)$ with $w_*\neq0$.  Then:
\begin{enumerate}
\item[\textup{(a)}] \textup{(Slice.)}  For every $(w,v)\in U$ with $w\neq0$, one has
$A_1^\top w\neq0$, and the fibre $U_w:=\{v'\in\R^d:(w,v')\in U\}$ is a nonempty
relatively open subset of the hyperplane $H_w$.
\item[\textup{(b)}] \textup{(Projection.)}  The set
$\Omega:=\{w\in\R^d:A_1^\top w\neq0,\ U_w\neq\varnothing\}$ is open and nonempty; in fact
it contains a neighbourhood of $w_*$.
\item[\textup{(c)}] \textup{(Parallel}$\Rightarrow$\textup{scalar.)}  If
$M,N\in\R^{d\times d}$ with $N$ invertible and $Mw\in\spn\{Nw\}$ for every $w$ in some
nonempty open $\Omega_0\subseteq\R^d$, then $M=cN$ for some scalar $c\in\R$.
\end{enumerate}
\end{lemma}

\begin{proof}
\emph{(a)}  Since $A_1$ is invertible, so is $A_1^\top$; hence $w\neq0$ implies
$A_1^\top w\neq0$, and $H_w$ is a genuine hyperplane.  For every $v'\in\R^d$,
\begin{equation}\label{eq:fibre}
(w,v')\in\cQ\iff w^\top A_1v'=(A_1^\top w)^\top v'=0\iff v'\in H_w .
\end{equation}
Since $U\subseteq\cQ$, \eqref{eq:fibre} gives $U_w\subseteq H_w$.  As $U$ is relatively
open, write $U=O\cap\cQ$ with $O\subseteq\R^{2d}$ open.  The map $\iota_w:H_w\to\R^{2d}$,
$\iota_w(v'):=(w,v')$, is continuous, and its image lies in $\cQ$ by \eqref{eq:fibre};
hence
\[
U_w=\{v'\in H_w:(w,v')\in U\}=\iota_w^{-1}(O),
\]
which is relatively open in $H_w$.  It contains $v$ (the second coordinate of the given
point $(w,v)\in U$), hence is nonempty.

\emph{(b)}  For $p\in\{1,\dots,d\}$ set $V_p:=\{w\in\R^d:(A_1^\top w)_p\neq0\}$, an open
subset of $\R^d$, and note
\begin{equation}\label{eq:cover}
\{w\in\R^d:A_1^\top w\neq0\}=\bigcup_{p=1}^d V_p .
\end{equation}
Fix $p$.  For $w\in V_p$ and $v\in\R^d$, writing
$v_{\widehat p}:=(v_i)_{i\neq p}\in\R^{d-1}$,
\begin{equation}\label{eq:solve}
(A_1^\top w)^\top v=0\iff
v_p=G_p(w,v_{\widehat p}):=-\frac{1}{(A_1^\top w)_p}\sum_{i\neq p}(A_1^\top w)_i v_i ,
\end{equation}
and $G_p$ is continuous on $V_p\times\R^{d-1}$.  Define
$\Psi_p:V_p\times\R^{d-1}\to\R^{2d}$ by $\Psi_p(w,v_{\widehat p}):=(w,v)$, where
$v_i:=(v_{\widehat p})_i$ for $i\neq p$ and $v_p:=G_p(w,v_{\widehat p})$.  By
\eqref{eq:solve}, $\Psi_p$ is a continuous bijection of $V_p\times\R^{d-1}$ onto
$\cQ_p:=\cQ\cap(V_p\times\R^d)$, with inverse $(w,v)\mapsto(w,v_{\widehat p})$, which is
continuous as the restriction of a linear projection.  Hence $\Psi_p$ is a homeomorphism
onto $\cQ_p$, and $\cQ_p$ is relatively open in $\cQ$ (intersection with the open set
$V_p\times\R^d$).

Since $U$ is relatively open in $\cQ$, the set $U\cap\cQ_p$ is relatively open in
$\cQ_p$, so
\[
W_p:=\Psi_p^{-1}(U)=\Psi_p^{-1}(U\cap\cQ_p)
\]
is open in $V_p\times\R^{d-1}$, hence open in $\R^d\times\R^{d-1}$.  The coordinate
projection $\pi(w,v_{\widehat p}):=w$ is an open map, so $\pi(W_p)$ is open in $\R^d$.
Moreover
\begin{equation}\label{eq:proj}
\pi(W_p)=\{w\in V_p:U_w\neq\varnothing\}.
\end{equation}
Indeed, if $(w,v_{\widehat p})\in W_p$ then $\Psi_p(w,v_{\widehat p})=(w,v)\in U$, so
$v\in U_w$ and $w\in V_p$; conversely, if $w\in V_p$ and $v\in U_w$, then
$(w,v)\in U\subseteq\cQ$ and $w\in V_p$, so $(w,v)\in U\cap\cQ_p$ and
$(w,v_{\widehat p})=\Psi_p^{-1}(w,v)\in W_p$.

Combining \eqref{eq:cover} and \eqref{eq:proj},
\[
\Omega=\{w:A_1^\top w\neq0,\ U_w\neq\varnothing\}
=\bigcup_{p=1}^d\{w\in V_p:U_w\neq\varnothing\}
=\bigcup_{p=1}^d\pi(W_p),
\]
a finite union of open sets, hence open.  Finally, $A_1^\top w_*\neq0$ by part (a) and
$v_*\in U_{w_*}$, so $w_*\in\Omega$; openness then yields a whole neighbourhood of $w_*$
inside $\Omega$.  In particular $\Omega\neq\varnothing$.

\emph{(c)}  The set $\Omega_0\setminus\{0\}$ is open, and nonempty: a nonempty open
subset of $\R^d$ contains a ball, hence more than one point, hence a point $w_0\neq0$.
Choose $r\in(0,|w_0|)$ small enough that the open ball $B:=B(w_0,r)$ satisfies
$B\subseteq\Omega_0$; then $B$ is nonempty, open, and $0\notin B$.

For $w\in B$ we have $w\neq0$, hence $Nw\neq0$ ($N$ invertible), and the hypothesis
$Mw\in\spn\{Nw\}$ provides a scalar $\mu(w)\in\R$ with $Mw=\mu(w)Nw$.  Setting
$P:=N^{-1}M$, this reads $Pw=\mu(w)w$; since $w\neq0$, $\mu(w)$ is a \emph{real}
eigenvalue of $P$ and $w\in\ker(P-\mu(w)I)$.  The real eigenvalues of $P$ are the real
roots of the degree-$d$ polynomial $\lambda\mapsto\det(P-\lambda I)$; list the distinct
ones as $\lambda_1,\dots,\lambda_k$ with $k\le d$.  (Non-real eigenvalues are irrelevant
here: each $\mu(w)$ is real by construction, and a nonzero real vector cannot be an
eigenvector for a non-real eigenvalue.)  Then
\[
B\subseteq\bigcup_{j=1}^k\ker(P-\lambda_jI),
\]
a finite union of linear subspaces of $\R^d$; in particular $k\ge1$, since
$B\neq\varnothing$.  The ball $B$ has nonempty interior, so by Lemma~\ref{lem:zariski}
the subspaces cannot all be proper: $\ker(P-\lambda_jI)=\R^d$ for some $j$, i.e.\
$P=\lambda_jI$.  Hence $M=NP=\lambda_jN$; take $c:=\lambda_j$.
\end{proof}

\begin{lemma}[Quadric rigidity, quadratic version]\label{lem:quadA}
Let $A_1$ be invertible, $d\ge2$, and $X,Y\in\R^{d\times d}$.  Suppose
$w^\top Xv+w^\top Yw=0$ for all $(w,v)$ in a relatively open subset $U\subseteq\cQ$
containing a point with $w\neq0$.  Then $X=cA_1$ for some $c\in\R$, and $\Sym(Y)=0$.
\end{lemma}

\begin{proof}
Let $(w_*,v_*)\in U$ with $w_*\neq0$ and apply Lemma~\ref{lem:slices}; let $\Omega$ be
the nonempty open set of part (b).

Fix $w\in\Omega$.  Then $A_1^\top w\neq0$ (so $w\neq0$) and $U_w\neq\varnothing$;
choosing any $v_0\in U_w$ and applying part (a) to the point $(w,v_0)\in U$, the fibre
$U_w$ is a nonempty relatively open subset of the hyperplane
$H_w=\{A_1^\top w\}^\perp$.  For every $v\in U_w$ we have $(w,v)\in U$, so the hypothesis
gives
\[
0=w^\top Xv+w^\top Yw=(X^\top w)^\top v+w^\top Yw\qquad(v\in U_w).
\]
Thus the affine function $f_w(v):=(X^\top w)^\top v+w^\top Yw$ vanishes on the nonempty
relatively open subset $U_w$ of $H_w$.  By Auxiliary Lemma~\ref{lem:auxaffine} (with
$a:=A_1^\top w$, $b:=X^\top w$, $\beta:=w^\top Yw$),
\begin{equation}\label{eq:quadAconc}
w^\top Yw=0\qquad\text{and}\qquad X^\top w\in\spn\{A_1^\top w\}
\qquad\text{for every }w\in\Omega .
\end{equation}

By Lemma~\ref{lem:slices}(c), applied with $M:=X^\top$, $N:=A_1^\top$ (invertible since
$A_1$ is) and $\Omega_0:=\Omega$, the second part of \eqref{eq:quadAconc} yields
$c\in\R$ with $X^\top=cA_1^\top$, i.e.\ $X=cA_1$.

For $Y$: the polynomial $w\mapsto w^\top Yw$ vanishes on the nonempty open set $\Omega$
by the first part of \eqref{eq:quadAconc}, hence vanishes identically on $\R^d$ by
Lemma~\ref{lem:zariski}.  Since $w^\top Yw=w^\top\Sym(Y)w$, the symmetric matrix
$S:=\Sym(Y)$ has identically vanishing quadratic form, and polarization
(\S\ref{subsec:notation}) gives $S=0$, i.e.\ $\Sym(Y)=0$.
\end{proof}

\begin{lemma}[Quadric rigidity, linear version]\label{lem:quadB}
Let $A_1$ be invertible, $d\ge2$, and $C,E\in\R^{d\times d}$.  Suppose $Cv+Ew=0$ for all
$(w,v)$ in a relatively open subset $U\subseteq\cQ$ containing a point with $w\neq0$.
Then $C=0$ and $E=0$.
\end{lemma}

\begin{proof}
Let $(w_*,v_*)\in U$ with $w_*\neq0$ and apply Lemma~\ref{lem:slices}; let $\Omega$ be
the nonempty open set of part (b).

Fix an index $j\in\{1,\dots,d\}$ and let $e_j$ be the $j$th standard basis vector.  For
$(w,v)\in U$ the hypothesis gives the scalar identity
\[
0=e_j^\top(Cv+Ew)=(C^\top e_j)^\top v+(E^\top e_j)^\top w .
\]
Fix $w\in\Omega$.  As in Lemma~\ref{lem:quadA}, parts (a)--(b) of
Lemma~\ref{lem:slices} show that $U_w$ is a nonempty relatively open subset of
$H_w=\{A_1^\top w\}^\perp$, and the affine function
$v\mapsto(C^\top e_j)^\top v+(E^\top e_j)^\top w$ vanishes on $U_w$.  By Auxiliary
Lemma~\ref{lem:auxaffine} (with $a:=A_1^\top w$, $b:=C^\top e_j$,
$\beta:=(E^\top e_j)^\top w$),
\begin{equation}\label{eq:quadBconc}
(E^\top e_j)^\top w=0\qquad\text{and}\qquad C^\top e_j\in\spn\{A_1^\top w\}
\qquad\text{for every }w\in\Omega .
\end{equation}

\emph{$E=0$.}  The (linear) polynomial $w\mapsto(E^\top e_j)^\top w$ vanishes on the
nonempty open set $\Omega$ by \eqref{eq:quadBconc}, hence identically on $\R^d$
(Lemma~\ref{lem:zariski}); therefore $E^\top e_j=0$.  As $j$ was arbitrary, $E=0$.

\emph{$C=0$.}  Set $u:=C^\top e_j$, a \emph{fixed} vector satisfying
$u\in\spn\{A_1^\top w\}$ for every $w\in\Omega$, by \eqref{eq:quadBconc}.  Since
$A_1^\top$ is invertible, it is a linear homeomorphism of $\R^d$, so $A_1^\top\Omega$ is
open and nonempty; moreover $0\notin A_1^\top\Omega$, since $A_1^\top w\neq0$ on
$\Omega$.  Pick any $a^{(1)}:=A_1^\top w^{(1)}\in A_1^\top\Omega$; then
$a^{(1)}\neq0$, and since $d\ge2$ the line $\spn\{a^{(1)}\}$ is a \emph{proper} subspace
of $\R^d$, hence has empty interior (Lemma~\ref{lem:zariski}).  The open set
$A_1^\top\Omega$ therefore cannot be contained in $\spn\{a^{(1)}\}$, so we may pick
$a^{(2)}:=A_1^\top w^{(2)}\in A_1^\top\Omega\setminus\spn\{a^{(1)}\}$; then
$a^{(1)},a^{(2)}$ are linearly independent.  Now
$u\in\spn\{a^{(1)}\}\cap\spn\{a^{(2)}\}$.  If $u\neq0$, write
$u=\alpha a^{(1)}=\beta a^{(2)}$ with $\alpha,\beta\neq0$; then
$a^{(2)}=(\alpha/\beta)\,a^{(1)}\in\spn\{a^{(1)}\}$, a contradiction.  Hence
$u=C^\top e_j=0$.  As $j$ was arbitrary, $C=0$.
\end{proof}

\begin{remark}\label{rem:rigidity}
The hypothesis $d\ge2$ is essential in Lemma~\ref{lem:quadB} (it is used exactly where
the line $\spn\{a^{(1)}\}$ is proper): for $d=1$, $A_1=(1)$, the set
$U:=\{(w,0):w\in(1,2)\}=\big((1,2)\times(-1,1)\big)\cap\cQ$ is relatively open in $\cQ$
and contains points with $w\neq0$, and $Cv+Ew=0$ holds on $U$ for $E=0$ and \emph{every}
$C$.  Invertibility of $A_1$ enters in Lemma~\ref{lem:slices}(a) ($A_1^\top w\neq0$), in
Lemma~\ref{lem:slices}(c) (as $N=A_1^\top$), and in the last step of
Lemma~\ref{lem:quadB} (openness of $A_1^\top\Omega$).
\end{remark}

\section{Step 1: identification of the attention matrix}\label{sec:A1}

\subsection*{Standing hypothesis and conventions for Sections 5--7}
Throughout Sections 5--7, $L\ge2$, $d\ge2$, $r\ge2$ are fixed,
$\theta'=(V_1',A_1',\dots,V_L',A_L')\in\cG^{(L)}$ (Definition~\ref{def:generic}),
$\theta=(V_1,A_1,\dots,V_L,A_L)\in\cP$ is arbitrary, and $\cO\subseteq\R^{2d}$ is open and
nonempty.  The \emph{standing hypothesis} is exactly the hypothesis of
Theorem~\ref{thm:IDL} at depth $L$:
\begin{equation}\tag{SH}\label{eq:SH}
\cO\cap\cM_L(\theta')\neq\varnothing,\qquad\text{and}\qquad
\forall\,P\in\cP(\cO)\ \ \exists\,T_P:\quad G^P_\theta=G^P_{\theta'}
\ \text{ on }(T_P,\infty).
\end{equation}
The present section uses only the second clause of \eqref{eq:SH}, and only through
constant probes: constant paths $P\equiv(w,v)$, $(w,v)\in\cO$, lie in $\cP(\cO)$
(Definition~\ref{def:paths}) and their observables are the constant-probe observables, so
\eqref{eq:SH} implies
\begin{equation}\tag{H}\label{eq:H}
\forall\,(w,v)\in\cO\ \ \exists\,T_{(w,v)}:\quad
F^{(L)}_\theta(w,v,\tau)=F^{(L)}_{\theta'}(w,v,\tau)
\quad\text{for all real }\tau>T_{(w,v)} .
\end{equation}
The anchor clause of \eqref{eq:SH} is consumed only in Section~\ref{sec:cascade}
(Lemma~\ref{lem:region}).

\emph{Stratification conventions.}  Given a probe $(w,v)\in\R^{2d}$, primed objects
$\lambda'_j,\phi'_j,s'_j,H'_j,S'^{\,j},\dots$ are those attached to $\theta'$ at $(w,v)$
by Proposition~\ref{prop:strat} and Lemma~\ref{lem:phi}, and unprimed ones the same for
$\theta$.  Write $\Omega'_0:=\C$ and
$\Omega'_j:=\C\setminus(S'^{\,1}\cup\dots\cup S'^{\,j})$, and
$S_{\theta'}:=S'^{\,1}\cup\dots\cup S'^{\,L}$, so $\Omega'_L=\C\setminus S_{\theta'}$;
analogously $S^j$, $\Omega_j$, $S_\theta$ for $\theta$.  By Lemma~\ref{lem:strataacc}
each partial union is closed, so each $\Omega'_j$ is open (indeed a domain, by the proof
of Proposition~\ref{prop:strat}); by Proposition~\ref{prop:strat}(b),(c), $s'_i$ is
holomorphic on $\Omega'_i$ and $S_{\theta'}$ is countable and closed.  Since
$\Omega'_{j-1}\subseteq\Omega'_i$ for $i\le j-1$, the truncated gate vector
$s'_{<j}:=(s'_1,\dots,s'_{j-1}):\Omega'_{j-1}\to\C^{j-1}$ is holomorphic and is in
particular defined at every point of $S'^{\,j}\subseteq\Omega'_{j-1}$.  Finally
$F_{\theta'}$ denotes the (unique) holomorphic $\C^d$-valued extension of
$\tau\mapsto F^{(L)}_{\theta'}(w,v,\tau)$ from $(0,\infty)$ to $\Omega'_L$ given by
Proposition~\ref{prop:strat}(c), and likewise $F_\theta$ on $\C\setminus S_\theta$.

\begin{auxlemma}[Blow-up at stratum points]\label{lem:stratpole}
Let $\vartheta$ be an $m$-tuple, $(w,v)$ a probe, and adopt the notation of
Proposition~\ref{prop:strat}.  Let $1\le j\le m$ and $\xi\in S^j$.  Then:
\begin{enumerate}
\item[(i)] there is $\varepsilon>0$ with
$D(\xi,\varepsilon)\setminus\{\xi\}\subseteq\Omega_j$;
\item[(ii)] $|s_j(\tau)|\to\infty$ as $\tau\to\xi$ with $\tau\in\Omega_j$.
\end{enumerate}
\end{auxlemma}

\begin{proof}
(i)  For $j=1$: $S^1$ is closed and discrete in $\C$ (Proposition~\ref{prop:strat},
hypothesis system (I1)), so some disc $D(\xi,\varepsilon)$ meets $S^1$ only in $\xi$,
whence $D(\xi,\varepsilon)\setminus\{\xi\}\subseteq\C\setminus S^1=\Omega_1$.  For
$j\ge2$: $\xi\in S^j\subseteq\Omega_{j-1}$, the partial union
$T_{j-1}=S^1\cup\dots\cup S^{j-1}$ is closed (Lemma~\ref{lem:strataacc}) and misses
$\xi$, and $S^j$ is closed and discrete in the open set $\Omega_{j-1}$; choose
$\varepsilon>0$ so small that $D(\xi,\varepsilon)\subseteq\Omega_{j-1}$ and
$D(\xi,\varepsilon)\cap S^j=\{\xi\}$.  Then
$D(\xi,\varepsilon)\setminus\{\xi\}\subseteq\Omega_{j-1}\setminus S^j=\Omega_j$.

(ii)  For $j=1$: by Proposition~\ref{prop:strat}(b), $s_1$ is meromorphic on $\C$ with
pole set exactly $S^1$, and at a pole $|s_1(\tau)|\to\infty$.  For $j\ge2$:
$s_j=\sig\circ H_j$ with $H_j$ holomorphic at $\xi$ and $H_j(\xi)\in\Pi$
(definition of $S^j$).  Writing $\sig=1/(1+e^{-z})$, the denominator $1+e^{-z}$ is
continuous and vanishes exactly on $\Pi$ (Lemma~\ref{lem:sigmoid}); hence as
$\tau\to\xi$ with $\tau\in\Omega_j$, $H_j(\tau)\to H_j(\xi)\in\Pi$ with
$H_j(\tau)\notin\Pi$, so $1+e^{-H_j(\tau)}\to0$ with nonzero values, and
$|s_j(\tau)|=|1+e^{-H_j(\tau)}|^{-1}\to\infty$.
\end{proof}

\begin{auxlemma}[Last-layer formula]\label{lem:lastlayer}
Let $\vartheta$ be an $m$-tuple, $(w,v)$ a probe, and let $F_\vartheta$ be the
holomorphic extension of Proposition~\ref{prop:strat}(c).  Then, on $\Omega_m$,
\begin{equation}\label{eq:lastlayer}
F_\vartheta(\tau)
= B_{m:1}v+\sum_{i=1}^{m}s_i(\tau)\,B_{m:i+1}V_i\,
W_{i-1}\big(s_1(\tau),\dots,s_{i-1}(\tau)\big)\,w ,
\end{equation}
where $W_{i-1}$ is the matrix polynomial of Lemma~\ref{lem:phi}.
\end{auxlemma}

\begin{proof}
By Lemma~\ref{lem:phi}(i), $v_m(z)=B_{m:1}v+T_m(z)w$ with
$T_m(z)=\sum_{i=1}^m z_iB_{m:i+1}V_iW_{i-1}(z)$.  Define $\widehat G(\tau)$ to be the
right-hand side of \eqref{eq:lastlayer}: each $s_i$ is holomorphic on
$\Omega_i\supseteq\Omega_m$, so $\widehat G$ is holomorphic on $\Omega_m$.  For real
$\tau>0$, the constructed functions $s_i(\tau)$ equal the recursion gates
(Proposition~\ref{prop:strat}(b)), so by Lemma~\ref{lem:lambdaphi} and
Proposition~\ref{prop:probe},
\[
\widehat G(\tau)=B_{m:1}v+T_m\big(s_1,\dots,s_m\big)w
=v_m(z)\big|_{z=(s_1,\dots,s_m)}=F^{(m)}_\vartheta(w,v,\tau).
\]
Thus $\widehat G$ is a holomorphic extension of $F^{(m)}_\vartheta(w,v,\cdot)$ from
$(0,\infty)$ to $\Omega_m$; by the uniqueness clause of
Proposition~\ref{prop:strat}(c), $\widehat G=F_\vartheta$ on $\Omega_m$.
\end{proof}

\begin{proposition}[Identification of the attention matrix]\label{prop:A1}
Assume \eqref{eq:H}.  Then $A_1=A_1'$.
\end{proposition}

\begin{proof}
\emph{Step 0: the probe, the coordinate, and the agreement.}
Define $\kappa_j(w):=w^\top\Sym\!\big(V_{j-1:1}'^{\,\top}A_j'V_{j-1:1}'\big)w$ for
$2\le j\le L$, and
\[
\cO_\star:=\big\{(w,v)\in\cO:\ w^\top A_1'v\neq0,\ \ \kappa_j(w)\neq0\ (2\le j\le L),
\ \ V'_{L:1}w\neq0\big\}.
\]
Each defining condition is the nonvanishing of a polynomial, and each polynomial is
nonzero: $(w,v)\mapsto w^\top A_1'v$ because $A_1'\neq0$ (by (G1), $\det A_1'\neq0$;
test on $(e_a,e_b)$); $w\mapsto\kappa_j(w)$ because
$\Sym(V_{j-1:1}'^{\,\top}A_j'V_{j-1:1}')\neq0$ by (G2) (a nonzero symmetric matrix has a
nonvanishing quadratic form, by polarization, \S\ref{subsec:notation}); and
$w\mapsto V'_{L:1}w$ componentwise because $V'_{L:1}\neq0$ by (G2).  By
Lemma~\ref{lem:zariski}, the finitely many polynomials have a common nonvanishing point
in the nonempty open set $\cO$; in fact $\cO_\star$ is open and dense in $\cO$.

Fix $(w,v)\in\cO_\star$ for Steps 1--4.  Write $\lambda_1':=w^\top A_1'v\neq0$, and
choose a coordinate $\iota\in\{1,\dots,d\}$ with $e_\iota^\top V'_{L:1}w\neq0$ (possible
since $V'_{L:1}w\neq0$).  Set $T_0:=\max\{T_{(w,v)},1\}$, so that by \eqref{eq:H}
\begin{equation}\label{eq:agree}
F^{(L)}_\theta(w,v,\tau)=F^{(L)}_{\theta'}(w,v,\tau)
\qquad\big(\tau\in(T_0,\infty)\subseteq(0,\infty)\big).
\end{equation}

\emph{Step 1: the polynomial family and the largeness regions.}
By Lemma~\ref{lem:phi}(iii), for $2\le m\le L$ the coefficient of $z_{m-1}^2$ in
$\phi'_m(z_1,\dots,z_{m-1};w,v)$ is the polynomial
$f_m(z_1,\dots,z_{m-2})\in\C[z_1,\dots,z_{m-2}]$ whose coefficient of
$z_1^2\cdots z_{m-2}^2$ equals $-\kappa_m(w)\neq0$; moreover
$\deg_{z_i}f_m\le2$ for each $i$ (the coefficient of $z_{m-1}^2$ in a polynomial of
per-variable degree $\le2$, Lemma~\ref{lem:phi}(ii)).  Hence $f_m$ has per-variable
degrees exactly $2$ in $z_1,\dots,z_{m-2}$ and a dominant top coefficient in the sense
of Definition~\ref{def:dtc} (for $m=2$, $f_2\equiv-\kappa_2(w)$ is a nonzero constant).
By Lemma~\ref{lem:phi}(v), the scalar polynomial
\[
g(z_1,\dots,z_{L-1}):=e_\iota^\top V_L'\,W'_{L-1}(z)\,w\in\C[z_1,\dots,z_{L-1}]
\]
is multilinear with coefficient of $z_1\cdots z_{L-1}$ equal to
$(-1)^{L-1}e_\iota^\top V'_{L:1}w\neq0$: again a dominant top coefficient.

Apply Lemma~\ref{lem:nested} to the family
$\cF:=\{f_2,\dots,f_L,\,g\}\subset\C[z_1,\dots,z_{L-1}]$ (so $K=L-1\ge1$), and let
$\cN_1,\dots,\cN_{L-1}$ and $R_1$, $R_k(\cdot)$ be the resulting regions and continuous
thresholds.  Recording the conclusion:
\begin{equation}\tag{R}\label{eq:R}
\text{for }1\le k\le L-1:\quad
\text{every member of }\cF\text{ lying in }\C[z_1,\dots,z_k]
\text{ is nonvanishing on }\cN_k .
\end{equation}
In particular $f_{k+2}\ (\in\C[z_1,\dots,z_k])$ and, for $k=L-1$, $g$ are nonvanishing
on the respective regions; the constant $f_2$ is nonvanishing everywhere.  Each $\cN_k$
is open in $\C^k$: by induction, $\cN_1$ is open, and
$\cN_k=\{(z',z_k):z'\in\cN_{k-1},\,|z_k|>R_k(z')\}$ is the preimage of
$(0,\infty)$ under the continuous map $(z',z_k)\mapsto|z_k|-R_k(z')$ on the open set
$\cN_{k-1}\times\C$.

\emph{Step 2: the tier sets.}  Define
\[
\cT_1:=S'^{\,1},\qquad
\cT_j:=\big\{\xi\in S'^{\,j}:\ s'_{<j}(\xi)\in\cN_{j-1}\big\}\quad(2\le j\le L).
\]
Since $\lambda_1'\neq0$, Proposition~\ref{prop:strat}(d) gives
\[
\cT_1=S'^{\,1}=\Bigl\{\tfrac{-b+(2n+1)\pi i}{\lambda_1'}:n\in\Z\Bigr\}\neq\varnothing,
\]
an infinite set; every $\xi\in\cT_1$ has
$\operatorname{Re}\xi=-b/\lambda_1'$ and
$\operatorname{Im}\xi=(2n+1)\pi/\lambda_1'\neq0$, so $\xi\notin\R$ and $\xi\neq0$.
More generally $S'^{\,j}\cap\R=\varnothing$ for every $j$
(Proposition~\ref{prop:strat}(a)), so every tier point is nonreal, in particular
nonzero.

\emph{Step 3: ignition and propagation.}

\emph{Claim B (propagation).}  For $1\le j\le L-1$ and every $\xi\in\cT_j$: every
punctured disc around $\xi$ contains infinitely many points of $\cT_{j+1}$; in
particular $\xi\in\acc(\cT_{j+1})$.

\emph{Proof of Claim B.}  Fix $\xi\in\cT_j$.  By Auxiliary
Lemma~\ref{lem:stratpole} there is $\varepsilon_0>0$ with
$D^*:=D(\xi,\varepsilon_0)\setminus\{\xi\}\subseteq\Omega'_j$ and
$|s'_j(\tau)|\to\infty$ as $\tau\to\xi$ in $D^*$.  By Lemma~\ref{lem:phi}(ii),(iii),
collect $\phi'_{j+1}$ by powers of $z_j$:
\[
\phi'_{j+1}(z_1,\dots,z_j;w,v)
=f_{j+1}(z_{<j})\,z_j^2+c_1(z_{<j})\,z_j+c_0(z_{<j}),
\qquad c_0,c_1\in\C[z_1,\dots,z_{j-1}],
\]
so that on $\Omega'_j$,
\[
H'_{j+1}(\tau)=\tau\Big[f_{j+1}\big(s'_{<j}(\tau)\big)\,s'_j(\tau)^2
+c_1\big(s'_{<j}(\tau)\big)\,s'_j(\tau)+c_0\big(s'_{<j}(\tau)\big)\Big]+b .
\]
As $\tau\to\xi$ in $D^*$: (i) $\tau\to\xi\neq0$; (ii) for $j\ge2$,
$s'_{<j}(\tau)\to\sigma^0:=s'_{<j}(\xi)\in\cN_{j-1}$ by continuity of the holomorphic
$s'_{<j}$ at $\xi\in\Omega'_{j-1}$, hence
$f_{j+1}(s'_{<j}(\tau))\to f_{j+1}(\sigma^0)\neq0$ by \eqref{eq:R}, and
$c_0,c_1$ evaluated along $s'_{<j}(\tau)$ stay bounded; for $j=1$,
$f_2\equiv-\kappa_2(w)\neq0$ and $c_0,c_1$ are constants; (iii)
$|s'_j(\tau)|\to\infty$.  Hence
\[
\big|\phi'_{j+1}\big(s'_{\le j}(\tau)\big)\big|
\ \ge\ \big|f_{j+1}(s'_{<j}(\tau))\big|\,|s'_j(\tau)|^2
-\big|c_1(s'_{<j}(\tau))\big|\,|s'_j(\tau)|-\big|c_0(s'_{<j}(\tau))\big|
\ \longrightarrow\ \infty,
\]
and $|H'_{j+1}(\tau)|\ge|\tau|\,|\phi'_{j+1}|-b\to\infty$ as $\tau\to\xi$ in $D^*$.

$H'_{j+1}$ is holomorphic on $\Omega'_j\supseteq D^*$, so Lemma~\ref{lem:T4} applies to
$H:=H'_{j+1}$ on the punctured disc $D^*$: for every $\varepsilon\in(0,\varepsilon_0)$
the set
\[
Z_\varepsilon:=\{\tau:\ 0<|\tau-\xi|<\varepsilon,\ H'_{j+1}(\tau)\in\Pi\}
\subseteq S'^{\,j+1}
\]
is infinite (the inclusion holds because $Z_\varepsilon\subseteq D^*\subseteq\Omega'_j$
and $S'^{\,j+1}=H_{j+1}'^{-1}(\Pi)\subseteq\Omega'_j$).  It remains to check that
$Z_\varepsilon\subseteq\cT_{j+1}$ for all sufficiently small $\varepsilon$, i.e.\ that
$s'_{<j+1}(\tau)=\big(s'_{<j}(\tau),s'_j(\tau)\big)\in\cN_j$ for $\tau$ near $\xi$.
For $j\ge2$: $s'_{<j}(\tau)$ lies in the open set $\cN_{j-1}$ for $\tau$ close to
$\xi$ (continuity, $\sigma^0\in\cN_{j-1}$); the threshold satisfies
$R_j\big(s'_{<j}(\tau)\big)\le R_j(\sigma^0)+1$ for $\tau$ close to $\xi$ (continuity of
$R_j$ at $\sigma^0$); and $|s'_j(\tau)|\to\infty$ exceeds this bound for $\tau$ close to
$\xi$.  For $j=1$: $\cN_1=\{|z_1|>R_1\}$ and $|s'_1(\tau)|\to\infty>R_1$ near $\xi$.
Hence for small $\varepsilon$, $Z_\varepsilon\subseteq\cT_{j+1}$, and every punctured
disc around $\xi$ (containing some small $Z_\varepsilon$) meets $\cT_{j+1}$ in
infinitely many points.  \hfill$\diamond$

\emph{Claim C (blow-up of the observable at tier $L$).}  For every $\xi\in\cT_L$:
$|e_\iota^\top F_{\theta'}(\tau)|\to\infty$ as $\tau\to\xi$ with $\tau\in\Omega'_L$.

\emph{Proof of Claim C.}  Let $\xi\in\cT_L$, so $\sigma^0:=s'_{<L}(\xi)\in\cN_{L-1}$.
By Auxiliary Lemma~\ref{lem:stratpole}, fix a punctured disc
$D^*\subseteq\Omega'_L$ around $\xi$ on which all $s'_i$ are holomorphic ($i<L$) and
$|s'_L(\tau)|\to\infty$.  By Auxiliary Lemma~\ref{lem:lastlayer} (with
$\vartheta=\theta'$, $m=L$; note $B'_{L:L+1}=I$),
\[
e_\iota^\top F_{\theta'}(\tau)
=e_\iota^\top B'_{L:1}v
+\sum_{i=1}^{L-1}s'_i(\tau)\,e_\iota^\top B'_{L:i+1}V'_i\,W'_{i-1}\big(s'_{<i}(\tau)\big)w
\;+\;s'_L(\tau)\,g\big(s'_{<L}(\tau)\big),
\]
with $g$ as in Step 1.  The terms with $i\le L-1$ are holomorphic at
$\xi\in\Omega'_{L-1}$, hence bounded on a neighbourhood; and
$g(s'_{<L}(\tau))\to g(\sigma^0)\neq0$ by \eqref{eq:R}.  Hence
\[
\big|e_\iota^\top F_{\theta'}(\tau)\big|
\ \ge\ |s'_L(\tau)|\,\big|g\big(s'_{<L}(\tau)\big)\big|-C
\ \longrightarrow\ \infty .  \hfill\diamond
\]

\emph{Claim D (transfer to the unknown side).}  $\cT_L\subseteq S_\theta$.

\emph{Proof of Claim D.}  Both $S_\theta$ and $S_{\theta'}$ are countable and closed
(Proposition~\ref{prop:strat}(c)), and both avoid $\R$
(Proposition~\ref{prop:strat}(a)); the scalar functions $F:=e_\iota^\top F_\theta$ and
$G:=e_\iota^\top F_{\theta'}$ are holomorphic on $\C\setminus S_\theta$ and
$\C\setminus S_{\theta'}$ respectively and agree on $(T_0,\infty)$ by \eqref{eq:agree},
an interval disjoint from $S_\theta\cup S_{\theta'}$.  Each $\xi\in\cT_L$ is an isolated
point of $S_{\theta'}$ (Auxiliary Lemma~\ref{lem:stratpole}(i): a punctured disc around
$\xi$ misses $S_{\theta'}$), and $|G(\tau)|\to\infty$ as $\tau\to\xi$ off $S_{\theta'}$
by Claim C.  Lemma~\ref{lem:transfer} yields $\xi\in S_\theta$.  \hfill$\diamond$

\emph{Step 4: descent and comparison of the first strata.}  By Claim B and the
monotonicity of iterated accumulation (Auxiliary Lemma~\ref{lem:discrete}(iii)), an
induction on $j$ gives $\cT_1\subseteq\acc^{(j-1)}(\cT_j)$ for $1\le j\le L$: the base
is trivial, and $\cT_1\subseteq\acc^{(j-1)}(\cT_j)$ together with
$\cT_j\subseteq\acc(\cT_{j+1})$ (Claim B, applied pointwise) gives
$\cT_1\subseteq\acc^{(j-1)}\big(\acc(\cT_{j+1})\big)=\acc^{(j)}(\cT_{j+1})$.  Hence,
using Claim D and monotonicity once more, and then Lemma~\ref{lem:strataacc} applied to
the unprimed stratification $S_\theta=S^1\cup\dots\cup S^L$ (whose structure is provided
by Proposition~\ref{prop:strat}) with $k=L-1$:
\[
S'^{\,1}=\cT_1\subseteq\acc^{(L-1)}(\cT_L)\subseteq\acc^{(L-1)}(S_\theta)
\subseteq S^1 .
\]
In particular $S^1\supseteq S'^{\,1}\neq\varnothing$; by
Proposition~\ref{prop:strat}(d) for $\theta$, this forces
$\lambda_1:=w^\top A_1v\neq0$ (else $S^1=\varnothing$) and every point of $S^1$ has real
part $-b/\lambda_1$.  Take the single point
$\zeta:=(-b+\pi i)/\lambda_1'\in S'^{\,1}\subseteq S^1$: comparing real parts,
\[
-\,\frac{b}{\lambda_1'}=\operatorname{Re}\zeta=-\,\frac{b}{\lambda_1},
\]
and since $b=\log r>0$, we conclude $\lambda_1=\lambda_1'$, i.e.
\begin{equation}\label{eq:bilinear}
w^\top(A_1-A_1')\,v=0 .
\end{equation}

\emph{Step 5: from one probe to the matrix.}  Steps 0--4 establish
\eqref{eq:bilinear} at every $(w,v)\in\cO_\star$, an open set, nonempty by Step 0.  The
polynomial $(w,v)\mapsto w^\top(A_1-A_1')v$ on $\R^{2d}$ vanishes on $\cO_\star$, hence
identically (Lemma~\ref{lem:zariski}); evaluating at $(w,v)=(e_a,e_b)$ for all $a,b$
gives $(A_1-A_1')_{ab}=0$, i.e.\ $A_1=A_1'$.
\end{proof}

\begin{corollary}[Depth certificate]\label{cor:depth}
Assume \eqref{eq:H}.  Then $V_1\neq0$.
\end{corollary}

\begin{proof}
Suppose $V_1=0$.  Fix the probe $(w,v)\in\cO_\star$ and the coordinate $\iota$ of
Step~0 of the proof of Proposition~\ref{prop:A1}, and keep the tier sets
$\cT_1,\dots,\cT_L$ of Steps 1--3, which depend only on the primed data.  By
Proposition~\ref{prop:strat}(e), $F^{(L)}_\theta(w,v,\cdot)
=F^{(L-1)}_{\theta_{\ge2}}(w,v,\cdot)$ on $(0,\infty)$, and the $(L-1)$-tuple
$\theta_{\ge2}$ carries its own stratification
$\widetilde S=\widetilde S^1\cup\dots\cup\widetilde S^{L-1}$ at $(w,v)$ with all the
properties of Proposition~\ref{prop:strat}; in particular $\widetilde S$ is countable
and closed, avoids $\R$, and $F^{(L)}_\theta(w,v,\cdot)$ extends holomorphically to
$\C\setminus\widetilde S$.  Claims C and D of the proof of
Proposition~\ref{prop:A1} go through verbatim with $E_F:=\widetilde S$ in place of
$S_\theta$ (Lemma~\ref{lem:transfer} is applied to the extension of
$e_\iota^\top F^{(L)}_\theta(w,v,\cdot)$ off $\widetilde S$, which agrees with
$e_\iota^\top F_{\theta'}$ on $(T_0,\infty)$ by \eqref{eq:agree}), giving
$\cT_L\subseteq\widetilde S$.  As in Step~4,
\[
\varnothing\neq S'^{\,1}=\cT_1\subseteq\acc^{(L-1)}(\cT_L)
\subseteq\acc^{(L-1)}(\widetilde S).
\]
But $\widetilde S$ has only $m=L-1$ strata, so Lemma~\ref{lem:strataacc} gives
$\acc^{(L-1)}(\widetilde S)=\acc^{(m)}(\widetilde S)=\varnothing$ --- a contradiction.
Hence $V_1\neq0$.
\end{proof}

From now on the common matrix $A_1=A_1'$ is denoted $A_1$; by (G1),
$\det A_1\neq0$ and $\Sym(A_1)\neq0$.  Consequently the quadric
$\cQ=\{q=0\}$ with $q(w,v)=w^\top A_1v$, and the covector
$\pi(w,v)=A_1^\top w-A_1v$, are common to both parameters, as are all objects of
\S\ref{sec:toolkit} built from $A_1$.

\section{Step 2: saturated limits, the trichotomy, and matching}\label{sec:cascade}

Throughout this section the standing hypothesis \eqref{eq:SH} of \S\ref{sec:A1} remains
in force.  In addition, by Proposition~\ref{prop:A1} and Corollary~\ref{cor:depth},
$A_1=A_1'$ (so $q$, $\pi$, $\cQ$, and the dial paths below are common to both
parameters) and $V_1\neq0$; by (G1), $\det A_1\neq0$ and $\Sym(A_1)\neq0$.  The output
of this section is Proposition~\ref{prop:matching}:
$B_{L:1}=B'_{L:1}$, $B_1=B_1'$ and $V_1=V_1'$.

\subsection{The sign region and the dial paths}

\begin{lemma}[Sign region]\label{lem:region}
Assume \eqref{eq:SH}.  There is a nonempty, relatively open, connected set
$\cU\subseteq\cQ\times(0,1)$ such that every $(w,v,t)\in\cU$ satisfies
\begin{enumerate}
\item[(i)] $w\neq0$, $A_1^\top w\neq0$ and $\pi(w,v)\neq0$;
\item[(ii)] $(w,v)\in\cO$;
\item[(iii)] $\phi'_\ell(t,1,\dots,1;w,v)>0$ for $\ell=2,\dots,L$.
\end{enumerate}
Moreover, every point of every nonempty relatively open subset
$\cU_0\subseteq\cU$ of $\cQ\times(0,1)$ admits a \emph{product neighbourhood}
$N=U\times J\subseteq\cU_0$, where $J\subseteq(0,1)$ is a nondegenerate open interval,
$U$ is a relatively open subset of $\cQ$ containing a point with $w\neq0$, and the
slice $U$ is the same for every $t\in J$.
\end{lemma}

\begin{proof}
\emph{Step 1 (the anchor).}  The proof uses only the anchor clause of \eqref{eq:SH}:
fix $(w^\sharp,v^\sharp)\in\cO\cap\cM_L(\theta')$.  (At the top level of the induction
this clause is supplied by Proposition~\ref{prop:anchorsexist} via the hypotheses of
Theorem~\ref{thm:main}; inside the induction of \S\ref{sec:descent} it is re-established
at each depth --- see step (3) of the proof of Theorem~\ref{thm:IDL}.)  Recall that
$\cM_L(\theta')$ is defined through $\theta'$'s first layer, whose attention matrix is
the common $A_1$.  By (M1): $q(w^\sharp,v^\sharp)=0$, $A_1^\top w^\sharp\neq0$ (hence
$w^\sharp\neq0$) and $\pi(w^\sharp,v^\sharp)\neq0$.  By (M2): there is
$s^\sharp\in(0,1)$ with $\phi'_\ell(s^\sharp,1,\dots,1;w^\sharp,v^\sharp)>0$ for
$\ell=2,\dots,L$.  Thus $p^\sharp:=(w^\sharp,v^\sharp,s^\sharp)\in\cQ\times(0,1)$
satisfies (i)--(iii).

\emph{Step 2 (the open conditions).}  Each function
$(w,v,t)\mapsto\phi'_\ell(t,1,\dots,1;w,v)$ is a polynomial in $(t,w,v)$ jointly: by
Lemma~\ref{lem:phi}(i) it equals $w^\top X'_\ell(t,1,\dots,1)v+w^\top
Y'_\ell(t,1,\dots,1)w$ with matrix entries polynomial in $t$.  Define
\[
\cV:=\big\{(w,v,t)\in\R^{2d}\times(0,1):\ w\neq0,\ A_1^\top w\neq0,\ \pi(w,v)\neq0,\
(w,v)\in\cO,\ \phi'_\ell(t,1,\dots,1;w,v)>0\ (2\le\ell\le L)\big\},
\]
an open subset of $\R^{2d+1}$ (finitely many conditions, each the nonvanishing or
positivity of a continuous function, or membership in the open sets $\cO$ and $(0,1)$).
By Step 1, $p^\sharp\in\cV\cap(\cQ\times(0,1))$.

\emph{Step 3 (graph charts).}  The gradient of $q$ is
$\nabla q(w,v)=(A_1v,\,A_1^\top w)$, which, since $A_1$ is invertible, vanishes exactly
at $(w,v)=(0,0)$.  At any point of $\cQ\setminus\{(0,0)\}$ some partial derivative of
$q$ is nonzero, so by the implicit function theorem (for the polynomial $q$) the
quadric is, near that point, the graph of a real-analytic function of the remaining
$2d-1$ coordinates: $\cQ\setminus\{(0,0)\}$ is a $(2d-1)$-dimensional real-analytic
submanifold of $\R^{2d}$, covered by graph charts (homeomorphisms with open subsets of
$\R^{2d-1}$).

\emph{Step 4 (local connectedness).}  Open subsets of $\R^{2d-1}$ are locally
connected; transporting along graph charts, $\cQ\setminus\{(0,0)\}$ is locally
connected, and so is
\[
S:=\big(\cQ\setminus\{(0,0)\}\big)\times(0,1)
\]
(around any point, the product of a connected chart neighbourhood with a subinterval is
a connected neighbourhood, and such products form a neighbourhood base).  Moreover $S$
is relatively open in $\cQ\times(0,1)$: its complement there is
$\{(0,0)\}\times(0,1)$, which is relatively closed.

\emph{Step 5 (definition of $\cU$).}  Let $\cW:=\cV\cap(\cQ\times(0,1))$, relatively
open in $\cQ\times(0,1)$ and contained in $S$ (its points have $w\neq0$); hence $\cW$
is open in the subspace $S$.  Define $\cU$ to be the connected component of $p^\sharp$
in $\cW$.  Since $S$ is locally connected, components of open subsets of $S$ are open
in $S$; as $S$ is relatively open in $\cQ\times(0,1)$, $\cU$ is relatively open there.
By construction $\cU$ is connected, nonempty, and every point satisfies (i)--(iii).

\emph{Step 6 (product neighbourhoods).}  Let $\cU_0\subseteq\cU$ be nonempty and
relatively open, say $\cU_0=O\cap(\cQ\times(0,1))$ with $O\subseteq\R^{2d+1}$ open, and
let $(w_1,v_1,t_1)\in\cU_0$.  Choose $\delta>0$ and $\varepsilon>0$ with
$B\big((w_1,v_1),\delta\big)\times(t_1-\varepsilon,t_1+\varepsilon)\subseteq O$ and
$(t_1-\varepsilon,t_1+\varepsilon)\subseteq(0,1)$, and set
$U:=B((w_1,v_1),\delta)\cap\cQ$, $J:=(t_1-\varepsilon,t_1+\varepsilon)$.  Then
$U\times J\subseteq O\cap(\cQ\times(0,1))=\cU_0$, $U$ is relatively open in $\cQ$,
contains $(w_1,v_1)$ (a point with $w_1\neq0$, by (i)), and is independent of $t\in J$.
\end{proof}

\begin{lemma}[Dial paths]\label{lem:dial}
Let $(w^0,v^0,t)\in\cU$.  Set $\pi^0:=\pi(w^0,v^0)\neq0$, $c:=\logit t-b$,
$y:=\pi^0/|\pi^0|^2$, and define the \emph{dial path}
$P_{w^0,v^0,t}=(w,v)$,
\[
w(\tau):=w^0-\frac{c}{\tau}\,y,\qquad v(\tau):=v^0+\frac{c}{\tau}\,y\qquad(\tau\neq0).
\]
Then:
\begin{enumerate}
\item[(a)] $u(\tau):=w(\tau)+v(\tau)\equiv w^0+v^0$;
$P:=P_{w^0,v^0,t}\in\cH_{T_0}$ for every $T_0>0$, with $P(\infty)=(w^0,v^0)\in\cO$;
hence $P\in\cP(\cO)$ and, for every parameter $\widetilde\theta\in\cP$, the observable
$G^P_{\widetilde\theta}$ is defined on $(T_0,\infty)$ for every $T_0>0$.
\item[(b)] With $\hat\rho:=y^\top A_1y$, the exact dial identity holds:
\begin{equation}\label{eq:exactdial}
\tau\,q\big(w(\tau),v(\tau)\big)+b=\logit t-\frac{c^2\hat\rho}{\tau}
\qquad(\tau\neq0).
\end{equation}
\item[(c)] Consequently, for real $\tau>0$ the first gate
$s_1(\tau)=\sig\big(\tau q(w(\tau),v(\tau))+b\big)$ satisfies $s_1(\tau)\in(0,1)$ and
\begin{equation}\label{eq:absorption}
|s_1(\tau)-t|\le\frac{C}{\tau},\qquad C:=\tfrac14\,c^2|\hat\rho|,
\end{equation}
so $s_1(\tau)\to t$ as $\tau\to+\infty$; and $(w(\tau),v(\tau))\to(w^0,v^0)$.
\item[(d)] All of (b)--(c) hold verbatim for the primed parameter, since $A_1=A_1'$
(Proposition~\ref{prop:A1}); in particular $s_1'=s_1$ along $P$.
\end{enumerate}
\end{lemma}

\begin{proof}
(a)  The sum is constant because the perturbations cancel.  Each component of
$w(\tau),v(\tau)$ is of the form $a+b'/\tau$ ($a,b'$ real constants), holomorphic on
$\{|\tau|>T_0\}$ for any $T_0>0$, bounded there ($|c/\tau|\le|c|/T_0$), and real for
real $\tau$; the limit at $\infty$ is $(w^0,v^0)$, which lies in $\cO$ by
Lemma~\ref{lem:region}(ii).  Hence $P\in\cH_{T_0}$ and $P\in\cP(\cO)$
(Definition~\ref{def:paths}); the observables are defined on $(T_0,\infty)$ as recorded
after Definition~\ref{def:paths}.

(b)  Using $q(w^0,v^0)=0$ (the base lies in $\cQ$) and bilinearity, with $a:=c/\tau$:
\begin{align*}
q\big(w(\tau),v(\tau)\big)
&=(w^0-ay)^\top A_1(v^0+ay)
=q(w^0,v^0)+a\,\big[w^{0\top}A_1y-y^\top A_1v^0\big]-a^2\,y^\top A_1y\\
&=a\,y^\top\big(A_1^\top w^0-A_1v^0\big)-a^2\hat\rho
=a\,y^\top\pi^0-a^2\hat\rho
=a-a^2\hat\rho,
\end{align*}
since $y^\top\pi^0=|\pi^0|^{-2}\pi^{0\top}\pi^0=1$.  Multiplying by $\tau$ and adding
$b$: $\tau q+b=c+b-(c^2/\tau)\hat\rho=\logit t-(c^2/\tau)\hat\rho$.

(c)  For real $\tau>0$ the argument of $\sig$ is real, so $s_1(\tau)\in(0,1)$
(Lemma~\ref{lem:sigmoid}).  Since $|\sig'|\le\frac14$ on $\R$
(\S\ref{subsec:notation}) and $t=\sig(\logit t)$, the mean value theorem applied to
\eqref{eq:exactdial} gives
$|s_1(\tau)-t|=\big|\sig(\logit t-c^2\hat\rho/\tau)-\sig(\logit t)\big|
\le\frac14\,c^2|\hat\rho|/\tau$.  The convergence of the probe point is clear from the
formulas.

(d)  All quantities in (b)--(c) are built from $A_1$, $b$ and the base data; since
$A_1'=A_1$, the first-layer slope $w(\tau)^\top A_1'v(\tau)=q(w(\tau),v(\tau))$ and
hence the first gate coincide for both parameters.
\end{proof}

\begin{auxlemma}[Smooth locus and components]\label{lem:smooth}
Let $S:=\big(\cQ\setminus\{(0,0)\}\big)\times(0,1)$.  Then:
\textup{(a)} $\cQ\setminus\{(0,0)\}$ is exactly the set of points of $\cQ$ at which
$\nabla q=(A_1v,A_1^\top w)\neq0$, and it is a $(2d-1)$-dimensional real-analytic
manifold with graph charts; \textup{(b)} $S$ is locally connected and relatively open
in $\cQ\times(0,1)$; \textup{(c)} $\cU\subseteq S$, and every relatively open subset
$\cW\subseteq\cU$ of $\cQ\times(0,1)$ is open in $S$, and each connected component of
$\cW$ is relatively open in $\cQ\times(0,1)$.
\end{auxlemma}

\begin{proof}
(a) and (b) are established in Steps 3--4 of the proof of Lemma~\ref{lem:region}.
For (c): $\cU\subseteq S$ because points of $\cU$ have $w\neq0$
(Lemma~\ref{lem:region}(i)); a relatively open $\cW\subseteq\cU\subseteq S$ of
$\cQ\times(0,1)$ is open in the subspace $S$; in a locally connected space the
connected components of an open set are open; and a subset open in $S$ is relatively
open in $\cQ\times(0,1)$ because $S$ itself is.
\end{proof}

\subsection{Cascade limits and the trichotomy}

In the next two results, $\Gamma^{(j)}:=(B_j-\varsigma_jV_j)\cdots(B_2-\varsigma_2V_2)$
for $j\ge2$ and $\Gamma^{(1)}:=I$ denote saturated partial products attached to
constants $\varsigma_2,\varsigma_3,\dots$ (empty product conventions as in
\S\ref{subsec:notation}).

\begin{lemma}[Cascade]\label{lem:cascade}
Let $2\le\ell\le L$, and let $\cU_0\subseteq\cU$ be nonempty, relatively open and
such that, for some constants $\varsigma_2,\dots,\varsigma_{\ell-1}\in[0,1]$, along
every dial path $P=P_{w^0,v^0,t}$ with base $(w^0,v^0,t)\in\cU_0$ one has
$s_j(\tau)=\varsigma_j+O(e^{-c_P\tau})$ for $2\le j\le\ell-1$ and some $c_P>0$
(unprimed parameter; the hypothesis is vacuous for $\ell=2$).  Define the polynomial
\[
\Lambda_\ell(w,v,t):=\phi_\ell(t,\varsigma_2,\dots,\varsigma_{\ell-1};w,v).
\]
Then:
\begin{enumerate}
\item[(a)] Along every such dial path,
$\lambda_\ell(\tau)\to\Lambda:=\Lambda_\ell(w^0,v^0,t)$ as $\tau\to+\infty$; and if
$\Lambda\neq0$, then $s_\ell(\tau)=\mathbf 1[\Lambda>0]+O(e^{-c\tau})$ for some $c>0$.
\item[(b)] If $\Lambda_\ell\equiv0$ on $\cU_0$, then there is a polynomial
$\gamma(z)=\gamma_0+\gamma_1z$ of degree $\le1$ such that
\begin{equation}\label{eq:curve}
\phi_\ell(z_1,\varsigma_2,\dots,\varsigma_{\ell-1};w,v)
=\gamma(z_1)\,q(w,v)
\qquad\text{identically in }(z_1,w,v)\in\R^{1+2d};
\end{equation}
equivalently, with $Y:=(\Gamma^{(\ell-1)})^\top A_\ell B_{\ell-1:1}$ and
$T_{\ell-1}(z_1,\varsigma)
:=z_1B_{\ell-1:2}V_1+\sum_{j=2}^{\ell-1}\varsigma_j\,
B_{\ell-1:j+1}V_j\,\Gamma^{(j-1)}(B_1-z_1V_1)$,
\begin{align}
(B_1-z_1V_1)^\top Y&=\gamma(z_1)\,A_1,
\tag{$\heartsuit$}\label{eq:heart}\\
\Sym\Big[(B_1-z_1V_1)^\top(\Gamma^{(\ell-1)})^\top A_\ell\,
T_{\ell-1}(z_1,\varsigma)\Big]&=0,
\tag{$\clubsuit$}\label{eq:club}
\end{align}
identically in $z_1$.  If moreover $\gamma\equiv0$, then along every dial path with
base in $\cU_0$, $s_\ell(\tau)=\alpha+O(e^{-c\tau})$ for every $0<c<c_P$.
\end{enumerate}
\end{lemma}

\begin{proof}
Throughout, fix a base $(w^0,v^0,t)\in\cU_0$, let $P=P_{w^0,v^0,t}$ be its dial path
(Lemma~\ref{lem:dial}), and let $s_j(\tau),w_j(\tau),v_j(\tau),\lambda_j(\tau)$ denote
the recursion quantities of Proposition~\ref{prop:probe} for $\theta$ at
$(w(\tau),v(\tau),\tau)$.

\emph{Specialization identities.}  Specializing the variables
$z=(z_1,\varsigma_2,\dots,\varsigma_{\ell-1})$ in Lemma~\ref{lem:phi}(i) and using
$W_{j-1}(z_1,\varsigma)=\Gamma^{(j-1)}(B_1-z_1V_1)$ (the factors $2,\dots,j-1$ carry the
constants $\varsigma$) gives
\begin{equation}\label{eq:Xspec}
\phi_\ell(z_1,\varsigma;w,v)=w^\top X_\ell(z_1,\varsigma)\,v+w^\top
Y_\ell(z_1,\varsigma)\,w,\quad
X_\ell(z_1,\varsigma)=(B_1-z_1V_1)^\top Y,\quad
Y_\ell(z_1,\varsigma)=W_{\ell-1}(z_1,\varsigma)^\top A_\ell T_{\ell-1}(z_1,\varsigma),
\end{equation}
with $Y$, $T_{\ell-1}(z_1,\varsigma)$ as in the statement (the formula for
$T_{\ell-1}(z_1,\varsigma)$ is the specialization of
$T_{\ell-1}(z)=\sum_{i=1}^{\ell-1}z_iB_{\ell-1:i+1}V_iW_{i-1}(z)$: the $i=1$ summand is
$z_1B_{\ell-1:2}V_1W_0=z_1B_{\ell-1:2}V_1$, and the $i=j\ge2$ summand is
$\varsigma_jB_{\ell-1:j+1}V_j\Gamma^{(j-1)}(B_1-z_1V_1)$).  In particular every entry of
$X_\ell(z_1,\varsigma)$ is an affine polynomial in $z_1$ and every entry of
$Y_\ell(z_1,\varsigma)$ has degree $\le2$ in $z_1$.

\emph{(a)}  Let $K:=[0,1]^{\ell-1}\times\bar B$, where $\bar B\subset\R^{2d}$ is a
closed ball containing $(w(\tau),v(\tau))$ for all $\tau\ge1$ (the path is bounded) ---
a compact convex set containing, for $\tau\ge1$, both argument tuples below (all gates
lie in $[0,1]$, being values of $\sig$ at real arguments, and
$\varsigma_j\in[0,1]$).  The polynomial $\phi_\ell$ is Lipschitz on $K$, with constant
$C_K$ say.  Hence
\[
\Delta(\tau):=\phi_\ell\big(s_1(\tau),s_2(\tau),\dots,s_{\ell-1}(\tau);
w(\tau),v(\tau)\big)
-\phi_\ell\big(s_1(\tau),\varsigma_2,\dots,\varsigma_{\ell-1};w(\tau),v(\tau)\big)
\]
satisfies $|\Delta(\tau)|\le C_K\max_{2\le j\le\ell-1}|s_j(\tau)-\varsigma_j|
=O(e^{-c_P\tau})$ by hypothesis (for $\ell=2$, $\Delta\equiv0$).  By
Lemma~\ref{lem:lambdaphi},
\[
\lambda_\ell(\tau)=\phi_\ell\big(s_1(\tau),\dots,s_{\ell-1}(\tau);
w(\tau),v(\tau)\big)
=\Lambda_\ell\big(w(\tau),v(\tau),s_1(\tau)\big)+\Delta(\tau).
\]
Since $s_1(\tau)\to t$ and $(w(\tau),v(\tau))\to(w^0,v^0)$ (Lemma~\ref{lem:dial}) and
$\Lambda_\ell$ is a polynomial, hence continuous,
$\lambda_\ell(\tau)\to\Lambda_\ell(w^0,v^0,t)=\Lambda$.

Suppose $\Lambda\neq0$ and choose $T_1\ge1$ with $|\lambda_\ell(\tau)-\Lambda|\le
|\Lambda|/4$ for $\tau>T_1$.  If $\Lambda>0$: for $\tau>T_1$,
$\tau\lambda_\ell(\tau)+b\ge\tfrac34\Lambda\tau+b\ge\tfrac34\Lambda\tau$, so by
\eqref{eq:tail},
$0<1-s_\ell(\tau)\le e^{-(\tau\lambda_\ell+b)}\le e^{-\frac34\Lambda\tau}$:
$s_\ell(\tau)=1+O(e^{-c\tau})$ with $c:=\tfrac34\Lambda$.  If $\Lambda<0$: for
$\tau>\max(T_1,\,4b/|\Lambda|)$,
$\tau\lambda_\ell(\tau)+b\le-\tfrac34|\Lambda|\tau+b\le-\tfrac12|\Lambda|\tau$, so by
\eqref{eq:tail}, $0<s_\ell(\tau)\le e^{\tau\lambda_\ell+b}\le e^{-|\Lambda|\tau/2}$:
$s_\ell(\tau)=0+O(e^{-c\tau})$ with $c:=|\Lambda|/2$.  This is (a).

\emph{(b)}  Assume $\Lambda_\ell\equiv0$ on $\cU_0$.  Fix any point of $\cU_0$ and, by
the product-neighbourhood clause of Lemma~\ref{lem:region}, a product neighbourhood
$N=U\times J\subseteq\cU_0$, with $J$ a nondegenerate open interval and $U$ a
$t$-independent relatively open subset of $\cQ$ containing a point with $w\neq0$.

\emph{Per-slice rigidity.}  Fix $\bar t\in J$.  For every $(w,v)\in U$ we have
$(w,v,\bar t)\in N\subseteq\cU_0$, so by hypothesis and \eqref{eq:Xspec},
\[
0=\Lambda_\ell(w,v,\bar t)
=w^\top X_\ell(\bar t,\varsigma)\,v+w^\top Y_\ell(\bar t,\varsigma)\,w
\qquad\big((w,v)\in U\big).
\]
Since $\det A_1\neq0$ and $d\ge2$, Lemma~\ref{lem:quadA} applies (with
$X=X_\ell(\bar t,\varsigma)$, $Y=Y_\ell(\bar t,\varsigma)$):
\[
X_\ell(\bar t,\varsigma)=c(\bar t)\,A_1\quad\text{for a scalar }c(\bar t),
\qquad\Sym Y_\ell(\bar t,\varsigma)=0\qquad(\bar t\in J).
\]
The scalar is unique ($A_1\neq0$) and is given by the trace formula
$c(\bar t)=\operatorname{tr}\!\big(A_1^\top X_\ell(\bar t,\varsigma)\big)/\|A_1\|_F^2$,
where $\|A_1\|_F^2:=\operatorname{tr}(A_1^\top A_1)>0$.

\emph{Polynomial upgrade.}  Each entry of $X_\ell(z_1,\varsigma)$ is affine in $z_1$
\eqref{eq:Xspec}, so
\[
\gamma(z_1):=\frac{\operatorname{tr}\!\big(A_1^\top X_\ell(z_1,\varsigma)\big)}
{\|A_1\|_F^2}
\]
is a polynomial of degree $\le1$, say $\gamma=\gamma_0+\gamma_1z$, with
$\gamma(\bar t)=c(\bar t)$ for $\bar t\in J$.  Each entry of
$X_\ell(z_1,\varsigma)-\gamma(z_1)A_1$ is a polynomial of degree $\le1$ in $z_1$
vanishing at every point of the infinite set $J$, hence identically zero: this is
\eqref{eq:heart}.  Likewise each entry of $\Sym Y_\ell(z_1,\varsigma)$ is a polynomial
of degree $\le2$ in $z_1$ vanishing on $J$, hence identically zero: this is
\eqref{eq:club}.  Adding the bilinear identity
$w^\top X_\ell(z_1,\varsigma)v=\gamma(z_1)\,w^\top A_1v=\gamma(z_1)q(w,v)$ and the
quadratic identity $w^\top Y_\ell(z_1,\varsigma)w
=w^\top\Sym\big(Y_\ell(z_1,\varsigma)\big)w=0$ to \eqref{eq:Xspec} gives
\eqref{eq:curve}, an identity of polynomials in \emph{all} of $(z_1,w,v)$: it holds on
$\R^{1+2d}$, not merely over the quadric.  (This upgrade is what makes
\eqref{eq:curve} usable along dial paths, whose points leave $\cQ$ for finite $\tau$.)
Conversely, \eqref{eq:heart}--\eqref{eq:club} imply \eqref{eq:curve} by the same
substitution, and \eqref{eq:curve} implies them: differentiating \eqref{eq:curve} in
$v$ at fixed $(z_1,w)$ gives $X_\ell(z_1,\varsigma)^\top w=\gamma(z_1)A_1^\top w$ for
all $w$, i.e.\ \eqref{eq:heart}, and then $w^\top Y_\ell(z_1,\varsigma)w\equiv0$ gives
\eqref{eq:club} by polarization.

\emph{The branch $\gamma\equiv0$.}  Suppose $\gamma\equiv0$, and let $P$ be a dial
path with base $(w^0,v^0,t)\in\cU_0$.  Evaluating \eqref{eq:curve} at
$(z_1,w,v)=\big(s_1(\tau),w(\tau),v(\tau)\big)$ --- legitimate by the previous
paragraph --- gives
$\phi_\ell\big(s_1(\tau),\varsigma;w(\tau),v(\tau)\big)
=\gamma\big(s_1(\tau)\big)\,q\big(w(\tau),v(\tau)\big)=0$, whence, in the notation of
part (a), $\lambda_\ell(\tau)=\Delta(\tau)=O(e^{-c_P\tau})$ and
$\tau\lambda_\ell(\tau)=O(\tau e^{-c_P\tau})$.  Since
$|\sig(b+x)-\sig(b)|\le|x|/4$ ($|\sig'|\le\tfrac14$, \S\ref{subsec:notation}) and
$\alpha=\sig(b)$,
\[
s_\ell(\tau)=\sig\big(\tau\lambda_\ell(\tau)+b\big)
=\alpha+O\big(\tau e^{-c_P\tau}\big)=\alpha+O(e^{-c\tau})
\quad\text{for every }0<c<c_P. \qedhere
\]
\end{proof}

\begin{proposition}[Trichotomy]\label{prop:trichotomy}
Assume \eqref{eq:SH}.  There exist a nonempty, relatively open, connected set
$\cU_*\subseteq\cU$ and constants $\varsigma_2,\dots,\varsigma_L\in\{0,1,\alpha\}$ such
that along every dial path $P=P_{w^0,v^0,t}$ with base $(w^0,v^0,t)\in\cU_*$ there is
$c_P>0$ with
\[
s_j(\tau)=\varsigma_j+O(e^{-c_P\tau})\ \ (2\le j\le L)
\qquad\text{and}\qquad
s'_j(\tau)=1+O(e^{-c_P\tau})\ \ (2\le j\le L).
\]
\end{proposition}

\begin{proof}
\emph{Primed side.}  We claim: for every base $(w^0,v^0,t)\in\cU$ and every
$2\le\ell\le L$, along the dial path $P_{w^0,v^0,t}$ one has
$s'_\ell(\tau)=1+O(e^{-c'\tau})$ for some path-dependent $c'>0$.  Induct on $\ell$,
assuming (vacuously for $\ell=2$) the statement for $2\le j\le\ell-1$.  The proof of
Lemma~\ref{lem:cascade}(a) is parameter-generic --- it uses only
Lemmas~\ref{lem:phi}, \ref{lem:lambdaphi} (valid for any tuple),
Lemma~\ref{lem:dial} ($s_1'=s_1\to t$) and \eqref{eq:tail} --- so it applies verbatim
to $\theta'$ with the constants $\varsigma'_j:=1$:
$\lambda'_\ell(\tau)\to\phi'_\ell(t,1,\dots,1;w^0,v^0)>0$, the positivity by
Lemma~\ref{lem:region}(iii) since $(w^0,v^0,t)\in\cU$, and the saturation estimate of
the case $\Lambda>0$ gives $s'_\ell(\tau)=1+O(e^{-c\tau})$.  Take $c'$ to be the
minimum over the finitely many levels.

\emph{Unprimed side.}  We construct, by induction on $\ell=1,\dots,L$, nonempty
relatively open connected sets
$\cU=:\cU^{(1)}_*\supseteq\cU^{(2)}_*\supseteq\dots\supseteq\cU^{(L)}_*$ and constants
$\varsigma_2,\dots,\varsigma_L\in\{0,1,\alpha\}$ such that, for each $\ell$, along
every dial path with base in $\cU^{(\ell)}_*$ one has
$s_j(\tau)=\varsigma_j+O(e^{-c_P\tau})$ for $2\le j\le\ell$ and some path-dependent
$c_P>0$.  The proposition follows with $\cU_*:=\cU^{(L)}_*$, combining with the primed
side ($\cU_*\subseteq\cU$) and shrinking each path's exponent to the minimum of the
finitely many constants involved.

\emph{Base $\ell=1$.}  $\cU^{(1)}_*:=\cU$ is nonempty, relatively open and connected
(Lemma~\ref{lem:region}); the gate condition is vacuous.

\emph{Step $\ell\ge2$.}  Suppose $\cU^{(\ell-1)}_*$ and
$\varsigma_2,\dots,\varsigma_{\ell-1}$ constructed.  The hypotheses of
Lemma~\ref{lem:cascade} hold with $\cU_0:=\cU^{(\ell-1)}_*$.  The function
$\Lambda_\ell(w,v,t)=\phi_\ell(t,\varsigma_2,\dots,\varsigma_{\ell-1};w,v)$ is a
polynomial in $(w,v,t)$, hence continuous.

\emph{Case 1: $\Lambda_\ell\not\equiv0$ on $\cU^{(\ell-1)}_*$.}  Pick
$(w_1,v_1,t_1)\in\cU^{(\ell-1)}_*$ with $\Lambda_\ell(w_1,v_1,t_1)\neq0$ and let
$\cU^{(\ell)}_*$ be the connected component of this point in
$\cW:=\{\Lambda_\ell\neq0\}\cap\cU^{(\ell-1)}_*$.  The set $\cW$ is relatively open in
$\cQ\times(0,1)$ (intersection of the open $\{\Lambda_\ell\neq0\}\subseteq\R^{2d+1}$
with the relatively open $\cU^{(\ell-1)}_*$) and contained in $\cU$; by Auxiliary
Lemma~\ref{lem:smooth}(c), the component $\cU^{(\ell)}_*$ is relatively open; it is
nonempty and connected by construction.  Since $\Lambda_\ell$ is continuous and
zero-free on the connected set $\cU^{(\ell)}_*$, its image is a connected subset of
$\R\setminus\{0\}$, hence $\Lambda_\ell$ has constant sign there; set
$\varsigma_\ell:=\mathbf 1\big[\Lambda_\ell>0\text{ on }\cU^{(\ell)}_*\big]\in\{0,1\}$.
By Lemma~\ref{lem:cascade}(a), along every dial path with base in
$\cU^{(\ell)}_*\subseteq\cU^{(\ell-1)}_*$ we get
$s_\ell(\tau)=\varsigma_\ell+O(e^{-c\tau})$; combined with the inductive rates for
$s_2,\dots,s_{\ell-1}$ (bases lie in $\cU^{(\ell-1)}_*$) and taking the minimum of the
finitely many per-path exponents, the inductive statement holds at level $\ell$.

\emph{Case 2: $\Lambda_\ell\equiv0$ on $\cU^{(\ell-1)}_*$.}
Lemma~\ref{lem:cascade}(b) provides $\gamma(z)=\gamma_0+\gamma_1z$ with
\eqref{eq:heart}, \eqref{eq:club}.  We show $\gamma\equiv0$; the final clause of
Lemma~\ref{lem:cascade}(b) then gives $s_\ell(\tau)=\alpha+O(e^{-c\tau})$ along every
dial path with base in $\cU^{(\ell)}_*:=\cU^{(\ell-1)}_*$, and the step closes with
$\varsigma_\ell:=\alpha$.

Matching coefficients of the affine (in $z_1$) matrix identity \eqref{eq:heart},
$(B_1-z_1V_1)^\top Y=(\gamma_0+\gamma_1z_1)A_1$:
\[
\text{(c0)}\quad B_1^\top Y=\gamma_0A_1,
\qquad\qquad
\text{(c1)}\quad -V_1^\top Y=\gamma_1A_1 .
\]

\emph{K2: $\gamma_1=0$, $\gamma_0\neq0$ is impossible.}  By (c1), $V_1^\top Y=0$.  By
(c0) and $\det A_1\neq0$,
$\det(B_1^\top)\det Y=\gamma_0^{\,d}\det A_1\neq0$, so $Y$ is invertible; hence
$V_1^\top=0$, i.e.\ $V_1=0$ --- contradicting Corollary~\ref{cor:depth}.

\emph{K3: $\gamma_1\neq0$ is impossible.}  By (c1), $V_1^\top Y=-\gamma_1A_1$ is
invertible, so $V_1$ and $Y$ are invertible.  If $\gamma_0=0$, then (c0) gives
$B_1^\top Y=0$ with $Y$ invertible, so $B_1=0$; but then
$Y=(\Gamma^{(\ell-1)})^\top A_\ell B_{\ell-1:1}
=(\Gamma^{(\ell-1)})^\top A_\ell B_{\ell-1:2}B_1=0$, a contradiction.  Hence
$\gamma_0\neq0$ and, by (c0), $B_1$ is invertible.  Combining (c0) and (c1):
$V_1^\top Y=-(\gamma_1/\gamma_0)\gamma_0A_1=-(\gamma_1/\gamma_0)B_1^\top Y$, so
$\big(V_1^\top+(\gamma_1/\gamma_0)B_1^\top\big)Y=0$ and, $Y$ being invertible,
$V_1=cB_1$ with $c:=-\gamma_1/\gamma_0\neq0$.  With $B_1=I+V_1$ this gives
$(1-c)V_1=cI$; $c\neq1$ (otherwise $0=I$), so
\[
V_1=\nu I,\qquad B_1=(1+\nu)I,\qquad
\nu:=\frac{c}{1-c}\neq0,\qquad 1+\nu=\frac{1}{1-c}\neq0 .
\]
Insert these scalars into \eqref{eq:club}.  With $\beta(z):=(1+\nu)-\nu z$ (so
$\beta\not\equiv0$, having $\beta(0)=1+\nu\neq0$) we have $B_1-z_1V_1=\beta(z_1)I$ and
\[
T_{\ell-1}(z_1,\varsigma)=\nu z_1\,B_{\ell-1:2}+\beta(z_1)\,\widetilde Q,
\qquad
\widetilde Q:=\sum_{j=2}^{\ell-1}\varsigma_j\,B_{\ell-1:j+1}V_j\,\Gamma^{(j-1)}
\]
(for $\ell=2$, $\widetilde Q=0$).  Hence, by linearity of $\Sym$,
\eqref{eq:club} reads
\[
\beta(z)\Big[\nu z\,S_1+\beta(z)\,S_2\Big]=0\quad\text{entrywise, identically in }z,
\qquad
S_1:=\Sym\big((\Gamma^{(\ell-1)})^\top A_\ell B_{\ell-1:2}\big),\quad
S_2:=\Sym\big((\Gamma^{(\ell-1)})^\top A_\ell\widetilde Q\big).
\]
For each entry $(a,b)$:
$\beta(z)\big(\nu z(S_1)_{ab}+\beta(z)(S_2)_{ab}\big)=0$ in $\R[z]$; since
$\beta\neq0$ and $\R[z]$ is an integral domain,
$\nu z(S_1)_{ab}+\beta(z)(S_2)_{ab}=0$.  The constant coefficient gives
$(1+\nu)(S_2)_{ab}=0$, so $S_2=0$; the coefficient of $z$ gives
$\nu(S_1)_{ab}-\nu(S_2)_{ab}=0$, so $S_1=S_2=0$.  On the other hand, evaluating
\eqref{eq:heart} at $z_1=0$, i.e.\ (c0), with $B_1^\top=(1+\nu)I$ and
$B_{\ell-1:1}=B_{\ell-1:2}B_1=(1+\nu)B_{\ell-1:2}$:
\[
(1+\nu)^2\,(\Gamma^{(\ell-1)})^\top A_\ell B_{\ell-1:2}=\gamma_0A_1 .
\]
Taking $\Sym$ of both sides: $(1+\nu)^2S_1=\gamma_0\Sym(A_1)$ --- the same $S_1$ as
above --- whence $\Sym(A_1)=0$ ($S_1=0$, $\gamma_0\neq0$, $(1+\nu)^2\neq0$),
contradicting (G1).

Both K2 and K3 being impossible, $\gamma_1=0$ and $\gamma_0=0$, i.e.\
$\gamma\equiv0$, as required.  (For $\ell=2$ the conventions give $\Gamma^{(1)}=I$,
$\widetilde Q=0$, $S_2=0$ and $Y=A_2B_1$; the argument is unchanged.)  Since at each
level exactly one of Case 1 / Case 2 holds, the constants
$\varsigma_\ell\in\{0,1,\alpha\}$ always exist, and the induction terminates at
$\ell=L$.
\end{proof}

\subsection{Matching and the $K=I$ extraction}

\begin{proposition}[Matching]\label{prop:matching}
Assume \eqref{eq:SH}.  Then $B_{L:1}=B'_{L:1}$, $B_1=B_1'$ and $V_1=V_1'$.
\end{proposition}

\begin{proof}
Let $\cU_*$ and $\varsigma_2,\dots,\varsigma_L\in\{0,1,\alpha\}$ be as in
Proposition~\ref{prop:trichotomy}, and retain
$\Gamma^{(j)}=(B_j-\varsigma_jV_j)\cdots(B_2-\varsigma_2V_2)$, $\Gamma^{(1)}=I$
(\emph{unprimed} matrices).

\emph{Step 0 (product neighbourhood).}  By the product-neighbourhood clause of
Lemma~\ref{lem:region} applied to $\cU_0:=\cU_*$, fix
$N=U\times J\subseteq\cU_*$ with $J\subseteq(0,1)$ a nondegenerate open interval and
$U\subseteq\cQ$ relatively open, $t$-independent, containing a point with $w\neq0$.

\emph{Step 1 (limits along dial paths, unprimed).}  Fix $(w^0,v^0,t)\in N$ and let
$P$ be its dial path.  By Proposition~\ref{prop:probe}, Lemma~\ref{lem:lambdaphi} and
Lemma~\ref{lem:phi}(i),
\[
G^P_\theta(\tau)=v_L(\tau)
=B_{L:1}v(\tau)+T_L\big(s_1(\tau),\dots,s_L(\tau)\big)\,w(\tau),
\]
with $T_L(z)=\sum_{i=1}^Lz_iB_{L:i+1}V_iW_{i-1}(z)$, a matrix polynomial; the
right-hand side is jointly continuous in the gates and the probe point.  By
Lemma~\ref{lem:dial} and Proposition~\ref{prop:trichotomy}, as $\tau\to+\infty$:
$s_1(\tau)\to t$, $s_j(\tau)\to\varsigma_j$ ($j\ge2$),
$(w(\tau),v(\tau))\to(w^0,v^0)$.  Hence, using
$W_{i-1}(t,\varsigma)=\Gamma^{(i-1)}(B_1-tV_1)$,
\begin{equation}\label{eq:limtheta}
G^P_\theta(\tau)\ \longrightarrow\
B_{L:1}v^0+\Big[t\,B_{L:2}V_1+D\,(B_1-tV_1)\Big]w^0,
\qquad
D:=\sum_{i=2}^{L}\varsigma_i\,B_{L:i+1}V_i\,\Gamma^{(i-1)} .
\end{equation}

\emph{Step 2 (limits, primed; telescoping).}  The same computation for $\theta'$,
whose gates converge to $(t,1,\dots,1)$, together with transparency
(Lemma~\ref{lem:phi}(iv): $W'_{i-1}(t,1,\dots,1)=B_1'-tV_1'$) gives
\begin{equation}\label{eq:limthetaprime-raw}
G^P_{\theta'}(\tau)\ \longrightarrow\
B'_{L:1}v^0+\Big[t\,B'_{L:2}V_1'
+\Big(\textstyle\sum_{i=2}^{L}B'_{L:i+1}V_i'\Big)(B_1'-tV_1')\Big]w^0 .
\end{equation}
By the skip identity $V_i'=B_i'-I$ \eqref{eq:transparency}, the sum telescopes:
$\sum_{i=2}^{L}B'_{L:i+1}(B_i'-I)
=\sum_{i=2}^{L}\big(B'_{L:i}-B'_{L:i+1}\big)=B'_{L:2}-I$.  Expanding,
\[
t\,B'_{L:2}V_1'+(B'_{L:2}-I)(B_1'-tV_1')
=t\,B'_{L:2}V_1'+B'_{L:1}-t\,B'_{L:2}V_1'-B_1'+tV_1'
=B'_{L:1}-B_1'+tV_1',
\]
so
\begin{equation}\label{eq:limthetaprime}
G^P_{\theta'}(\tau)\ \longrightarrow\
B'_{L:1}v^0+\big[B'_{L:1}-B_1'+t\,V_1'\big]w^0 :
\end{equation}
on the primed side the unknown saturation data is entirely absent and the
$t$-dependence is affine with coefficient $V_1'$.

\emph{Step 3 (matching the limits).}  $P\in\cP(\cO)$ by Lemma~\ref{lem:dial}(a), so
the agreement clause of \eqref{eq:SH} gives $G^P_\theta=G^P_{\theta'}$ on
$(T_P,\infty)$; letting $\tau\to+\infty$ and using
\eqref{eq:limtheta}, \eqref{eq:limthetaprime}:
\begin{equation}\label{eq:diff}
\big(B_{L:1}-B'_{L:1}\big)\,v^0
+\Big[t\,B_{L:2}V_1+D(B_1-tV_1)-B'_{L:1}+B_1'-t\,V_1'\Big]\,w^0=0
\qquad\text{for every }(w^0,v^0,t)\in N=U\times J .
\end{equation}

\emph{Step 4 (separating the probe coefficients).}  Fix $t\in J$ and apply
Lemma~\ref{lem:quadB} to \eqref{eq:diff} on the relatively open
$U\subseteq\cQ$ (which contains a point with $w\neq0$; $\det A_1\neq0$, $d\ge2$): with
$C:=B_{L:1}-B'_{L:1}$ and $E_t$ the bracketed matrix,
\[
C=0\qquad\text{and}\qquad E_t=0\quad\text{for every }t\in J .
\]
In particular $B_{L:1}=B'_{L:1}$.

\emph{Step 5 (coefficients in $t$).}  $E_t$ is affine in $t$:
\[
E_t=\big[B_{L:2}V_1-DV_1-V_1'\big]\,t+\big[DB_1-B'_{L:1}+B_1'\big].
\]
An affine matrix function of $t$ vanishing on the infinite set $J$ has both
coefficients zero:
\begin{equation}\label{eq:coeffs}
\big(B_{L:2}-D\big)V_1=V_1',
\qquad
D\,B_1=B'_{L:1}-B_1' .
\end{equation}

\emph{Step 6 (the $K=I$ cancellation).}  Set $K:=B_{L:2}-D$.  The first identity of
\eqref{eq:coeffs} is $KV_1=V_1'$.  Substituting $D=B_{L:2}-K$ into the second:
$B_{L:2}B_1-KB_1=B'_{L:1}-B_1'$, i.e.\ $B_{L:1}-KB_1=B'_{L:1}-B_1'$; since
$B_{L:1}=B'_{L:1}$ (Step 4), this reads $KB_1=B_1'$.  Subtracting the two identities
and using the skip identity on \emph{both} sides
($B_1-V_1=I=B_1'-V_1'$, \eqref{eq:transparency}):
\[
K=K(B_1-V_1)=KB_1-KV_1=B_1'-V_1'=I .
\]
Hence $V_1=KV_1=V_1'$ and $B_1=KB_1=B_1'$.
\end{proof}

\begin{remark}[The case $L=2$]\label{rem:L2}
All conventions degenerate gracefully at $L=2$: the trichotomy concerns the single
level $\ell=2$ (and the cascade hypothesis at $\ell=2$ is vacuous, $\Delta\equiv0$);
$\Gamma^{(1)}=I$, $D=\varsigma_2V_2$, $B_{L:2}=B_2$, and Step 6 reads
$(B_2-\varsigma_2V_2)V_1=V_1'$, $(B_2-\varsigma_2V_2)B_1=B_2'$\,$\cdots$\,$=B_1'$
modulo $B_{2:1}=B'_{2:1}$, with $K=B_2-\varsigma_2V_2=I$ extracted exactly as in the
general case.  Note that $K=I$ is a statement about the indicated \emph{combination}
of $\theta$'s deeper layers and saturation constants; the trichotomy is still needed
beforehand, to know that the limits \eqref{eq:limtheta} exist at all.
\end{remark}

\section{Step 3: realization, the sweep, and the proof of Theorem~\ref{thm:IDL}}
\label{sec:descent}

After \S\S\ref{sec:A1}--\ref{sec:cascade}, the entire first layer is common under
\eqref{eq:SH}: $(V_1,A_1)=(V_1',A_1')$, and, by (G3) and
Proposition~\ref{prop:matching}, $\det B_1=\det B_1'\neq0$.  Throughout this section we
write, for the common first layer,
\[
\Phi_t(w,v):=\big((B_1-tV_1)w,\ B_1v+tV_1w\big),
\qquad
\vartheta_0(w,v,t):=\pi(w,v)^\top(B_1-tV_1)^{-1}V_1w
\]
(the latter wherever $\det(B_1-tV_1)\neq0$); these maps, and the quadric data $q,\pi$,
agree for $\theta$ and $\theta'$.  Given a path $P=(w,v)\in\cH_T$, its
\emph{first-layer effective path} is
\[
\tau\longmapsto\Big((B_1-s_1(\tau)V_1)\,w(\tau),\ \ B_1v(\tau)+s_1(\tau)V_1w(\tau)\Big),
\qquad
s_1(\tau)=\sig\big(\tau\,q(w(\tau),v(\tau))+b\big),
\]
defined for real $\tau>T$; by Lemma~\ref{lem:peeling}, the depth-$L$ observable of
either parameter along $P$ equals the depth-$(L-1)$ observable of its tail along the
effective path.

\subsection{Realization}

\begin{lemma}[Realization]\label{lem:realization}
Let $(w^0,v^0)\in\R^{2d}$ and $t\in(0,1)$ satisfy $q(w^0,v^0)=0$,
$\pi^0:=\pi(w^0,v^0)\neq0$, $V_1w^0\neq0$, $\det(B_1-tV_1)\neq0$, and
$\vartheta_0(w^0,v^0,t)\neq0$.  Let $\widetilde P=(\widetilde w,\widetilde v)\in
\cH_{T_0}$ with $\widetilde P(\infty)=\Phi_t(w^0,v^0)$.  Then there exist $T\ge T_0$
and $P=(w,v)\in\cH_T$ with $P(\infty)=(w^0,v^0)$ such that for every real $\tau>T$,
\[
\big(B_1-s_1(\tau)V_1\big)\,w(\tau)=\widetilde w(\tau),
\qquad
B_1v(\tau)+s_1(\tau)V_1w(\tau)=\widetilde v(\tau),
\]
where $s_1(\tau)=\sig\big(\tau\,q(w(\tau),v(\tau))+b\big)$.  In other words, the
first-layer effective path of $P$ equals $\widetilde P$ on $(T,\infty)$.
\end{lemma}

\begin{proof}
Write $\vartheta_0:=\vartheta_0(w^0,v^0,t)\neq0$ and $c:=\logit t-b$.  The proof is a
scalar Banach fixed point in the dial value of the first gate.

\emph{Step 0 (real symmetry).}  Call a function $f$, holomorphic on
$\{|\tau|>T'\}$, \emph{real-symmetric} if
$\overline{f(\bar\tau)}=f(\tau)$ there.  Members of $\cH_{T'}$ are real-symmetric: for
such $f$, the function $g(\tau):=\overline{f(\bar\tau)}$ is holomorphic on the same set
(a standard verification from the power series, or from the Cauchy--Riemann equations),
which is a \emph{domain} (the complement of a closed disc is open and connected), and
$g=f$ on $(T',\infty)$ since $f$ is real there; the identity theorem gives $g\equiv f$.
The extensions to $\tau=\infty$ inherit the symmetry.  Throughout, all identity-theorem
arguments are applied componentwise to vector-valued functions.

\emph{Step 1 (the auxiliary family $v_S$).}  Since $\det B_1\neq0$ ((G3), the first
layer being common by Proposition~\ref{prop:matching}), define
\[
u(\tau):=B_1^{-1}\big(\widetilde w(\tau)+\widetilde v(\tau)\big),
\]
holomorphic, bounded and real-symmetric on $\{|\tau|>T_0\}$, with
\[
u(\infty)=B_1^{-1}\big((B_1-tV_1)w^0+B_1v^0+tV_1w^0\big)
=B_1^{-1}B_1(w^0+v^0)=w^0+v^0=:u^0 .
\]
Set $M:=(B_1-tV_1)^{-1}V_1$ and $\varsigma_*:=\tfrac1{2(1+\|M\|)}>0$ (operator norm).
For $S\in\C$ with $|S-t|<\varsigma_*$, the factorization
$B_1-SV_1=(B_1-tV_1)\big(I-(S-t)M\big)$ and the Neumann series
$\big(I-(S-t)M\big)^{-1}=\sum_{k\ge0}(S-t)^kM^k$ (convergent since
$|S-t|\,\|M\|<\tfrac12$) show that $B_1-SV_1$ is invertible and
$S\mapsto(B_1-SV_1)^{-1}$ is holomorphic on $D(t,\varsigma_*)$, with \emph{real} matrix
Taylor coefficients at the real point $t$ ($M$ and $(B_1-tV_1)^{-1}$ are real).  Define
\[
v_S(\tau):=(B_1-SV_1)^{-1}\big(\widetilde v(\tau)-S\,V_1u(\tau)\big),
\]
i.e.\ the function $V(\zeta,S):=v_S(1/\zeta)$ (with the convention $1/0=\infty$,
the paths extending across $\infty$ by Definition~\ref{def:paths}), jointly holomorphic
on $D(0,1/T_0)\times D(t,\varsigma_*)$ as a product and composition of holomorphic
maps.  By construction the exact identity
\begin{equation}\label{eq:vSident}
B_1v_S(\tau)+S\,V_1\big(u(\tau)-v_S(\tau)\big)=\widetilde v(\tau)
\end{equation}
holds: $(B_1-SV_1)v_S=\widetilde v-SV_1u$ rearranges to \eqref{eq:vSident}.  Values
and derivative at $(\zeta,S)=(0,t)$:
\[
v_t(\infty)=(B_1-tV_1)^{-1}\big(B_1v^0+tV_1w^0-tV_1(w^0+v^0)\big)
=(B_1-tV_1)^{-1}(B_1-tV_1)v^0=v^0,
\]
and, differentiating $(B_1-SV_1)v_S=\widetilde v-SV_1u$ in $S$ and solving,
$\partial_Sv_S=-(B_1-SV_1)^{-1}V_1\,(u-v_S)$, so that
\begin{equation}\label{eq:mu}
\mu:=\partial_Sv_S\big|_{(\infty,\,t)}=-(B_1-tV_1)^{-1}V_1\,(u^0-v^0)
=-(B_1-tV_1)^{-1}V_1w^0 .
\end{equation}

\emph{Claim (real symmetry of $V$): }
$\overline{V(\bar\zeta,\bar S)}=V(\zeta,S)$ on $D(0,1/T_0)\times D(t,\varsigma_*)$.
Indeed, fix first a real $S\in(t-\varsigma_*,t+\varsigma_*)$: the maps
$\zeta\mapsto V(\zeta,S)$ and $\zeta\mapsto\overline{V(\bar\zeta,S)}$ are holomorphic
on $D(0,1/T_0)$ and agree on $(0,1/T_0)$ (for real $\zeta>0$, $\tau=1/\zeta\in
(T_0,\infty)$, where $\widetilde v(\tau)$, $u(\tau)$ are real by Step 0 and the matrix
$(B_1-SV_1)^{-1}$ is real, having real Taylor coefficients at real $S$); the identity
theorem gives agreement on $D(0,1/T_0)$.  Now fix $\zeta\in D(0,1/T_0)$: the maps
$S\mapsto V(\zeta,S)$ and $S\mapsto\overline{V(\bar\zeta,\bar S)}$ are holomorphic on
$D(t,\varsigma_*)$ and, by the first step, agree for all real
$S\in(t-\varsigma_*,t+\varsigma_*)$; the identity theorem concludes.
\hfill$\diamond$

\emph{Step 2 (the scalar data).}  Define
\[
Q(\tau,\varsigma):=q\big(u(\tau)-v_{t+\varsigma}(\tau),\ v_{t+\varsigma}(\tau)\big),
\qquad
\widehat Q(\zeta,\varsigma):=Q(1/\zeta,\varsigma),
\]
jointly holomorphic on $D(0,1/T_0)\times D(0,\varsigma_*)$ (composition of the
polynomial $q$ with holomorphic vector functions), real-symmetric in the two-variable
sense ($\overline{\widehat Q(\bar\zeta,\bar\varsigma)}=\widehat Q(\zeta,\varsigma)$, by
the Claim and realness of $q$'s coefficients).  Its anchor values:
\[
\widehat Q(0,0)=q(u^0-v^0,v^0)=q(w^0,v^0)=0,
\]
and, by bilinearity of $q$ and \eqref{eq:mu},
\[
\partial_\varsigma\widehat Q(0,0)=q(-\mu,v^0)+q(w^0,\mu)
=\mu^\top\big(A_1^\top w^0-A_1v^0\big)
=\pi^{0\top}\mu=-\,\pi^{0\top}(B_1-tV_1)^{-1}V_1w^0=-\,\vartheta_0\neq0 .
\]

Set $\delta_1\le\tfrac1{2T_0}$ and
$\varsigma_1\le\tfrac14\min(\varsigma_*,t,1-t)$, to be shrunk once below.  Define the
divided difference
\[
\widehat\cR(\zeta,\varsigma):=\frac{1}{2\pi i}\oint_{|\omega|=3\varsigma_1}
\frac{\widehat Q(\zeta,\omega)}{(\omega-\varsigma)\,\omega}\,d\omega,
\qquad(\zeta,\varsigma)\in D(0,1/T_0)\times D(0,3\varsigma_1),
\]
which is jointly holomorphic (differentiation under the integral) and satisfies, by
the Cauchy formula,
$\varsigma\,\widehat\cR(\zeta,\varsigma)
=\widehat Q(\zeta,\varsigma)-\widehat Q(\zeta,0)$ and
$\widehat\cR(\zeta,0)=\partial_\varsigma\widehat Q(\zeta,0)$; in particular
$\widehat\cR(0,0)=-\vartheta_0$.  It is real-symmetric: off $\varsigma=0$ it is the
quotient of real-symmetric functions by $\varsigma$, and the symmetry extends by
continuity.  Next, $a_0(\zeta):=\widehat Q(\zeta,0)$ is holomorphic on $D(0,1/T_0)$
with $a_0(0)=0$ --- this is exactly where the hypothesis $q(w^0,v^0)=0$ is consumed
--- so
\[
\widehat h(\zeta):=\frac{a_0(\zeta)}{\zeta}
\]
extends holomorphically to $D(0,1/T_0)$ (removable singularity), is real-symmetric, and
satisfies $\widehat h(1/\tau)=\tau\,Q(\tau,0)$.  Set
$H_\infty:=\sup_{|\zeta|\le\delta_1}|\widehat h|<\infty$.  Finally,
$\logit$ is holomorphic on $D(t,\min(t,1-t))$ (both $z$ and $1-z$ stay in the right
half-plane there), so
\[
\psi(\varsigma):=\frac{\logit(t+\varsigma)-\logit t}{\varsigma}
\]
extends holomorphically to $D(0,\min(t,1-t))$, with real Taylor coefficients at $0$
($\logit$ is real-analytic at the real point $t$); thus
$\logit(t+\varsigma)=\logit t+\varsigma\psi(\varsigma)$ there.  Set
$C_\psi:=\sup_{|\varsigma|\le\varsigma_1}|\psi|<\infty$.

Since $\widehat\cR$ is continuous at $(0,0)$ with $|\widehat\cR(0,0)|=|\vartheta_0|>0$,
shrink $\delta_1$ and $\varsigma_1$ (preserving the constraints above) so that
\[
|\widehat\cR(\zeta,\varsigma)|\ge\tfrac12|\vartheta_0|
\qquad\text{on }\ \bar P:=\overline{D(0,\delta_1)}\times\overline{D(0,2\varsigma_1)} .
\]
Define on a neighbourhood of $\bar P$ the holomorphic, real-symmetric function
\[
F(\zeta,\varsigma):=\frac{c-\widehat h(\zeta)+\varsigma\,\psi(\varsigma)}
{\widehat\cR(\zeta,\varsigma)},
\qquad
M_F:=\sup_{\bar P}|F|<\infty,
\]
and record the Cauchy estimate: for $|\zeta|\le\delta_1$, $|\varsigma|\le\varsigma_1$,
the disc $D(\varsigma,\varsigma_1)$ lies in $D(0,2\varsigma_1)$, so
$|\partial_\varsigma F(\zeta,\varsigma)|\le M_F/\varsigma_1=:L_F$; by integrating along
segments (the disc $\overline{D(0,\varsigma_1)}$ is convex),
$|F(\zeta,a)-F(\zeta,b)|\le L_F|a-b|$ for $a,b\in\overline{D(0,\varsigma_1)}$.

\emph{Step 3 (the fixed point).}  Let
\[
T:=\max\Big\{\tfrac1{\delta_1},\ \tfrac{\rho}{\varsigma_1},\
\tfrac{4C_\psi}{|\vartheta_0|},\ 2L_F\Big\},
\qquad
\rho:=\max\Big(1,\ \tfrac{4}{|\vartheta_0|}\big(|c|+H_\infty\big)\Big),
\]
noting $T\ge1/\delta_1\ge2T_0\ge T_0$.  Let $\widehat X_T$ be the complex Banach space
of bounded holomorphic functions on $A_T:=\{|\tau|>T\}$ with the sup norm (complete:
a sup-norm Cauchy sequence converges uniformly on $A_T$, and uniform limits of
holomorphic functions are holomorphic, by Weierstrass), and let
$\bar B_\rho\subseteq\widehat X_T$ be the closed ball of radius $\rho$, a complete
metric space.  Define
\[
\cT(x)(\tau):=F\Big(\tfrac1\tau,\ \tfrac{x(\tau)}{\tau}\Big),\qquad x\in\bar B_\rho .
\]
This is well defined and holomorphic on $A_T$: $|1/\tau|<1/T\le\delta_1$ and
$|x(\tau)/\tau|\le\rho/T\le\varsigma_1$, so the argument stays in
$\overline{D(0,\delta_1)}\times\overline{D(0,\varsigma_1)}\subset\bar P$.
\emph{Self-map:}
\[
|\cT(x)(\tau)|
\le\frac{|c|+H_\infty+C_\psi\,\rho/T}{\tfrac12|\vartheta_0|}
=\frac{2}{|\vartheta_0|}\big(|c|+H_\infty\big)
+\frac{2C_\psi}{|\vartheta_0|}\cdot\frac{\rho}{T}
\le\frac{\rho}{2}+\frac{\rho}{2}=\rho,
\]
using $\rho\ge\frac4{|\vartheta_0|}(|c|+H_\infty)$ and
$T\ge4C_\psi/|\vartheta_0|$.  \emph{Contraction:} for $x,y\in\bar B_\rho$ and
$\tau\in A_T$, the points $x(\tau)/\tau,\,y(\tau)/\tau$ lie in the convex disc
$\overline{D(0,\varsigma_1)}$, so
\[
|\cT(x)(\tau)-\cT(y)(\tau)|
\le L_F\,\frac{|x(\tau)-y(\tau)|}{|\tau|}
\le\frac{L_F}{T}\,\|x-y\|\le\tfrac12\|x-y\| .
\]
By the Banach fixed point theorem there is a unique $x^*\in\bar B_\rho$ with
$\cT(x^*)=x^*$.  \emph{Reality:} the map $x^\dagger(\tau):=\overline{x^*(\bar\tau)}$
lies in $\bar B_\rho$, and by real symmetry of $F$,
\[
\cT(x^\dagger)(\tau)=F\big(\tfrac1\tau,\tfrac{\overline{x^*(\bar\tau)}}{\tau}\big)
=\overline{F\big(\tfrac1{\bar\tau},\tfrac{x^*(\bar\tau)}{\bar\tau}\big)}
=\overline{\cT(x^*)(\bar\tau)}=\overline{x^*(\bar\tau)}=x^\dagger(\tau),
\]
so $x^\dagger$ is also a fixed point in $\bar B_\rho$; uniqueness gives
$x^\dagger=x^*$, i.e.\ $x^*$ is real-symmetric, in particular real on $(T,\infty)$.

\emph{Step 4 (the gate consistency).}  Put $S(\tau):=t+x^*(\tau)/\tau$ on $A_T$;
then $|S(\tau)-t|\le\rho/T\le\varsigma_1<\min(t,1-t)$, so for real $\tau>T$ the value
$S(\tau)$ is real and lies in $(0,1)$.  Multiplying the fixed-point equation
$x^*=F(1/\tau,x^*/\tau)$ by $\widehat\cR(1/\tau,x^*/\tau)$ and rearranging:
\[
\widehat h\big(\tfrac1\tau\big)+x^*\,\widehat\cR\big(\tfrac1\tau,\tfrac{x^*}{\tau}\big)
=c+\tfrac{x^*}{\tau}\,\psi\big(\tfrac{x^*}{\tau}\big).
\]
The left side equals
$\tau Q(\tau,0)+\tau\big[\widehat Q(\tfrac1\tau,\tfrac{x^*}\tau)
-\widehat Q(\tfrac1\tau,0)\big]=\tau\,Q\big(\tau,\tfrac{x^*(\tau)}{\tau}\big)$
(by $\varsigma\widehat\cR=\widehat Q(\cdot,\varsigma)-\widehat Q(\cdot,0)$ and
$\widehat h(1/\tau)=\tau Q(\tau,0)$); the right side equals
$\logit\big(t+\tfrac{x^*}{\tau}\big)-b=\logit S(\tau)-b$ (the defining expansion of
$\psi$, valid as $|x^*/\tau|\le\varsigma_1$).  Hence
\begin{equation}\label{eq:fixed}
\logit S(\tau)-b=\tau\,Q\Big(\tau,\frac{x^*(\tau)}{\tau}\Big)
\qquad(\text{real }\tau>T).
\end{equation}

\emph{Step 5 (assembling the path).}  Define, on $A_T$,
\[
v(\tau):=v_{S(\tau)}(\tau)=V\Big(\tfrac1\tau,\,S(\tau)\Big),
\qquad
w(\tau):=u(\tau)-v(\tau).
\]
Both are holomorphic (compositions; the curve $(1/\tau,S(\tau))$ stays in
$\overline{D(0,\delta_1)}\times\overline{D(t,\varsigma_1)}$, a compact subset of $V$'s
domain since $\varsigma_1<\varsigma_*$), bounded ($V$ and $u$ are bounded on the
relevant compacts), and real on $(T,\infty)$ ($u$, $x^*$, hence $S$, are real there,
and $V$ is real at real arguments with $\zeta\in(0,\delta_1)$, by the Claim).  Hence
$P:=(w,v)\in\cH_T$, and, $x^*$ being bounded,
$S(\tau)\to t$, so by joint continuity of $V$ at $(0,t)$:
$v(\tau)\to V(0,t)=v^0$ and $w(\tau)\to u^0-v^0=w^0$, i.e.\ $P(\infty)=(w^0,v^0)$.

For real $\tau>T$: by the definitions of $Q$ and of $w,v$,
$Q\big(\tau,x^*(\tau)/\tau\big)=q\big(u(\tau)-v_{S(\tau)}(\tau),
v_{S(\tau)}(\tau)\big)=q\big(w(\tau),v(\tau)\big)$, so \eqref{eq:fixed} reads
$\logit S(\tau)-b=\tau\,q(w(\tau),v(\tau))$, i.e.\
$S(\tau)=\sig\big(\tau q(w(\tau),v(\tau))+b\big)=s_1(\tau)$, using $S(\tau)\in(0,1)$
and $\sig\circ\logit=\mathrm{id}_{(0,1)}$.  Identity \eqref{eq:vSident} at
$S=S(\tau)=s_1(\tau)$ gives $B_1v(\tau)+s_1(\tau)V_1w(\tau)=\widetilde v(\tau)$; and
\[
\big(B_1-s_1(\tau)V_1\big)w(\tau)
=B_1u(\tau)-\big(B_1v(\tau)+s_1(\tau)V_1w(\tau)\big)
=\big(\widetilde w(\tau)+\widetilde v(\tau)\big)-\widetilde v(\tau)
=\widetilde w(\tau),
\]
by the definition of $u$.  This is the assertion, with $T\ge T_0$.

\emph{Closing remarks.}  (i) The hypotheses $\pi^0\neq0$ and $V_1w^0\neq0$ are implied
by $\vartheta_0\neq0$ (if either vanished, so would $\vartheta_0$); they are retained
in the statement for citation convenience.  (ii) Invertibility of $A_1$ is not used.
(iii) $\det B_1\neq0$ ((G3)) is used exactly once, in the definition of $u$.
\end{proof}

\subsection{The sweep}

\begin{lemma}[Sweep]\label{lem:sweep}
Let the first layer be common with $\det B_1\neq0$, let $(w^\sharp,v^\sharp)\in\cQ$
with $\pi(w^\sharp,v^\sharp)\neq0$, and let $t_\sharp\in(0,1)$ satisfy
$\det(B_1-t_\sharp V_1)\neq0$ and $\vartheta_0(w^\sharp,v^\sharp,t_\sharp)\neq0$.
Assume also $(w^\sharp,v^\sharp)\in\cO$ with $\cO\subseteq\R^{2d}$ open.  Then there
are a relatively open set $N\subseteq\cQ\times(0,1)$ containing
$(w^\sharp,v^\sharp,t_\sharp)$ and an open set $\widetilde\cO\subseteq\R^{2d}$
containing $p^\flat:=\Phi_{t_\sharp}(w^\sharp,v^\sharp)$ such that:
\begin{equation}\label{eq:persist}
\text{every }(w,v,t)\in N\text{ satisfies }
(w,v)\in\cO,\ q(w,v)=0,\ \pi(w,v)\neq0,\ V_1w\neq0,\
\det(B_1-tV_1)\neq0,\ \vartheta_0(w,v,t)\neq0;
\end{equation}
and $\widetilde\cO\subseteq\Psi(N)$, where $\Psi(w,v,t):=\Phi_t(w,v)$.
\end{lemma}

\begin{proof}
\emph{Step 0 (persistence).}  Note first that
$\vartheta_0(w^\sharp,v^\sharp,t_\sharp)\neq0$ forces $V_1w^\sharp\neq0$ and
$\pi(w^\sharp,v^\sharp)\neq0$.  The functions
$(w,v)\mapsto\pi(w,v)$, $(w,v)\mapsto V_1w$, $(w,v,t)\mapsto\det(B_1-tV_1)$ are
continuous, and $\vartheta_0$ is continuous on the open set
$\{\det(B_1-tV_1)\neq0\}$ (a ratio of polynomials via the adjugate); each of the
finitely many quantities in \eqref{eq:persist} is nonzero, or the membership open, at
$(w^\sharp,v^\sharp,t_\sharp)$.  Hence there is a relatively open
$N_0\subseteq\cQ\times(0,1)$ containing the point on which \eqref{eq:persist} holds
(the equality $q=0$ holds on all of $\cQ\times(0,1)$).

\emph{Step 1 (chart on the quadric).}  $\nabla q(w^\sharp,v^\sharp)
=(A_1v^\sharp,A_1^\top w^\sharp)\neq0$: if both components vanished, then
$\pi(w^\sharp,v^\sharp)=A_1^\top w^\sharp-A_1v^\sharp=0$, contrary to Step 0.  By the
implicit function theorem applied to the polynomial $q$, there are an open
$\Xi\subseteq\R^{2d-1}$, a point $\xi^\sharp\in\Xi$ and a real-analytic map
$\gamma:\Xi\to\R^{2d}$ which is a homeomorphism onto a relatively open neighbourhood of
$(w^\sharp,v^\sharp)$ in $\cQ$, with $\gamma(\xi^\sharp)=(w^\sharp,v^\sharp)$ and
$D\gamma(\xi^\sharp)$ of rank $2d-1$ with image the tangent space
\[
\cT:=\ker\nabla q(w^\sharp,v^\sharp)^\top
=\big\{(\delta w,\delta v):\ (A_1v^\sharp)^\top\delta w
+(A_1^\top w^\sharp)^\top\delta v=0\big\}.
\]

\emph{Step 2 (the Jacobian and the inverse function theorem).}  Define the
real-analytic map
\[
G:\Xi\times(0,1)\to\R^{2d},\qquad G(\xi,t):=\Phi_t\big(\gamma(\xi)\big).
\]
Its differential at $(\xi^\sharp,t_\sharp)$ consists of the $2d-1$ columns
$\Phi_{t_\sharp}\big(D\gamma(\xi^\sharp)\,e_i\big)$ ($\Phi_{t_\sharp}$ is linear in
$(w,v)$) together with the $t$-column
$\partial_t\Phi_t(w^\sharp,v^\sharp)\big|_{t_\sharp}=(-V_1w^\sharp,\,V_1w^\sharp)$.
The linear map $\Phi_{t_\sharp}$ has block-triangular matrix
$\big[\begin{smallmatrix}B_1-t_\sharp V_1&0\\ t_\sharp V_1&B_1\end{smallmatrix}\big]$,
hence determinant $\det(B_1-t_\sharp V_1)\det B_1\neq0$: it is an isomorphism of
$\R^{2d}$.  The first $2d-1$ columns therefore span the $(2d-1)$-dimensional subspace
$\Phi_{t_\sharp}(\cT)$, and $DG(\xi^\sharp,t_\sharp)$ is invertible if and only if
\[
h:=\Phi_{t_\sharp}^{-1}\big(-V_1w^\sharp,\ V_1w^\sharp\big)\notin\cT .
\]
Compute $h=(\delta w,\delta v)$: $(B_1-t_\sharp V_1)\delta w=-V_1w^\sharp$ gives
$\delta w=-m$ with $m:=(B_1-t_\sharp V_1)^{-1}V_1w^\sharp$; and
$B_1\delta v+t_\sharp V_1\delta w=V_1w^\sharp$ gives
\[
\delta v=B_1^{-1}\big(I+t_\sharp V_1(B_1-t_\sharp V_1)^{-1}\big)V_1w^\sharp
=B_1^{-1}\,B_1(B_1-t_\sharp V_1)^{-1}V_1w^\sharp=m,
\]
using $I+t_\sharp V_1(B_1-t_\sharp V_1)^{-1}
=(B_1-t_\sharp V_1+t_\sharp V_1)(B_1-t_\sharp V_1)^{-1}=B_1(B_1-t_\sharp V_1)^{-1}$.
Hence $h=(-m,m)$ and
\[
\nabla q(w^\sharp,v^\sharp)\cdot h
=-(A_1v^\sharp)^\top m+(A_1^\top w^\sharp)^\top m
=m^\top\big(A_1^\top w^\sharp-A_1v^\sharp\big)
=\pi(w^\sharp,v^\sharp)^\top(B_1-t_\sharp V_1)^{-1}V_1w^\sharp
=\vartheta_0(w^\sharp,v^\sharp,t_\sharp)\neq0,
\]
so $h\notin\cT$ and $DG(\xi^\sharp,t_\sharp)$ is invertible.  By the inverse function
theorem, after shrinking $\Xi$ to $\Xi'\ni\xi^\sharp$ and $(0,1)$ to an open interval
$J'\ni t_\sharp$, the map $G$ is a diffeomorphism of $\Xi'\times J'$ onto an open set
$\widetilde\cO\subseteq\R^{2d}$ containing $G(\xi^\sharp,t_\sharp)=p^\flat$; shrinking
further, we may assume $N:=\gamma(\Xi')\times J'\subseteq N_0$ (so \eqref{eq:persist}
holds on $N$; $\gamma(\Xi')$ is relatively open in $\cQ$, so $N$ is relatively open in
$\cQ\times(0,1)$).  Every point of $\widetilde\cO$ is of the form
$G(\xi,t)=\Psi(\gamma(\xi),t)$ with $(\gamma(\xi),t)\in N$, i.e.\
$\widetilde\cO\subseteq\Psi(N)$.
\end{proof}

\subsection{Proof of Theorem~\ref{thm:IDL}}

\begin{proof}[Proof of Theorem~\ref{thm:IDL}]
Induction on $L$.

\emph{Base case $L=1$.}  Here $\theta'=(V_1',A_1')\in\cG^{(1)}$, so $V_1'\neq0$ and
$A_1'\neq0$; $\theta=(V_1,A_1)$ is arbitrary; $\cO$ is open and nonempty, with no
anchor condition.  By Proposition~\ref{prop:probe}, for any $\vartheta=(\widetilde
V,\widetilde A)$ and any path,
\begin{equation}\label{eq:onelayer}
G^P_\vartheta(\tau)=\widetilde B\,v(\tau)
+\sig\big(\tau\,w(\tau)^\top\widetilde A\,v(\tau)+b\big)\,\widetilde V\,w(\tau),
\qquad \widetilde B:=I+\widetilde V .
\end{equation}

Choose an entry $(A_1')_{jk}\neq0$ and a row $\iota$ with $e_\iota^\top V_1'\neq0$, and
set
\[
\cO_g:=\{(w,v)\in\cO:\ w^\top A_1'v\neq0,\ e_\iota^\top V_1'w\neq0\}.
\]
The two polynomials are nonzero (the first takes the value $(A_1')_{jk}\neq0$ at
$(e_j,e_k)$; the second is a nonzero linear form in $w$), so $\cO_g$ is open and
nonempty (Lemma~\ref{lem:zariski}).  Fix $(w,v)\in\cO_g$ and write
$\lambda:=w^\top A_1v$, $\lambda':=w^\top A_1'v\neq0$.  The constant path at $(w,v)$
lies in $\cP(\cO)$, so by hypothesis and \eqref{eq:onelayer}, for $\tau>T:=T_{(w,v)}$,
\begin{equation}\label{eq:onelayeragree}
B_1v+\sig(\tau\lambda+b)\,V_1w
=B_1'v+\sig(\tau\lambda'+b)\,V_1'w .
\end{equation}

\emph{$\lambda=\lambda'$.}  Take the $\iota$-th coordinate of
\eqref{eq:onelayeragree}.  Both sides are restrictions to $(T,\infty)$ of functions
meromorphic on $\C$, holomorphic off the pole sets
$E_F\subseteq\{(-b+(2n+1)\pi i)/\lambda\}_{n\in\Z}$ (empty if $\lambda=0$ or
$e_\iota^\top V_1w=0$) and $E_G=\{\tau_k:=((2k+1)\pi i-b)/\lambda'\}_{k\in\Z}$
respectively --- Lemma~\ref{lem:sigmoid} --- and the pole sets are countable, closed,
disjoint from $\R$ (nonzero imaginary parts), and consist of isolated points.  At each
$\tau_k$ the right side blows up: $|\sig(\tau\lambda'+b)|\to\infty$
(Lemma~\ref{lem:sigmoid}: simple poles) while its coefficient
$e_\iota^\top V_1'w\neq0$ and the constant term is bounded.  Lemma~\ref{lem:transfer}
gives $\tau_k\in E_F$ for every $k$; in particular $E_F\neq\varnothing$, forcing
$\lambda\neq0$ and $e_\iota^\top V_1w\neq0$, and $E_G\subseteq E_F$ with every point of
$E_F$ of real part $-b/\lambda$ and every point of $E_G$ of real part $-b/\lambda'$.
Comparing real parts at the single point $\tau_0$:
$-b/\lambda'=-b/\lambda$, and $b=\log r>0$ yields $\lambda=\lambda'$.  Thus
$w^\top(A_1-A_1')v=0$ on the nonempty open set $\cO_g$; by Lemma~\ref{lem:zariski} the
bilinear polynomial vanishes identically, and testing on $(e_a,e_b)$: $A_1=A_1'$.

\emph{$V_1=V_1'$ and $B_1=B_1'$.}  With $A_1=A_1'$, \eqref{eq:onelayeragree} becomes,
for $(w,v)\in\cO_g$ and $\tau>T$,
\[
(B_1-B_1')\,v=\sig(\tau\lambda'+b)\,(V_1'-V_1)\,w .
\]
Since $\lambda'\neq0$ and $\sig'(x)=\sig(x)(1-\sig(x))>0$ on $\R$, the map
$\tau\mapsto\sig(\tau\lambda'+b)$ is strictly monotone on $\R$; pick
$\tau_1\neq\tau_2>T$, so $\sig(\tau_1\lambda'+b)\neq\sig(\tau_2\lambda'+b)$, and
subtract the two instances of the display:
$0=\big[\sig(\tau_1\lambda'+b)-\sig(\tau_2\lambda'+b)\big](V_1'-V_1)w$, whence
$(V_1'-V_1)w=0$ and then $(B_1-B_1')v=0$.  Both polynomial maps vanish on the nonempty
open $\cO_g$, hence identically (Lemma~\ref{lem:zariski}, componentwise):
$V_1=V_1'$ and $B_1=B_1'$, i.e.\ $\theta=\theta'$.

\emph{Inductive step $L\ge2$.}  Assume the theorem at depth $L-1$.  The hypothesis of
the theorem at depth $L$ is verbatim the standing hypothesis \eqref{eq:SH} of
\S\S\ref{sec:A1}--\ref{sec:cascade}, whose results are proved independently of the
present theorem.

(1)  By Proposition~\ref{prop:A1}, Corollary~\ref{cor:depth} and
Proposition~\ref{prop:matching}: $A_1=A_1'$, $V_1=V_1'\,(\neq0)$, $B_1=B_1'$; by (G3),
$\det B_1\neq0$.  The maps $q,\pi,\Phi_t,\vartheta_0$ coincide for both parameters.

(2)  By the anchor clause of \eqref{eq:SH}, fix
$(w^\sharp,v^\sharp)\in\cO\cap\cM_L(\theta')$, and by (M3) of
Definition~\ref{def:anchor} (formulated through $\theta'$'s first layer, i.e.\ the
common one) a value $t_\sharp\in(0,1)$ with $\det(B_1-t_\sharp V_1)\neq0$,
$\vartheta_0(w^\sharp,v^\sharp,t_\sharp)\neq0$ and
$p^\flat:=\Phi_{t_\sharp}(w^\sharp,v^\sharp)\in\cM_{L-1}(\theta'_{\ge2})$.  By (M1),
$q(w^\sharp,v^\sharp)=0$ and $\pi(w^\sharp,v^\sharp)\neq0$.

(3)  Lemma~\ref{lem:sweep} provides a relatively open $N\subseteq\cQ\times(0,1)$ with
the persistence properties \eqref{eq:persist} (in particular $N\subseteq\{(w,v)\in\cO\}$)
and an open $\widetilde\cO\ni p^\flat$ with $\widetilde\cO\subseteq\Psi(N)$.  Note that
\[
\widetilde\cO\cap\cM_{L-1}(\theta'_{\ge2})\ni p^\flat\neq\varnothing :
\]
the anchor clause at depth $L-1$ is re-established.

(4)  Let $\widetilde P\in\cP(\widetilde\cO)$ be arbitrary, say
$\widetilde P\in\cH_{T_0}$ with $\widetilde p:=\widetilde P(\infty)\in\widetilde\cO$.
By (3), $\widetilde p=\Phi_t(w^0,v^0)$ for some $(w^0,v^0,t)\in N$, and by
\eqref{eq:persist} the hypotheses of Lemma~\ref{lem:realization} hold at
$(w^0,v^0,t)$.  Lemma~\ref{lem:realization} produces $T\ge T_0$ and $P\in\cH_T$ with
$P(\infty)=(w^0,v^0)\in\cO$ --- so $P\in\cP(\cO)$ --- whose first-layer effective path
is $\widetilde P$ on $(T,\infty)$.  Since the first layers of $\theta$ and $\theta'$
coincide, the first gate $s_1(\tau)$ along $P$ is the same for both parameters, and
Lemma~\ref{lem:peeling} gives, for real $\tau>T$,
\[
G^P_\theta(\tau)=F^{(L-1)}_{\theta_{\ge2}}
\big(\widetilde w(\tau),\widetilde v(\tau),\tau\big)
=G^{\widetilde P}_{\theta_{\ge2}}(\tau),
\qquad
G^P_{\theta'}(\tau)=G^{\widetilde P}_{\theta'_{\ge2}}(\tau).
\]
By \eqref{eq:SH} there is $T_P$ with $G^P_\theta=G^P_{\theta'}$ on $(T_P,\infty)$;
hence, with $T_{\widetilde P}:=\max(T,T_P,0)$,
\[
G^{\widetilde P}_{\theta_{\ge2}}(\tau)=G^{\widetilde P}_{\theta'_{\ge2}}(\tau)
\qquad\text{for all real }\tau>T_{\widetilde P}.
\]

(5)  The data $\big(\theta'_{\ge2},\theta_{\ge2},\widetilde\cO\big)$ satisfy the
hypotheses of the theorem at depth $L-1$: $\theta'_{\ge2}\in\cG^{(L-1)}$
(Definition~\ref{def:generic}), $d\ge2$ and $r$ are unchanged, $\widetilde\cO$ is open
and nonempty, $\widetilde\cO\cap\cM_{L-1}(\theta'_{\ge2})\neq\varnothing$ by (3) (a
condition that is anyway vacuous when $L-1=1$), and the agreement along every
$\widetilde P\in\cP(\widetilde\cO)$ was established in (4).  By the inductive
hypothesis, $\theta_{\ge2}=\theta'_{\ge2}$; together with (1), $\theta=\theta'$.  This
completes the induction and the proof.
\end{proof}

\section{Genericity: certificates, anchors, and witnesses}\label{sec:generic}

\emph{Standing notation.}  Throughout, $\theta'\in\cP=\cP^{(L)}$, with tails
$\theta'_{\ge k+1}$ (layer $i$ of $\theta'_{\ge k+1}$ is layer $k+i$ of $\theta'$);
the primed formal streams $w_k(z),v_k(z)$ and matrices $W'_k(z)$, $T'_k(z)$ are those
of \S\ref{subsec:slope} for $\theta'$, with stream index up to $L-1$.  We write
$\adj(M)$ for the adjugate: $M\adj(M)=\det(M)\,I$, $\adj(M)=\det(M)M^{-1}$ for
invertible $M$, and every entry of $\adj(M)$ is a polynomial in the entries of $M$.
(The certificate rows $T_k(v)$ below should not be confused with the stream tails
$T'_k(z)$ or the sequence length $T=r+1$; the arguments disambiguate.)  This section
uses only \S\ref{subsec:slope} and Lemma~\ref{lem:zariski}; it is independent of
\S\S\ref{sec:A1}--\ref{sec:descent}.

\begin{auxlemma}[Composition identity]\label{lem:comp}
Let $\theta'$ be an $L$-tuple, $(w,v)\in\R^{2d}$, $0\le k\le L-1$ and
$t_1,\dots,t_k\in\R$.  Put
$p_k:=\big(w_{(k)},v_{(k)}\big):=\big(w_k(z),v_k(z)\big)\big|_{z=(t_1,\dots,t_k)}$
(streams of $\theta'$; $p_0=(w,v)$).  Then:
\begin{enumerate}
\item[(i)] for $k\le L-2$ and any $t_{k+1}$,
$p_{k+1}=\Phi^{\theta'_{\ge k+1}}_{t_{k+1}}(p_k)$, where $\Phi^\vartheta_t$ is the map
of Definition~\ref{def:anchor} for the tuple $\vartheta$;
\item[(ii)] the formal streams of the tail $\theta'_{\ge k+1}$ started at $p_k$, in
variables $\zeta=(\zeta_1,\zeta_2,\dots)$, satisfy
$\big(w^{\mathrm{tail}}_j(\zeta),v^{\mathrm{tail}}_j(\zeta)\big)
=\big(w_{k+j},v_{k+j}\big)(t_1,\dots,t_k,\zeta_1,\dots,\zeta_j)$ for
$0\le j\le L-k$;
\item[(iii)] consequently, for $1\le\ell\le L-k$,
\[
\phi^{\theta'_{\ge k+1}}_\ell\big(\zeta_1,\dots,\zeta_{\ell-1};p_k\big)
=\phi'_{k+\ell}\big(t_1,\dots,t_k,\zeta_1,\dots,\zeta_{\ell-1};w,v\big);
\]
in particular $q_{\theta'_{\ge k+1}}(p_k)=\phi'_{k+1}(t_1,\dots,t_k;w,v)$.
\end{enumerate}
\end{auxlemma}

\begin{proof}
(ii)  Induction on $j$.  The base $j=0$ is the definition of $p_k$.  For the step,
layer $j+1$ of the tail is layer $k+j+1$ of $\theta'$, so both recursions apply the
\emph{same} affine map
$(w,v)\mapsto\big((B'_{k+j+1}-z\,V'_{k+j+1})w,\ B'_{k+j+1}v+z\,V'_{k+j+1}w\big)$ with
$z=\zeta_{j+1}$, the tail to its $j$-th streams and $\theta'$ to its $(k+j)$-th streams
specialized at $(t_1,\dots,t_k,\zeta_1,\dots,\zeta_j)$, which agree by the inductive
hypothesis.

(i)  is the case $j=1$, $\zeta_1=t_{k+1}$ of (ii), by the definition of
$\Phi^\vartheta_t$ (the first-layer recursion step of the tuple $\vartheta$).

(iii)  By the definition of the slope polynomials and (ii) at $j=\ell-1$:
$\phi^{\mathrm{tail}}_\ell(\zeta;p_k)
=w^{\mathrm{tail}}_{\ell-1}(\zeta)^\top A'_{k+\ell}\,v^{\mathrm{tail}}_{\ell-1}(\zeta)
=w_{k+\ell-1}^\top A'_{k+\ell}\,v_{k+\ell-1}
=\phi'_{k+\ell}(t,\zeta;w,v)$, the middle streams specialized as stated.
\end{proof}

\subsection{Unwinding the anchor set}

\begin{lemma}[Unwinding]\label{lem:unwind}
Let $L\ge2$, $\theta'\in\cP$ and $(w,v)\in\R^{2d}$.  Then $(w,v)\in\cM_L(\theta')$ if
and only if there exist $t_1,\dots,t_{L-1}\in(0,1)$ and
$s_0,\dots,s_{L-2}\in(0,1)$ such that, with
$p_k=(w_{(k)},v_{(k)})$ as in Auxiliary Lemma~\ref{lem:comp} and
$\pi^{(k)}:=(A'_{k+1})^\top w_{(k)}-A'_{k+1}v_{(k)}$, the following hold for every
$0\le k\le L-2$:
\begin{enumerate}
\item[(a)] $\phi'_{k+1}(t_1,\dots,t_k;w,v)=0$;
\item[(b)] $(A'_{k+1})^\top w_{(k)}\neq0$ and $\pi^{(k)}\neq0$;
\item[(c)] $\phi'_{k+\ell}(t_1,\dots,t_k,s_k,1,\dots,1;w,v)>0$ for
$\ell=2,\dots,L-k$;
\item[(d)] $\det\big(B'_{k+1}-t_{k+1}V'_{k+1}\big)\neq0$ and
$\big(\pi^{(k)}\big)^\top\big(B'_{k+1}-t_{k+1}V'_{k+1}\big)^{-1}V'_{k+1}w_{(k)}\neq0$.
\end{enumerate}
\end{lemma}

\begin{proof}
Induction on $L\ge2$, over all parameters.

\emph{Translation at stage $0$.}  By Auxiliary Lemma~\ref{lem:comp}(iii) (with $k=0$,
where the conditions are statements about $\theta'$ itself), the conditions
(M1)--(M3) of Definition~\ref{def:anchor} for $\cM_L(\theta')$ at $(w,v)$ read:
(M1) $\Leftrightarrow$ (a)$_{k=0}$ and (b)$_{k=0}$; (M2) $\Leftrightarrow$ the
existence of $s_0\in(0,1)$ with (c)$_{k=0}$; and (M3) $\Leftrightarrow$ the existence
of $t_1\in(0,1)$ with (d)$_{k=0}$ and
$p_1=\Phi^{\theta'}_{t_1}(w,v)\in\cM_{L-1}(\theta'_{\ge2})$ --- the last identification
of the transported point by Auxiliary Lemma~\ref{lem:comp}(i).

\emph{Base $L=2$.}  $\cM_1(\theta'_{\ge2})=\R^{2d}$, so the clause
$p_1\in\cM_1$ is vacuous and the translation above is exactly the statement (the index
$k$ takes only the value $0$).

\emph{Step $L\ge3$.}  By the inductive hypothesis applied to the $(L-1)$-tuple
$\theta'_{\ge2}$ at the point $p_1$: $p_1\in\cM_{L-1}(\theta'_{\ge2})$ iff there exist
$t_2,\dots,t_{L-1},\,s_1,\dots,s_{L-2}\in(0,1)$ such that conditions (a)--(d) hold for
the tail at the iterated points
$p^{\mathrm{tail}}_{k'}$ ($0\le k'\le L-3$) built from $p_1$ with gates
$t_2,\dots,t_{1+k'}$.  By Auxiliary Lemma~\ref{lem:comp}(ii) (applied with $k=1$),
$p^{\mathrm{tail}}_{k'}=p_{k'+1}$, and by (iii) every slope polynomial of the tail (and
of the tails of the tail, since
$(\theta'_{\ge2})_{\ge k'+1}=\theta'_{\ge k'+2}$) at these points equals the
corresponding slope polynomial of $\theta'$ at $(w,v)$ with the concatenated gate
string; likewise the data of (b) and (d) for the tail at stage $k'$ are literally the
data of (b) and (d) for $\theta'$ at stage $k=k'+1$ (same matrices
$A'_{k+1},B'_{k+1},V'_{k+1}$, same points $p_k$).  Hence the tail conditions at stages
$0\le k'\le L-3$ are exactly the conditions (a)--(d) at stages $1\le k\le L-2$, and
combining with the stage-$0$ translation proves the equivalence.
\end{proof}

\subsection{The anchor certificate}

\begin{definition}[Anchor certificate]\label{def:certificate}
Let $\theta'\in\cP$, $t=(t_1,\dots,t_{L-1})\in\R^{L-1}$,
$s=(s_0,\dots,s_{L-2})\in\R^{L-1}$ and $w^\circ\in\R^d$.  With
$w_{(k)}:=W'_k(t_1,\dots,t_k)\,w^\circ$ and
$\pi^{(k)}(v):=(A'_{k+1})^\top w_{(k)}
-A'_{k+1}\big(B'_{k:1}v+T'_k(t_1,\dots,t_k)w^\circ\big)$, define the affine-in-$v$
\emph{certificate rows}
\begin{align*}
E_j(v)&:=\phi'_j(t_1,\dots,t_{j-1};w^\circ,v), &&1\le j\le L-1,\\
S_{k,\ell}(v)&:=\phi'_{k+\ell}(t_1,\dots,t_k,s_k,1,\dots,1;w^\circ,v),
&&0\le k\le L-2,\ \ 2\le\ell\le L-k,\\
T_k(v)&:=\pi^{(k)}(v)^\top\,\adj\!\big(B'_{k+1}-t_{k+1}V'_{k+1}\big)\,
V'_{k+1}\,w_{(k)},&&0\le k\le L-2 .
\end{align*}
Their number is
$N=(L-1)+\binom L2+(L-1)=\binom L2+2(L-1)$ (there are $L-1$ rows $E_j$,
$\sum_{k=0}^{L-2}(L-k-1)=\binom L2$ rows $S_{k,\ell}$, and $L-1$ rows $T_k$).  Let
$\mathcal G=\mathcal G(\theta';t,s,w^\circ)\in\R^{d\times N}$ be the matrix whose
columns are the $v$-gradients of the $N$ rows, in a fixed order, and set
\[
\mathfrak c(\theta';t,s,w^\circ)
:=\det\!\big(\mathcal G^\top\mathcal G\big)\cdot
\prod_{k=0}^{L-2}\det\!\big(B'_{k+1}-t_{k+1}V'_{k+1}\big)\cdot
\prod_{k=0}^{L-2}\big|(A'_{k+1})^\top w_{(k)}\big|^2 .
\]
We say $\theta'$ \emph{satisfies the anchor certificate} --- this is condition
\textup{(G4)} of Definition~\ref{def:generic} --- if
$\mathfrak c(\theta';\cdot,\cdot,\cdot)\not\equiv0$ as a polynomial in
$(t,s,w^\circ)\in\R^{L-1}\times\R^{L-1}\times\R^d$.
\end{definition}

By Auxiliary Lemma~\ref{lem:grad} below, $\mathfrak c$ is a polynomial in
$(\theta',t,s,w^\circ)$ jointly; hence (G4) is the condition that the finitely many
coefficients of $\mathfrak c$ as a polynomial in $(t,s,w^\circ)$ --- themselves
polynomials in $\theta'$ --- do not all vanish, as asserted after
Definition~\ref{def:generic}.

\begin{auxlemma}[Gradients]\label{lem:grad}
Each certificate row is an affine function of $v$, with constant $v$-gradients
\begin{align*}
\nabla_vE_j&=(B'_{j-1:1})^\top(A'_j)^\top\,W'_{j-1}(t_1,\dots,t_{j-1})\,w^\circ,\\
\nabla_vS_{k,\ell}&=(B'_{k+\ell-1:1})^\top(A'_{k+\ell})^\top\,
\big(B'_{k+1}-s_kV'_{k+1}\big)\,W'_k(t_1,\dots,t_k)\,w^\circ,\\
\nabla_vT_k&=-\,(B'_{k:1})^\top(A'_{k+1})^\top\,
\adj\!\big(B'_{k+1}-t_{k+1}V'_{k+1}\big)\,V'_{k+1}\,W'_k(t_1,\dots,t_k)\,w^\circ .
\end{align*}
\textup{(}For $j=1$ the products are empty and $\nabla_vE_1=(A'_1)^\top w^\circ$;
for $\ell=2$ the trailing string $1,\dots,1$ in $S_{k,2}$ is empty.\textup{)}  Every
entry of every gradient, and every other factor of $\mathfrak c$, is a polynomial in
$(\theta',t,s,w^\circ)$ jointly.
\end{auxlemma}

\begin{proof}
\emph{$E_j$.}  By Lemma~\ref{lem:phi}(i) applied to $\theta'$ with $\ell=j$ and
$z=(t_1,\dots,t_{j-1})$, the map $v\mapsto\phi'_j(z;w^\circ,v)$ is affine with
gradient $(B'_{j-1:1})^\top(A'_j)^\top W'_{j-1}(z)w^\circ$.

\emph{$S_{k,\ell}$.}  The same formula with $\ell$ replaced by $k+\ell$ and
$z=(t_1,\dots,t_k,s_k,1,\dots,1)$ gives the gradient
$(B'_{k+\ell-1:1})^\top(A'_{k+\ell})^\top W'_{k+\ell-1}(z)w^\circ$.  By transparency
(Lemma~\ref{lem:phi}(iv), or directly from \eqref{eq:transparency}), the factors
of $W'_{k+\ell-1}(z)$ carrying the trailing ones collapse to $I$:
\[
W'_{k+\ell-1}(t_1,\dots,t_k,s_k,1,\dots,1)
=\big(B'_{k+1}-s_kV'_{k+1}\big)\,W'_k(t_1,\dots,t_k),
\]
yielding the stated form.

\emph{$T_k$.}  Only $\pi^{(k)}(v)$ depends on $v$, affinely, with Jacobian of
$v\mapsto B'_{k:1}v+T'_k(t)w^\circ$ equal to $B'_{k:1}$; hence
$v\mapsto T_k(v)$ is affine with gradient
$\big(-A'_{k+1}B'_{k:1}\big)^\top\adj(B'_{k+1}-t_{k+1}V'_{k+1})V'_{k+1}w_{(k)}$, which
is the stated expression.

\emph{Polynomiality.}  The entries of $W'_k(z)$ and $T'_k(z)$ are polynomials in
$(\theta',z)$ (\S\ref{subsec:slope}); the entries of the adjugate are polynomials in
the entries of its argument; determinants and squared Euclidean norms are polynomials.
Composing, every quantity above, and hence $\mathfrak c$, is a polynomial in
$(\theta',t,s,w^\circ)$.
\end{proof}

\subsection{Anchors exist generically}

\begin{proposition}[Anchors exist]\label{prop:anchorsexist}
Let $L\ge2$, $d\ge N$, and let $\theta'\in\cP$ satisfy \textup{(G4)} (in particular,
any $\theta'\in\cG^{(L)}$).  Then $\cM_L(\theta')\neq\varnothing$.
\end{proposition}

\begin{proof}
(F1) \emph{A good evaluation point.}  By (G4), $\mathfrak c(\theta';\cdot)\not\equiv0$
as a polynomial in $(t,s,w^\circ)$.  The box
$(0,1)^{L-1}\times(0,1)^{L-1}\times\R^d$ is a nonempty open subset of
$\R^{L-1}\times\R^{L-1}\times\R^d$, so by Lemma~\ref{lem:zariski} there is a point
$(t,s,w^\circ)$ in it with $\mathfrak c(\theta';t,s,w^\circ)\neq0$.  Unpacking the
factors of $\mathfrak c$: the $N$ gradient columns of $\mathcal G$ are linearly
independent ($\det\mathcal G^\top\mathcal G\neq0$ is the Gram criterion);
$\det(B'_{k+1}-t_{k+1}V'_{k+1})\neq0$ and $(A'_{k+1})^\top w_{(k)}\neq0$ for
$0\le k\le L-2$; and $t_k,s_k\in(0,1)$.

(F2) \emph{The linear solve.}  Consider the affine map
$\R^d\to\R^N$, $v\mapsto\big(E_j(v);\,S_{k,\ell}(v);\,T_k(v)\big)$.  Its linear part
is $v\mapsto\mathcal G^\top v$, of rank $N$ (independent columns; here $d\ge N$ is
used), hence surjective; so the affine map is surjective.  Choose $v\in\R^d$ with
\[
E_j(v)=0\ \ (1\le j\le L-1),\qquad
S_{k,\ell}(v)=1\ \ (\text{all }k,\ell),\qquad
T_k(v)=1\ \ (0\le k\le L-2).
\]

(F3) \emph{Verification via the unwinding.}  We check the conditions of
Lemma~\ref{lem:unwind} for $(w^\circ,v)$ with the gate values $t$ and dials $s$ chosen
in (F1).  Note that the points $p_k$ and covectors $\pi^{(k)}$ of
Lemma~\ref{lem:unwind} are exactly those of Definition~\ref{def:certificate} (same
streams, same specializations).  For $0\le k\le L-2$:
(a) $\phi'_{k+1}(t_{\le k};w^\circ,v)=E_{k+1}(v)=0$;
(b) $(A'_{k+1})^\top w_{(k)}\neq0$ by (F1), and $\pi^{(k)}\neq0$ because
$T_k(v)=1\neq0$ and $T_k(v)$ is a linear functional of $\pi^{(k)}$;
(c) $\phi'_{k+\ell}(t_{\le k},s_k,1,\dots,1;w^\circ,v)=S_{k,\ell}(v)=1>0$ for
$\ell=2,\dots,L-k$, with the single dial $s_k\in(0,1)$;
(d) $\det(B'_{k+1}-t_{k+1}V'_{k+1})\neq0$ by (F1), and, by
$\adj(M)=\det(M)M^{-1}$,
\[
\big(\pi^{(k)}\big)^\top\big(B'_{k+1}-t_{k+1}V'_{k+1}\big)^{-1}V'_{k+1}w_{(k)}
=\frac{T_k(v)}{\det\big(B'_{k+1}-t_{k+1}V'_{k+1}\big)}\neq0 .
\]
By Lemma~\ref{lem:unwind}, $(w^\circ,v)\in\cM_L(\theta')$.
\end{proof}

\subsection{The genericity set is generic}

\begin{proposition}[Genericity]\label{prop:genericnonempty}
Let $L\ge1$ and $d\ge\max(2,N_L)$, where $N_m:=\binom m2+2(m-1)$.  Then
$\cP^{(L)}\setminus\cG^{(L)}$ is contained in a finite union of proper algebraic
subsets of $\cP^{(L)}$; in particular $\cG^{(L)}$ contains the complement of such a
union.
\end{proposition}

\begin{proof}
Since $N_m$ is nondecreasing in $m$, the hypothesis implies $d\ge\max(2,N_m)$ for
every $1\le m\le L$.  Induct on $L$.

\emph{Base $L=1$.}  $\cP^{(1)}\setminus\cG^{(1)}=\{V_1'=0\}\cup\{A_1'=0\}$, a union of
two proper linear subspaces (nonzero matrices exist), each algebraic.

\emph{Step $L\ge2$.}  By Definition~\ref{def:generic},
\[
\cP\setminus\cG^{(L)}\ \subseteq\
Z_1\cup Z_2\cup Z_3\cup Z_4\ \cup\ p^{-1}\big(\cP^{(L-1)}\setminus\cG^{(L-1)}\big),
\]
where $p(\theta'):=\theta'_{\ge2}$ is the (linear, surjective) projection and:

$Z_1:=\{\det A_1'=0\}\cup\{\Sym A_1'=0\}$: both constituents are algebraic ($\det$ is
a polynomial; $\Sym A_1'=0$ is a system of linear equations) and proper --- witness
any $\theta'$ with $A_1'=I$ ($\det=1$, $\Sym=I$).

$Z_2:=\bigcup_{j=2}^L\{\Sym(V_{j-1:1}'^{\,\top}A_j'V_{j-1:1}')=0\}\cup\{V'_{L:1}=0\}$:
each constituent is algebraic (entries are polynomials in $\theta'$) and proper ---
witness $V'_\ell=e_1e_1^\top$, $A'_\ell=I$ for all $\ell$: then
$V'_{j-1:1}=(e_1e_1^\top)^{j-1}=e_1e_1^\top$ (idempotent), so
$\Sym(V_{j-1:1}'^{\,\top}A_j'V_{j-1:1}')=\Sym(e_1e_1^\top)=e_1e_1^\top\neq0$ and
$V'_{L:1}=e_1e_1^\top\neq0$.

$Z_3:=\{\det B_1'=0\}$: algebraic; proper --- witness $V_1'=e_1e_1^\top$,
$\det(I+e_1e_1^\top)=2\neq0$.

$Z_4:=\{\theta':\ \mathfrak c(\theta';\cdot,\cdot,\cdot)\equiv0\}$: the common zero
set of the finitely many coefficients of $\mathfrak c$ in $(t,s,w^\circ)$, each a
polynomial in $\theta'$ (Auxiliary Lemma~\ref{lem:grad}) --- algebraic; proper by
Lemma~\ref{lem:witness} below, whose proof is independent of the present proposition.

\emph{Recursive part.}  By the inductive hypothesis,
$\cP^{(L-1)}\setminus\cG^{(L-1)}\subseteq Y_1\cup\dots\cup Y_M$ with each $Y_i$ a
proper algebraic subset of $\cP^{(L-1)}$.  Each pullback $p^{-1}(Y_i)$ is algebraic in
$\cP$ (compositions of the defining polynomials with the projection) and proper:
choose $\eta\in\cP^{(L-1)}\setminus Y_i$ and any first layer; the resulting parameter
lies outside $p^{-1}(Y_i)$.

Altogether $\cP\setminus\cG^{(L)}$ is contained in a finite union of proper algebraic
subsets.  \emph{For later use in the proof of Corollary~\ref{cor:factors} we record
that every constituent listed above admits a witness all of whose $V$-matrices have
rank at most one:} for $Z_1,Z_2,Z_3$ the witness $V'_\ell=e_1e_1^\top$,
$A'_\ell=I$ just verified (it avoids each of these constituents simultaneously); for
$Z_4$ the witness $\theta^\circ$ of Lemma~\ref{lem:witness}, whose $V$-matrices are
rank one by construction; and, for the recursive constituents $p^{-1}(Y_i)$, by
induction a depth-$(L-1)$ witness with rank-$\le1$ $V$-matrices, prepended by the
first layer $(e_1e_1^\top,\,I)$.
\end{proof}

\begin{lemma}[Witness for the certificate]\label{lem:witness}
Let $L\ge2$ and $d\ge\max(2,N)$.  There is $\theta^\circ\in\cP$, all of whose
$V$-matrices have rank one, with
$\mathfrak c(\theta^\circ;\cdot,\cdot,\cdot)\not\equiv0$; in particular $Z_4$ is a
proper algebraic subset of $\cP$.
\end{lemma}

\begin{proof}
We exhibit $\theta^\circ$ and a single evaluation point at which
$\mathfrak c\neq0$.  Note $N\ge L+1$ for $L\ge2$ (indeed
$N-(L+1)=\binom L2+L-3\ge0$), so the standard vectors $e_1,\dots,e_N$, and in
particular $e_1,\dots,e_{L}$, are available in $\R^d$.

\emph{Step 0 (the layers and the evaluation point).}  Set
\[
V^\circ_\ell:=e_{\ell+1}e_\ell^\top\ \ (1\le\ell\le L-1),\qquad
V^\circ_L:=e_1e_1^\top,
\]
all of rank one, so $B^\circ_\ell=I+e_{\ell+1}e_\ell^\top$ for $\ell\le L-1$; each
$B^\circ_\ell$ is unipotent ($e_{\ell+1}e_\ell^\top$ is nilpotent of order two, as
$e_\ell^\top e_{\ell+1}=0$), with
$\det B^\circ_\ell=1$, $\det(B^\circ_\ell-aV^\circ_\ell)
=\det\big(I+(1-a)e_{\ell+1}e_\ell^\top\big)=1$ for \emph{every} scalar $a$, and
\[
\big(I+a\,e_{\ell+1}e_\ell^\top\big)^{-1}=I-a\,e_{\ell+1}e_\ell^\top .
\]
The matrices $A^\circ_j$ are specified in Step 4.  Evaluate at
\[
t=(0,\dots,0),\qquad s=(\tfrac12,\dots,\tfrac12),\qquad w^\circ=e_1
\]
--- a legitimate evaluation point: $\mathfrak c(\theta^\circ;\cdot)$ is a polynomial
on $\R^{L-1}\times\R^{L-1}\times\R^d$, and nonvanishing at \emph{any} point
establishes $\not\equiv0$ (the open box $(0,1)^{2(L-1)}$ is only needed later, in
Proposition~\ref{prop:anchorsexist}, where Lemma~\ref{lem:zariski} supplies interior
points).

\emph{Step 1 (the transported probes).}  At $t=0$, $W'_k(0,\dots,0)=B^\circ_{k:1}$,
so $w_{(k)}=B^\circ_{k:1}e_1$.  By induction,
\begin{equation}\label{eq:Wrec}
w_{(k)}=e_1+e_2+\dots+e_{k+1}\qquad(0\le k\le L-1):
\end{equation}
indeed $w_{(0)}=e_1$ and
$B^\circ_{k+1}w_{(k)}=w_{(k)}+e_{k+2}\big(e_{k+1}^\top w_{(k)}\big)
=w_{(k)}+e_{k+2}$, the last coordinate pairing being $1$ by \eqref{eq:Wrec}.

\emph{Step 2 (the source vectors).}  By Auxiliary Lemma~\ref{lem:grad}, at the
evaluation point every gradient takes the form
$\pm(B^\circ_{m:1})^\top(A^\circ_j)^\top f$ with a \emph{group index} $j$ (the index of
the $A$-matrix involved), a common prefix $m=j-1$, and a \emph{source vector} $f$:
\begin{itemize}
\item $\nabla_vE_j=(B^\circ_{j-1:1})^\top(A^\circ_j)^\top\,w_{(j-1)}$
\hfill($1\le j\le L-1$),
\item $\nabla_vS_{k,\ell}=(B^\circ_{j-1:1})^\top(A^\circ_j)^\top\,a_k$ with
$j:=k+\ell$, where, by \eqref{eq:Wrec},
\begin{equation}\label{eq:Sid}
a_k:=\big(B^\circ_{k+1}-s_kV^\circ_{k+1}\big)w_{(k)}
=w_{(k)}+(1-s_k)\,e_{k+2}\qquad(0\le k\le L-2),
\end{equation}
\item $\nabla_vT_k=-(B^\circ_{k:1})^\top(A^\circ_{k+1})^\top\,
\adj(B^\circ_{k+1})\,V^\circ_{k+1}w_{(k)}
=-(B^\circ_{j-1:1})^\top(A^\circ_j)^\top\,e_{j+1}$ with $j:=k+1$, since
\begin{equation}\label{eq:Tid}
V^\circ_{k+1}w_{(k)}=e_{k+2}\big(e_{k+1}^\top w_{(k)}\big)=e_{k+2},\qquad
\adj(B^\circ_{k+1})\,e_{k+2}
=\big(I-e_{k+2}e_{k+1}^\top\big)e_{k+2}=e_{k+2}
\end{equation}
($\adj=\det\cdot\,\mathrm{inverse}$ and $\det B^\circ_{k+1}=1$).
\end{itemize}
Thus the rows organize into \emph{groups} by the index $j$ of the attention matrix:
group $j$ ($1\le j\le L-1$) consists of $E_j$, of $S_{k,j-k}$ for $0\le k\le j-2$, and
of $T_{j-1}$ --- that is $1+(j-1)+1=j+1$ rows --- and group $L$ consists of
$S_{k,L-k}$ for $0\le k\le L-2$, i.e.\ $L-1$ rows; the total is
$\sum_{j=1}^{L-1}(j+1)+(L-1)=N$, as it must be.  The source families are
\begin{equation}\label{eq:Svec}
\cF_j:=\{w_{(j-1)}\}\cup\{a_k:0\le k\le j-2\}\cup\{e_{j+1}\}\ \ (1\le j\le L-1),
\qquad
\cF_L:=\{a_k:0\le k\le L-2\}.
\end{equation}

\emph{Step 3 (each $\cF_j$ is linearly independent).}  All members of $\cF_j$ lie in
$\spn\{e_1,\dots,e_{j+1}\}$ (for $j\le L-1$), resp.\ $\spn\{e_1,\dots,e_L\}$ (for
$j=L$), and $|\cF_j|=j+1$, resp.\ $L-1$.  For $j\le L-1$: $e_{j+1}$ is the only member
with a nonzero $e_{j+1}$-coordinate (the $a_k$ have top index $k+2\le j$, and
$w_{(j-1)}$ has top index $j$), so it suffices to prove independence of
$\{w_{(j-1)}\}\cup\{a_k\}_{k=0}^{j-2}$ inside $\spn\{e_1,\dots,e_j\}$.  Form the
$j\times j$ matrix with rows indexed by $e_1,\dots,e_j$ and columns
$w_{(j-1)},a_0,\dots,a_{j-2}$, and subtract from each row the previous one (from the
bottom up): the column $w_{(j-1)}$ (all ones) becomes $e_1$, and the column $a_k$
acquires the entries $-s_k$ in row $k+2$, $s_k-1$ in row $k+3$ (when $\le j$), and $0$
elsewhere in rows $2,\dots,j$.  Expanding along the first column, the determinant is
$\pm$ the determinant of a lower bidiagonal matrix with diagonal entries $-s_k\neq0$:
\[
\det=\pm\prod_{k=0}^{j-2}s_k\neq0\qquad(s_k=\tfrac12),
\]
so $\cF_j$ is independent.  For $j=L$: the $a_k$ have pairwise distinct top indices
$k+2\in\{2,\dots,L\}$ with top coefficients $1-s_k\neq0$, so they are independent by
triangularity.

\emph{Step 4 (the attention matrices and joint independence).}  Partition
$\{1,\dots,N\}$ into consecutive blocks of sizes $n_1,\dots,n_L$, where
$n_j:=j+1$ ($j\le L-1$) and $n_L:=L-1$, with offsets $o_j:=\sum_{j'<j}n_{j'}$
(possible since $\sum_jn_j=N\le d$).  For each $j$ choose a linear map
$C_j\in\R^{d\times d}$ that is injective on $\spn\cF_j$ with image inside
$\spn\{e_{o_j+1},\dots,e_{o_j+n_j}\}$: concretely, fix the basis $\cF_j$ of
$\spn\cF_j$ (Step 3), let $\Lambda_j:\R^d\to\R^{n_j}$ extract the coefficients in this
basis after orthogonal projection onto $\spn\cF_j$, and let $C_j$ map
$x\mapsto\sum_{i=1}^{n_j}(\Lambda_jx)_i\,e_{o_j+i}$.  Now define
\[
A^\circ_j:=C_j^\top\,\big(B^\circ_{j-1:1}\big)^{-1}\qquad(1\le j\le L),
\]
so that $(B^\circ_{j-1:1})^\top(A^\circ_j)^\top=C_j$ and every gradient of group $j$
equals $\pm C_jf$ for its source $f\in\cF_j$.  Suppose a linear combination of all $N$
gradients vanishes.  Decomposing along the pairwise disjoint coordinate blocks, the
combination within each group vanishes separately; since $C_j$ is injective on
$\spn\cF_j$ and $\cF_j$ is independent, the vectors $\{C_jf:f\in\cF_j\}$ are
independent, so all coefficients vanish.  Hence the $N$ gradients are linearly
independent and $\det(\mathcal G^\top\mathcal G)\neq0$ at the evaluation point.

\emph{Step 5 (the remaining factors).}  $\det(B^\circ_{k+1}-t_{k+1}V^\circ_{k+1})=1$
for every $k$ (Step 0, with $a=t_{k+1}=0$; in fact for every $t$).  And for
$0\le k\le L-2$, with $j=k+1\le L-1$,
\[
(A^\circ_{k+1})^\top w_{(k)}
=\big((B^\circ_{k:1})^{-1}\big)^\top C_{k+1}\,w_{(k)}\neq0,
\]
because $w_{(k)}=w_{(j-1)}\in\cF_j$ is a basis vector of $\spn\cF_j$, so
$C_{j}w_{(j-1)}\neq0$, and the left factor is invertible.  All three factors of
$\mathfrak c(\theta^\circ;0,s,e_1)$ are therefore nonzero, so
$\mathfrak c(\theta^\circ;\cdot)\not\equiv0$.
\end{proof}

\begin{remark}[The bound $d\ge N$ is intrinsic to (G4)]\label{rem:dlessN}
If $d<N$, then for \emph{every} $\theta'$ and every $(t,s,w^\circ)$ the $N$ gradient
columns of $\mathcal G$ lie in $\R^d$ and are linearly dependent, so
$\det(\mathcal G^\top\mathcal G)\equiv0$ and (G4) fails identically: the certificate
route to anchors requires $d\ge N$.  This says nothing about identifiability itself
below that dimension; cf.\ Remark~\ref{rem:dimension}.
\end{remark}

\begin{remark}[Interfaces]\label{rem:interfaces}
Section~\ref{sec:generic} consumes only the slope-polynomial structure of
\S\ref{subsec:slope} and Lemma~\ref{lem:zariski}; it is logically independent of
\S\S\ref{sec:A1}--\ref{sec:descent}.  Conversely, \S\ref{sec:descent} never invokes
Propositions~\ref{prop:anchorsexist} or \ref{prop:genericnonempty}: the anchor needed
inside the induction is supplied by the hypothesis of Theorem~\ref{thm:IDL} and
re-established by the sweep.  The only forward citation in this section ---
Proposition~\ref{prop:genericnonempty} citing Lemma~\ref{lem:witness} --- is harmless:
the witness lemma is self-contained.  The dependency graph of the paper is therefore
acyclic.
\end{remark}

\subsection{Proofs of Theorem~\ref{thm:main} and Corollary~\ref{cor:factors}}
\label{subsec:wrapup}

\begin{proof}[Proof of Theorem~\ref{thm:main}]
Let $\cN$ be the finite union of proper algebraic subsets of $\cP$ provided by
Proposition~\ref{prop:genericnonempty} (applicable: $d\ge\max(2,d_*(L))$ and
$N=d_*(L)$), so that $\cP\setminus\cN\subseteq\cG^{(L)}$; by
Lemma~\ref{lem:zariski}, $\cN$ is closed and Lebesgue null (a finite union of zero
sets of nonzero polynomials).  Fix $\theta'\in\cP\setminus\cN$ and $\theta\in\cP$ with
$\TF_\theta=\TF_{\theta'}$ on all probe inputs $X_{w,v}(\tau)$,
$(w,v,\tau)\in\R^{2d}\times(0,\infty)$, in the last output column --- the case of
agreement on all of $\R^{d\times(r+1)}$ being a fortiori.  By \eqref{eq:F},
\[
F^{(L)}_\theta(w,v,\tau)=F^{(L)}_{\theta'}(w,v,\tau)
\qquad\text{for all }(w,v,\tau)\in\R^{2d}\times(0,\infty).
\]
Set $\cO:=\R^{2d}$.  For every $P\in\cP(\cO)$, say $P\in\cH_{T_0}$ with $T_0>0$, and
every real $\tau>T_0$, the point $(w(\tau),v(\tau))$ is real, so
\[
G^P_\theta(\tau)=F^{(L)}_\theta\big(w(\tau),v(\tau),\tau\big)
=F^{(L)}_{\theta'}\big(w(\tau),v(\tau),\tau\big)=G^P_{\theta'}(\tau):
\]
the agreement clause of Theorem~\ref{thm:IDL} holds with $T_P:=T_0$.  The anchor
clause holds as well: for $L=1$ no condition is required (and
$\cM_1(\theta')=\R^{2d}$ in any case), while for $L\ge2$,
Proposition~\ref{prop:anchorsexist} --- applicable since $d\ge N$ and
$\theta'\in\cG^{(L)}$ --- gives
$\cO\cap\cM_L(\theta')=\cM_L(\theta')\neq\varnothing$.  Theorem~\ref{thm:IDL} yields
$\theta=\theta'$.
\end{proof}

\begin{proof}[Proof of Corollary~\ref{cor:factors}]
Let
$\mathfrak F:=\prod_{\ell=1}^L\big(\R^{d\times d_v}\times\R^{d_v\times d}
\times\R^{d\times d}\times\R^{d\times d}\big)$ be the factor parameter space and
$\mu:\mathfrak F\to\cP$ the polynomial map
\[
\mu\Big(\big(W_O^{(\ell)},W_V^{(\ell)},W_K^{(\ell)},W_Q^{(\ell)}\big)_\ell\Big)
:=\Big(W_O^{(\ell)}W_V^{(\ell)},\ (W_K^{(\ell)})^\top W_Q^{(\ell)}\Big)_\ell .
\]
Write $\cN=Z_1\cup\dots\cup Z_M$ for the finite union of proper algebraic
constituents constructed in the proof of Proposition~\ref{prop:genericnonempty}
(each $Z_i$ the common zero set of finitely many polynomials), and define
\[
\cN_{\mathrm{fac}}:=\bigcup_{i=1}^M\mu^{-1}(Z_i)\ \cup\
\bigcup_{\ell=1}^L\big\{\rank\big(W_O^{(\ell)}W_V^{(\ell)}\big)<d_v\big\}\ \cup\
\bigcup_{\ell=1}^L\big\{\rank\big((W_K^{(\ell)})^\top W_Q^{(\ell)}\big)<d\big\}.
\]

\emph{$\cN_{\mathrm{fac}}$ is a finite union of proper algebraic subsets.}
Algebraicity: each $\mu^{-1}(Z_i)$ is cut out by the compositions of $Z_i$'s defining
polynomials with the polynomial map $\mu$; each rank condition
$\rank(M)<\rho$ is the vanishing of all $\rho\times\rho$ minors of the indicated
product, polynomials in the factors.  Properness: for the rank constituents, the
witness $W_O^{(\ell)}=\big[\begin{smallmatrix}I_{d_v}\\0\end{smallmatrix}\big]$,
$W_V^{(\ell)}=[\,I_{d_v}\ 0\,]$, $W_K^{(\ell)}=W_Q^{(\ell)}=I$ gives
$\rank(W_OW_V)=\rank\big(\begin{smallmatrix}I_{d_v}&0\\0&0\end{smallmatrix}\big)=d_v$
and $\rank(W_K^\top W_Q)=\rank(I)=d$.  For the pullbacks: by the recording at the end
of the proof of Proposition~\ref{prop:genericnonempty}, each $Z_i$ admits a witness
$\theta^{(i)}\in\cP\setminus Z_i$ all of whose $V$-matrices have rank $\le1\le d_v$;
every rank-$\le1$ matrix $ab^\top$ factors as $W_OW_V$ with
$W_O:=[\,a\ 0\ \cdots\ 0\,]\in\R^{d\times d_v}$ and $W_V$ the matrix with first row
$b^\top$ and zero rows below; and every $A\in\R^{d\times d}$ equals
$I^\top A=(W_K)^\top W_Q$ with $W_K:=I$, $W_Q:=A$.  Hence
$\theta^{(i)}\in\mu(\mathfrak F)$, and any preimage of $\theta^{(i)}$ lies outside
$\mu^{-1}(Z_i)$.

\emph{Forward direction.}  Let the factor tuple $F'$ of $\theta'$ lie outside
$\cN_{\mathrm{fac}}$ and suppose $\TF_\theta=\TF_{\theta'}$ on
$\R^{d\times(r+1)}$, where $\theta=\mu(F)$ for the unprimed factor tuple $F$.  Then
$\theta'=\mu(F')\notin Z_1\cup\dots\cup Z_M=\cN$, so Theorem~\ref{thm:main} (whose
hypotheses on $L,d,r$ are assumed) gives $\theta=\theta'$, i.e.\ for every $\ell$:
\[
W_O^{(\ell)}W_V^{(\ell)}=W_O'^{(\ell)}W_V'^{(\ell)}=:V_\ell',
\qquad
(W_K^{(\ell)})^\top W_Q^{(\ell)}=(W_K'^{(\ell)})^\top W_Q'^{(\ell)}=:A_\ell',
\]
with $\rank V_\ell'=d_v$ and $\rank A_\ell'=d$ ($F'\notin\cN_{\mathrm{fac}}$).  Apply
Lemma~\ref{lem:lift} with $\rho=d_v$, $A:=W_O'^{(\ell)}$, $B:=W_V'^{(\ell)}$,
$A':=W_O^{(\ell)}$, $B':=W_V^{(\ell)}$: there is a unique
$G^{(v)}_\ell\in GL_{d_v}(\R)$ with $W_O^{(\ell)}=W_O'^{(\ell)}G^{(v)}_\ell$ and
$W_V^{(\ell)}=(G^{(v)}_\ell)^{-1}W_V'^{(\ell)}$.  Likewise with $\rho=d$,
$A:=(W_K'^{(\ell)})^\top$, $B:=W_Q'^{(\ell)}$, $A':=(W_K^{(\ell)})^\top$,
$B':=W_Q^{(\ell)}$: a unique $G^{(a)}_\ell\in GL_d(\R)$ with
$(W_K^{(\ell)})^\top=(W_K'^{(\ell)})^\top G^{(a)}_\ell$ and
$W_Q^{(\ell)}=(G^{(a)}_\ell)^{-1}W_Q'^{(\ell)}$.

\emph{Converse.}  Any such change of factors leaves every product
$W_O^{(\ell)}W_V^{(\ell)}$ and $(W_K^{(\ell)})^\top W_Q^{(\ell)}$, hence
$\theta$ and $\TF_\theta$, unchanged.  Layers being ordered with one head each, no
further (e.g.\ permutation) symmetries occur.
\end{proof}

\section{Remarks}\label{sec:remarks}

\begin{remark}[The two architectural features]\label{rem:architecture}
Causal masking is what closes the probe family under the layer map
(Proposition~\ref{prop:probe}): the $r$ identical leading columns are transported
rigidly regardless of the scores, and only the last column carries a gate.  The skip
connections enter twice, both times through the single identity $B_\ell-V_\ell=I$
\eqref{eq:transparency}: as \emph{transparency} of fully attending layers on the
generic side, and as the \emph{$K=I$ cancellation} (Proposition~\ref{prop:matching},
Step 6) that eliminates all dependence on the unknown parameter's saturation pattern.
\end{remark}

\begin{remark}[The complex-analytic core]\label{rem:complexcore}
All complex analysis is confined to \S\ref{sec:A1}; it follows Fefferman's singular-set
strategy \cite{fefferman1994reconstructing}, with arithmetic progressions of poles of
the logistic sigmoid in place of the poles of $\tanh$ and with the tier hierarchy
replacing his iterated function-theoretic reconstruction.  The bookkeeping device that
makes the hierarchy uniform in the unknown parameter is the largeness region of
Lemma~\ref{lem:nested}, whose thresholds are continuous rather than constant ---
constants cannot work, because the deeper gates blow up at rates governed by the
varying values of the earlier gates.
\end{remark}

\begin{remark}[Sequence length]\label{rem:seqlength}
The probe family uses sequences of length $T=r+1$ with $r\ge2$, and the offset
$b=\log r$ in every gate is forced exactly by the softmax normalization over the $r$
identical key columns (proof of Proposition~\ref{prop:probe}).  Since $b>0$ is all
that the arguments use (via $\alpha=\sig(b)\in(0,1)$ and the real-part comparison in
\S\ref{sec:A1}, which requires $b\neq0$), any fixed $r\ge2$ works, and the
identification is already forced by this one probe family at one sequence length.
\end{remark}

\begin{remark}[The dimension hypothesis]\label{rem:dimension}
The bound $d\ge d_*(L)=\binom L2+2(L-1)$ enters only through the anchor certificate:
Proposition~\ref{prop:anchorsexist} produces anchors by a single full-rank affine
solve against $N=d_*(L)$ certificate rows, and Lemma~\ref{lem:witness} spends the
coordinate vectors $e_1,\dots,e_N$ on its block construction.  We have made no attempt
to optimize it: anchors are points subject to $L-1$ equalities and finitely many open
conditions, so one expects them to exist for far smaller $d$; any such improvement
feeds directly into Theorem~\ref{thm:main} through the anchor clause of
Theorem~\ref{thm:IDL}, the rest of the proof requiring only $d\ge2$.
\end{remark}

\begin{remark}[Exotic saturation profiles]\label{rem:exotic}
A priori the unknown parameter's gates could oscillate, saturate at different values
on different bases, or converge to values outside $\{0,1,\alpha\}$.  The trichotomy
(Proposition~\ref{prop:trichotomy}) eliminates all of this by linear algebra: on a
connected sign component each level either saturates by sign rigidity, or its limit
slope vanishes identically, in which case the curve identity \eqref{eq:curve} survives
off the quadric and the coefficient extraction against $\det A_1\neq0$,
$\Sym A_1\neq0$ and the depth certificate $V_1\neq0$ (Corollary~\ref{cor:depth})
forces the neutral limit $\alpha$.  It is here, and only here, that the genericity
items (G1)--(G2) beyond mere nonvanishing of $A_1'$ earn their keep.
\end{remark}

\begin{thebibliography}{9}

\bibitem{fefferman1994reconstructing}
C.~Fefferman,
\emph{Reconstructing a neural net from its output},
Revista Matem\'atica Iberoamericana \textbf{10} (1994), 507--555.

\bibitem{sussmann1992}
H.~J. Sussmann,
\emph{Uniqueness of the weights for minimal feedforward nets with a given
input--output map},
Neural Networks \textbf{5} (1992), 589--593.

\bibitem{albertinisontag}
F.~Albertini and E.~D. Sontag,
\emph{For neural networks, function determines form},
Neural Networks \textbf{6} (1993), 975--990.

\bibitem{henry2025geometry}
A.~Henry,
\emph{On the geometry of low-rank factorizations},
preprint (2025).

\end{thebibliography}

\end{document}